%!TEX program = xelatex
\documentclass[11pt,a4paper]{article}
\usepackage[margin=2.2cm]{geometry}
\usepackage{amsmath,amssymb,amsthm,mathtools}
\usepackage{bm}
\usepackage{enumitem}
\usepackage{booktabs}
\usepackage{xcolor}
\newcommand{\mserr}[1]{\textcolor{red}{#1}}
\newcommand{\errpanel}[1]{%
\begin{center}
\fcolorbox{red!75!black}{red!8}{%
\begin{minipage}{0.94\linewidth}
\setlength{\parskip}{0.45em plus 0.1em}
#1
\end{minipage}}
\end{center}
}
\usepackage[round,authoryear]{natbib}
\usepackage{hyperref}
\usepackage{fancyhdr}
\usepackage{xeCJK}

\hypersetup{
    colorlinks=true,
    linkcolor=blue!60!black,
    citecolor=blue!60!black,
    urlcolor=blue!60!black
}

\pagestyle{fancy}
\fancyhf{}
\rhead{\small Qiuzhen QE Probability \& Statistics --- Spring 2025}
\lhead{\small Official Qualifying Exam}
\rfoot{\small Page \thepage}

\DeclareMathOperator*{\argmax}{arg\,max}
\DeclareMathOperator*{\argmin}{arg\,min}
\DeclareMathOperator{\Var}{Var}
\DeclareMathOperator{\E}{\mathbb{E}}
\DeclareMathOperator{\Pbb}{\mathbb{P}}
\DeclareMathOperator{\Cov}{Cov}

\setlist[itemize]{topsep=3pt,itemsep=2pt,parsep=0pt}
\setlist[enumerate]{topsep=3pt,itemsep=2pt,parsep=0pt}

\newcommand{\examstatement}[1]{%
  \par\addvspace{0.85em}%
  \noindent{\bfseries #1}\par\nobreak\addvspace{0.28em}%
}

% Shared amsthm note/remark/theorem styles + cleveref
% Shared math extras for qe-review.
% Requires already-loaded: amsmath, amssymb.
%
% Classic pitfall: AMS blackboard-bold fonts have no digit glyphs, so
% \mathbb{1} renders as overlapping garbage. Map the digit 1 to dsfont's
% \mathds{1}; leave \mathbb{R}, \mathbb{P}, etc. on AMS.
%
% Usage (path relative to the main .tex being compiled):
%   notes:  % Shared math extras for qe-review.
% Requires already-loaded: amsmath, amssymb.
%
% Classic pitfall: AMS blackboard-bold fonts have no digit glyphs, so
% \mathbb{1} renders as overlapping garbage. Map the digit 1 to dsfont's
% \mathds{1}; leave \mathbb{R}, \mathbb{P}, etc. on AMS.
%
% Usage (path relative to the main .tex being compiled):
%   notes:  % Shared math extras for qe-review.
% Requires already-loaded: amsmath, amssymb.
%
% Classic pitfall: AMS blackboard-bold fonts have no digit glyphs, so
% \mathbb{1} renders as overlapping garbage. Map the digit 1 to dsfont's
% \mathds{1}; leave \mathbb{R}, \mathbb{P}, etc. on AMS.
%
% Usage (path relative to the main .tex being compiled):
%   notes:  \input{../../common/qe-math-extras.tex}
%   exams:  \input{../../../common/qe-math-extras.tex}

\usepackage{dsfont}
\usepackage{pdftexcmds}

\makeatletter
\let\qe@amsmmathbb\mathbb
\renewcommand{\mathbb}[1]{%
  \ifnum\pdf@strcmp{\unexpanded{#1}}{1}=\z@
    \mathds{1}%
  \else
    \qe@amsmmathbb{#1}%
  \fi
}
\makeatother

% Preferred indicator macros (either style is safe after the remap above).
\newcommand{\bbone}{\mathds{1}}
\newcommand{\ind}{\mathbf{1}}

%   exams:  % Shared math extras for qe-review.
% Requires already-loaded: amsmath, amssymb.
%
% Classic pitfall: AMS blackboard-bold fonts have no digit glyphs, so
% \mathbb{1} renders as overlapping garbage. Map the digit 1 to dsfont's
% \mathds{1}; leave \mathbb{R}, \mathbb{P}, etc. on AMS.
%
% Usage (path relative to the main .tex being compiled):
%   notes:  \input{../../common/qe-math-extras.tex}
%   exams:  \input{../../../common/qe-math-extras.tex}

\usepackage{dsfont}
\usepackage{pdftexcmds}

\makeatletter
\let\qe@amsmmathbb\mathbb
\renewcommand{\mathbb}[1]{%
  \ifnum\pdf@strcmp{\unexpanded{#1}}{1}=\z@
    \mathds{1}%
  \else
    \qe@amsmmathbb{#1}%
  \fi
}
\makeatother

% Preferred indicator macros (either style is safe after the remap above).
\newcommand{\bbone}{\mathds{1}}
\newcommand{\ind}{\mathbf{1}}


\usepackage{dsfont}
\usepackage{pdftexcmds}

\makeatletter
\let\qe@amsmmathbb\mathbb
\renewcommand{\mathbb}[1]{%
  \ifnum\pdf@strcmp{\unexpanded{#1}}{1}=\z@
    \mathds{1}%
  \else
    \qe@amsmmathbb{#1}%
  \fi
}
\makeatother

% Preferred indicator macros (either style is safe after the remap above).
\newcommand{\bbone}{\mathds{1}}
\newcommand{\ind}{\mathbf{1}}

%   exams:  % Shared math extras for qe-review.
% Requires already-loaded: amsmath, amssymb.
%
% Classic pitfall: AMS blackboard-bold fonts have no digit glyphs, so
% \mathbb{1} renders as overlapping garbage. Map the digit 1 to dsfont's
% \mathds{1}; leave \mathbb{R}, \mathbb{P}, etc. on AMS.
%
% Usage (path relative to the main .tex being compiled):
%   notes:  % Shared math extras for qe-review.
% Requires already-loaded: amsmath, amssymb.
%
% Classic pitfall: AMS blackboard-bold fonts have no digit glyphs, so
% \mathbb{1} renders as overlapping garbage. Map the digit 1 to dsfont's
% \mathds{1}; leave \mathbb{R}, \mathbb{P}, etc. on AMS.
%
% Usage (path relative to the main .tex being compiled):
%   notes:  \input{../../common/qe-math-extras.tex}
%   exams:  \input{../../../common/qe-math-extras.tex}

\usepackage{dsfont}
\usepackage{pdftexcmds}

\makeatletter
\let\qe@amsmmathbb\mathbb
\renewcommand{\mathbb}[1]{%
  \ifnum\pdf@strcmp{\unexpanded{#1}}{1}=\z@
    \mathds{1}%
  \else
    \qe@amsmmathbb{#1}%
  \fi
}
\makeatother

% Preferred indicator macros (either style is safe after the remap above).
\newcommand{\bbone}{\mathds{1}}
\newcommand{\ind}{\mathbf{1}}

%   exams:  % Shared math extras for qe-review.
% Requires already-loaded: amsmath, amssymb.
%
% Classic pitfall: AMS blackboard-bold fonts have no digit glyphs, so
% \mathbb{1} renders as overlapping garbage. Map the digit 1 to dsfont's
% \mathds{1}; leave \mathbb{R}, \mathbb{P}, etc. on AMS.
%
% Usage (path relative to the main .tex being compiled):
%   notes:  \input{../../common/qe-math-extras.tex}
%   exams:  \input{../../../common/qe-math-extras.tex}

\usepackage{dsfont}
\usepackage{pdftexcmds}

\makeatletter
\let\qe@amsmmathbb\mathbb
\renewcommand{\mathbb}[1]{%
  \ifnum\pdf@strcmp{\unexpanded{#1}}{1}=\z@
    \mathds{1}%
  \else
    \qe@amsmmathbb{#1}%
  \fi
}
\makeatother

% Preferred indicator macros (either style is safe after the remap above).
\newcommand{\bbone}{\mathds{1}}
\newcommand{\ind}{\mathbf{1}}


\usepackage{dsfont}
\usepackage{pdftexcmds}

\makeatletter
\let\qe@amsmmathbb\mathbb
\renewcommand{\mathbb}[1]{%
  \ifnum\pdf@strcmp{\unexpanded{#1}}{1}=\z@
    \mathds{1}%
  \else
    \qe@amsmmathbb{#1}%
  \fi
}
\makeatother

% Preferred indicator macros (either style is safe after the remap above).
\newcommand{\bbone}{\mathds{1}}
\newcommand{\ind}{\mathbf{1}}


\usepackage{dsfont}
\usepackage{pdftexcmds}

\makeatletter
\let\qe@amsmmathbb\mathbb
\renewcommand{\mathbb}[1]{%
  \ifnum\pdf@strcmp{\unexpanded{#1}}{1}=\z@
    \mathds{1}%
  \else
    \qe@amsmmathbb{#1}%
  \fi
}
\makeatother

% Preferred indicator macros (either style is safe after the remap above).
\newcommand{\bbone}{\mathds{1}}
\newcommand{\ind}{\mathbf{1}}

% Shared amsthm environments for qe-review (minimal).
% Requires already-loaded: amsthm.
% Safe to \input after \usepackage{amsmath,amssymb,amsthm,...}.
%
% Shared numbering uses aliascnt so cleveref keys off the environment
% name (Lemma, Proposition, Exercise, ...) rather than the parent
% counter (Theorem, Definition, Remark). Without aliascnt, \cref{lem:...}
% prints ``Theorem 9.13'' instead of ``Lemma 9.13''.

\usepackage{aliascnt}

\theoremstyle{plain}
\newtheorem{theorem}{Theorem}[section]

\newaliascnt{lemma}{theorem}
\newtheorem{lemma}[lemma]{Lemma}
\aliascntresetthe{lemma}

\newaliascnt{proposition}{theorem}
\newtheorem{proposition}[proposition]{Proposition}
\aliascntresetthe{proposition}

\newaliascnt{corollary}{theorem}
\newtheorem{corollary}[corollary]{Corollary}
\aliascntresetthe{corollary}

\newaliascnt{claim}{theorem}
\newtheorem{claim}[claim]{Claim}
\aliascntresetthe{claim}

\newaliascnt{fact}{theorem}
\newtheorem{fact}[fact]{Fact}
\aliascntresetthe{fact}

\theoremstyle{definition}
\newtheorem{definition}{Definition}[section]

\newaliascnt{exercise}{definition}
\newtheorem{exercise}[exercise]{Exercise}
\aliascntresetthe{exercise}

\newaliascnt{assumption}{definition}
\newtheorem{assumption}[assumption]{Assumption}
\aliascntresetthe{assumption}

\newaliascnt{notation}{definition}
\newtheorem{notation}[notation]{Notation}
\aliascntresetthe{notation}

\newaliascnt{construction}{definition}
\newtheorem{construction}[construction]{Construction}
\aliascntresetthe{construction}

% Example: document-wide numbering (not per section / chapter)
\newtheorem{example}{Example}

\theoremstyle{remark}
\newtheorem{remark}{Remark}[section]
\newtheorem*{remark*}{Remark}

\newaliascnt{heuristic}{remark}
\newtheorem{heuristic}[heuristic]{Heuristic}
\aliascntresetthe{heuristic}

\newaliascnt{convention}{remark}
\newtheorem{convention}[convention]{Convention}
\aliascntresetthe{convention}

\newaliascnt{warning}{remark}
\newtheorem{warning}[warning]{Warning}
\aliascntresetthe{warning}

\newaliascnt{proofsketch}{remark}
\newtheorem{proofsketch}[proofsketch]{Proof sketch}
\aliascntresetthe{proofsketch}

\newaliascnt{checkpoint}{remark}
\newtheorem{checkpoint}[checkpoint]{Checkpoint}
\aliascntresetthe{checkpoint}

% Note: document-wide numbering (not per section)
\newtheorem{note}{Note}
\newtheorem*{note*}{Note}

\newaliascnt{examnote}{note}
\newtheorem{examnote}[examnote]{Exam note}
\aliascntresetthe{examnote}

\newaliascnt{study}{note}
\newtheorem{study}[study]{Study note}
\aliascntresetthe{study}

\newaliascnt{summary}{note}
\newtheorem{summary}[summary]{Summary}
\aliascntresetthe{summary}

\newtheorem{examproblem}{Problem}[section]
\newtheorem*{examproblem*}{Problem}

% Bilingual aliases
\newenvironment{dingli}{\begin{theorem}}{\end{theorem}}
\newenvironment{yinli}{\begin{lemma}}{\end{lemma}}
\newenvironment{mingti}{\begin{proposition}}{\end{proposition}}
\newenvironment{tuilun}{\begin{corollary}}{\end{corollary}}
\newenvironment{dingyi}{\begin{definition}}{\end{definition}}
\newenvironment{li}{\begin{example}}{\end{example}}
\newenvironment{zhu}{\begin{remark}}{\end{remark}}
\newenvironment{biji}{\begin{note}}{\end{note}}
\newenvironment{kaoshi}{\begin{examnote}}{\end{examnote}}

\renewcommand{\qedsymbol}{\ensuremath{\square}}

% Shared cleveref setup for qe-review.
% Load AFTER: hyperref, and AFTER all \newtheorem / amsthm environments.
% Shared theorem counters are aliased in qe-amsthm-envs.tex (aliascnt),
% otherwise \cref on a lemma/proposition prints ``Theorem''.
%
% Usage (path relative to the main .tex being compiled):
%   notes:  % Shared cleveref setup for qe-review.
% Load AFTER: hyperref, and AFTER all \newtheorem / amsthm environments.
% Shared theorem counters are aliased in qe-amsthm-envs.tex (aliascnt),
% otherwise \cref on a lemma/proposition prints ``Theorem''.
%
% Usage (path relative to the main .tex being compiled):
%   notes:  % Shared cleveref setup for qe-review.
% Load AFTER: hyperref, and AFTER all \newtheorem / amsthm environments.
% Shared theorem counters are aliased in qe-amsthm-envs.tex (aliascnt),
% otherwise \cref on a lemma/proposition prints ``Theorem''.
%
% Usage (path relative to the main .tex being compiled):
%   notes:  \input{../../common/qe-cleveref.tex}
%   exams:  \input{../../../common/qe-cleveref.tex}
%
% Prefer \cref / \Cref over \ref. Use \Cref at sentence start.
% Unnumbered sectioning (e.g. \paragraph with default secnumdepth)
% produces an empty cleveref type; map that to "Paragraph".

\usepackage[nameinlink,noabbrev]{cleveref}

% Fallback for unnumbered \paragraph / \subsubsection labels
\crefname{}{Paragraph}{Paragraphs}
\Crefname{}{Paragraph}{Paragraphs}

% Numbered theorem-like (plain)
\crefname{theorem}{Theorem}{Theorems}
\Crefname{theorem}{Theorem}{Theorems}
\crefname{lemma}{Lemma}{Lemmas}
\Crefname{lemma}{Lemma}{Lemmas}
\crefname{proposition}{Proposition}{Propositions}
\Crefname{proposition}{Proposition}{Propositions}
\crefname{corollary}{Corollary}{Corollaries}
\Crefname{corollary}{Corollary}{Corollaries}
\crefname{claim}{Claim}{Claims}
\Crefname{claim}{Claim}{Claims}
\crefname{fact}{Fact}{Facts}
\Crefname{fact}{Fact}{Facts}

% Definition-like
\crefname{definition}{Definition}{Definitions}
\Crefname{definition}{Definition}{Definitions}
\crefname{example}{Example}{Examples}
\Crefname{example}{Example}{Examples}
\crefname{exercise}{Exercise}{Exercises}
\Crefname{exercise}{Exercise}{Exercises}
\crefname{assumption}{Assumption}{Assumptions}
\Crefname{assumption}{Assumption}{Assumptions}
\crefname{notation}{Notation}{Notations}
\Crefname{notation}{Notation}{Notations}
\crefname{construction}{Construction}{Constructions}
\Crefname{construction}{Construction}{Constructions}

% Remark-like
\crefname{remark}{Remark}{Remarks}
\Crefname{remark}{Remark}{Remarks}
\crefname{heuristic}{Heuristic}{Heuristics}
\Crefname{heuristic}{Heuristic}{Heuristics}
\crefname{convention}{Convention}{Conventions}
\Crefname{convention}{Convention}{Conventions}
\crefname{warning}{Warning}{Warnings}
\Crefname{warning}{Warning}{Warnings}
\crefname{proofsketch}{Proof sketch}{Proof sketches}
\Crefname{proofsketch}{Proof sketch}{Proof sketches}
\crefname{checkpoint}{Checkpoint}{Checkpoints}
\Crefname{checkpoint}{Checkpoint}{Checkpoints}

% Document-wide notes / exam problems
\crefname{note}{Note}{Notes}
\Crefname{note}{Note}{Notes}
\crefname{examnote}{Exam note}{Exam notes}
\Crefname{examnote}{Exam note}{Exam notes}
\crefname{study}{Study note}{Study notes}
\Crefname{study}{Study note}{Study notes}
\crefname{summary}{Summary}{Summaries}
\Crefname{summary}{Summary}{Summaries}
\crefname{examproblem}{Problem}{Problems}
\Crefname{examproblem}{Problem}{Problems}

% Sectioning / floats (explicit, noabbrev-friendly)
\crefname{section}{Section}{Sections}
\Crefname{section}{Section}{Sections}
\crefname{subsection}{Subsection}{Subsections}
\Crefname{subsection}{Subsection}{Subsections}
\crefname{subsubsection}{Subsubsection}{Subsubsections}
\Crefname{subsubsection}{Subsubsection}{Subsubsections}
\crefname{paragraph}{Paragraph}{Paragraphs}
\Crefname{paragraph}{Paragraph}{Paragraphs}
\crefname{equation}{Equation}{Equations}
\Crefname{equation}{Equation}{Equations}
\crefname{figure}{Figure}{Figures}
\Crefname{figure}{Figure}{Figures}
\crefname{table}{Table}{Tables}
\Crefname{table}{Table}{Tables}
\crefname{enumi}{Item}{Items}
\Crefname{enumi}{Item}{Items}

%   exams:  % Shared cleveref setup for qe-review.
% Load AFTER: hyperref, and AFTER all \newtheorem / amsthm environments.
% Shared theorem counters are aliased in qe-amsthm-envs.tex (aliascnt),
% otherwise \cref on a lemma/proposition prints ``Theorem''.
%
% Usage (path relative to the main .tex being compiled):
%   notes:  \input{../../common/qe-cleveref.tex}
%   exams:  \input{../../../common/qe-cleveref.tex}
%
% Prefer \cref / \Cref over \ref. Use \Cref at sentence start.
% Unnumbered sectioning (e.g. \paragraph with default secnumdepth)
% produces an empty cleveref type; map that to "Paragraph".

\usepackage[nameinlink,noabbrev]{cleveref}

% Fallback for unnumbered \paragraph / \subsubsection labels
\crefname{}{Paragraph}{Paragraphs}
\Crefname{}{Paragraph}{Paragraphs}

% Numbered theorem-like (plain)
\crefname{theorem}{Theorem}{Theorems}
\Crefname{theorem}{Theorem}{Theorems}
\crefname{lemma}{Lemma}{Lemmas}
\Crefname{lemma}{Lemma}{Lemmas}
\crefname{proposition}{Proposition}{Propositions}
\Crefname{proposition}{Proposition}{Propositions}
\crefname{corollary}{Corollary}{Corollaries}
\Crefname{corollary}{Corollary}{Corollaries}
\crefname{claim}{Claim}{Claims}
\Crefname{claim}{Claim}{Claims}
\crefname{fact}{Fact}{Facts}
\Crefname{fact}{Fact}{Facts}

% Definition-like
\crefname{definition}{Definition}{Definitions}
\Crefname{definition}{Definition}{Definitions}
\crefname{example}{Example}{Examples}
\Crefname{example}{Example}{Examples}
\crefname{exercise}{Exercise}{Exercises}
\Crefname{exercise}{Exercise}{Exercises}
\crefname{assumption}{Assumption}{Assumptions}
\Crefname{assumption}{Assumption}{Assumptions}
\crefname{notation}{Notation}{Notations}
\Crefname{notation}{Notation}{Notations}
\crefname{construction}{Construction}{Constructions}
\Crefname{construction}{Construction}{Constructions}

% Remark-like
\crefname{remark}{Remark}{Remarks}
\Crefname{remark}{Remark}{Remarks}
\crefname{heuristic}{Heuristic}{Heuristics}
\Crefname{heuristic}{Heuristic}{Heuristics}
\crefname{convention}{Convention}{Conventions}
\Crefname{convention}{Convention}{Conventions}
\crefname{warning}{Warning}{Warnings}
\Crefname{warning}{Warning}{Warnings}
\crefname{proofsketch}{Proof sketch}{Proof sketches}
\Crefname{proofsketch}{Proof sketch}{Proof sketches}
\crefname{checkpoint}{Checkpoint}{Checkpoints}
\Crefname{checkpoint}{Checkpoint}{Checkpoints}

% Document-wide notes / exam problems
\crefname{note}{Note}{Notes}
\Crefname{note}{Note}{Notes}
\crefname{examnote}{Exam note}{Exam notes}
\Crefname{examnote}{Exam note}{Exam notes}
\crefname{study}{Study note}{Study notes}
\Crefname{study}{Study note}{Study notes}
\crefname{summary}{Summary}{Summaries}
\Crefname{summary}{Summary}{Summaries}
\crefname{examproblem}{Problem}{Problems}
\Crefname{examproblem}{Problem}{Problems}

% Sectioning / floats (explicit, noabbrev-friendly)
\crefname{section}{Section}{Sections}
\Crefname{section}{Section}{Sections}
\crefname{subsection}{Subsection}{Subsections}
\Crefname{subsection}{Subsection}{Subsections}
\crefname{subsubsection}{Subsubsection}{Subsubsections}
\Crefname{subsubsection}{Subsubsection}{Subsubsections}
\crefname{paragraph}{Paragraph}{Paragraphs}
\Crefname{paragraph}{Paragraph}{Paragraphs}
\crefname{equation}{Equation}{Equations}
\Crefname{equation}{Equation}{Equations}
\crefname{figure}{Figure}{Figures}
\Crefname{figure}{Figure}{Figures}
\crefname{table}{Table}{Tables}
\Crefname{table}{Table}{Tables}
\crefname{enumi}{Item}{Items}
\Crefname{enumi}{Item}{Items}

%
% Prefer \cref / \Cref over \ref. Use \Cref at sentence start.
% Unnumbered sectioning (e.g. \paragraph with default secnumdepth)
% produces an empty cleveref type; map that to "Paragraph".

\usepackage[nameinlink,noabbrev]{cleveref}

% Fallback for unnumbered \paragraph / \subsubsection labels
\crefname{}{Paragraph}{Paragraphs}
\Crefname{}{Paragraph}{Paragraphs}

% Numbered theorem-like (plain)
\crefname{theorem}{Theorem}{Theorems}
\Crefname{theorem}{Theorem}{Theorems}
\crefname{lemma}{Lemma}{Lemmas}
\Crefname{lemma}{Lemma}{Lemmas}
\crefname{proposition}{Proposition}{Propositions}
\Crefname{proposition}{Proposition}{Propositions}
\crefname{corollary}{Corollary}{Corollaries}
\Crefname{corollary}{Corollary}{Corollaries}
\crefname{claim}{Claim}{Claims}
\Crefname{claim}{Claim}{Claims}
\crefname{fact}{Fact}{Facts}
\Crefname{fact}{Fact}{Facts}

% Definition-like
\crefname{definition}{Definition}{Definitions}
\Crefname{definition}{Definition}{Definitions}
\crefname{example}{Example}{Examples}
\Crefname{example}{Example}{Examples}
\crefname{exercise}{Exercise}{Exercises}
\Crefname{exercise}{Exercise}{Exercises}
\crefname{assumption}{Assumption}{Assumptions}
\Crefname{assumption}{Assumption}{Assumptions}
\crefname{notation}{Notation}{Notations}
\Crefname{notation}{Notation}{Notations}
\crefname{construction}{Construction}{Constructions}
\Crefname{construction}{Construction}{Constructions}

% Remark-like
\crefname{remark}{Remark}{Remarks}
\Crefname{remark}{Remark}{Remarks}
\crefname{heuristic}{Heuristic}{Heuristics}
\Crefname{heuristic}{Heuristic}{Heuristics}
\crefname{convention}{Convention}{Conventions}
\Crefname{convention}{Convention}{Conventions}
\crefname{warning}{Warning}{Warnings}
\Crefname{warning}{Warning}{Warnings}
\crefname{proofsketch}{Proof sketch}{Proof sketches}
\Crefname{proofsketch}{Proof sketch}{Proof sketches}
\crefname{checkpoint}{Checkpoint}{Checkpoints}
\Crefname{checkpoint}{Checkpoint}{Checkpoints}

% Document-wide notes / exam problems
\crefname{note}{Note}{Notes}
\Crefname{note}{Note}{Notes}
\crefname{examnote}{Exam note}{Exam notes}
\Crefname{examnote}{Exam note}{Exam notes}
\crefname{study}{Study note}{Study notes}
\Crefname{study}{Study note}{Study notes}
\crefname{summary}{Summary}{Summaries}
\Crefname{summary}{Summary}{Summaries}
\crefname{examproblem}{Problem}{Problems}
\Crefname{examproblem}{Problem}{Problems}

% Sectioning / floats (explicit, noabbrev-friendly)
\crefname{section}{Section}{Sections}
\Crefname{section}{Section}{Sections}
\crefname{subsection}{Subsection}{Subsections}
\Crefname{subsection}{Subsection}{Subsections}
\crefname{subsubsection}{Subsubsection}{Subsubsections}
\Crefname{subsubsection}{Subsubsection}{Subsubsections}
\crefname{paragraph}{Paragraph}{Paragraphs}
\Crefname{paragraph}{Paragraph}{Paragraphs}
\crefname{equation}{Equation}{Equations}
\Crefname{equation}{Equation}{Equations}
\crefname{figure}{Figure}{Figures}
\Crefname{figure}{Figure}{Figures}
\crefname{table}{Table}{Tables}
\Crefname{table}{Table}{Tables}
\crefname{enumi}{Item}{Items}
\Crefname{enumi}{Item}{Items}

%   exams:  % Shared cleveref setup for qe-review.
% Load AFTER: hyperref, and AFTER all \newtheorem / amsthm environments.
% Shared theorem counters are aliased in qe-amsthm-envs.tex (aliascnt),
% otherwise \cref on a lemma/proposition prints ``Theorem''.
%
% Usage (path relative to the main .tex being compiled):
%   notes:  % Shared cleveref setup for qe-review.
% Load AFTER: hyperref, and AFTER all \newtheorem / amsthm environments.
% Shared theorem counters are aliased in qe-amsthm-envs.tex (aliascnt),
% otherwise \cref on a lemma/proposition prints ``Theorem''.
%
% Usage (path relative to the main .tex being compiled):
%   notes:  \input{../../common/qe-cleveref.tex}
%   exams:  \input{../../../common/qe-cleveref.tex}
%
% Prefer \cref / \Cref over \ref. Use \Cref at sentence start.
% Unnumbered sectioning (e.g. \paragraph with default secnumdepth)
% produces an empty cleveref type; map that to "Paragraph".

\usepackage[nameinlink,noabbrev]{cleveref}

% Fallback for unnumbered \paragraph / \subsubsection labels
\crefname{}{Paragraph}{Paragraphs}
\Crefname{}{Paragraph}{Paragraphs}

% Numbered theorem-like (plain)
\crefname{theorem}{Theorem}{Theorems}
\Crefname{theorem}{Theorem}{Theorems}
\crefname{lemma}{Lemma}{Lemmas}
\Crefname{lemma}{Lemma}{Lemmas}
\crefname{proposition}{Proposition}{Propositions}
\Crefname{proposition}{Proposition}{Propositions}
\crefname{corollary}{Corollary}{Corollaries}
\Crefname{corollary}{Corollary}{Corollaries}
\crefname{claim}{Claim}{Claims}
\Crefname{claim}{Claim}{Claims}
\crefname{fact}{Fact}{Facts}
\Crefname{fact}{Fact}{Facts}

% Definition-like
\crefname{definition}{Definition}{Definitions}
\Crefname{definition}{Definition}{Definitions}
\crefname{example}{Example}{Examples}
\Crefname{example}{Example}{Examples}
\crefname{exercise}{Exercise}{Exercises}
\Crefname{exercise}{Exercise}{Exercises}
\crefname{assumption}{Assumption}{Assumptions}
\Crefname{assumption}{Assumption}{Assumptions}
\crefname{notation}{Notation}{Notations}
\Crefname{notation}{Notation}{Notations}
\crefname{construction}{Construction}{Constructions}
\Crefname{construction}{Construction}{Constructions}

% Remark-like
\crefname{remark}{Remark}{Remarks}
\Crefname{remark}{Remark}{Remarks}
\crefname{heuristic}{Heuristic}{Heuristics}
\Crefname{heuristic}{Heuristic}{Heuristics}
\crefname{convention}{Convention}{Conventions}
\Crefname{convention}{Convention}{Conventions}
\crefname{warning}{Warning}{Warnings}
\Crefname{warning}{Warning}{Warnings}
\crefname{proofsketch}{Proof sketch}{Proof sketches}
\Crefname{proofsketch}{Proof sketch}{Proof sketches}
\crefname{checkpoint}{Checkpoint}{Checkpoints}
\Crefname{checkpoint}{Checkpoint}{Checkpoints}

% Document-wide notes / exam problems
\crefname{note}{Note}{Notes}
\Crefname{note}{Note}{Notes}
\crefname{examnote}{Exam note}{Exam notes}
\Crefname{examnote}{Exam note}{Exam notes}
\crefname{study}{Study note}{Study notes}
\Crefname{study}{Study note}{Study notes}
\crefname{summary}{Summary}{Summaries}
\Crefname{summary}{Summary}{Summaries}
\crefname{examproblem}{Problem}{Problems}
\Crefname{examproblem}{Problem}{Problems}

% Sectioning / floats (explicit, noabbrev-friendly)
\crefname{section}{Section}{Sections}
\Crefname{section}{Section}{Sections}
\crefname{subsection}{Subsection}{Subsections}
\Crefname{subsection}{Subsection}{Subsections}
\crefname{subsubsection}{Subsubsection}{Subsubsections}
\Crefname{subsubsection}{Subsubsection}{Subsubsections}
\crefname{paragraph}{Paragraph}{Paragraphs}
\Crefname{paragraph}{Paragraph}{Paragraphs}
\crefname{equation}{Equation}{Equations}
\Crefname{equation}{Equation}{Equations}
\crefname{figure}{Figure}{Figures}
\Crefname{figure}{Figure}{Figures}
\crefname{table}{Table}{Tables}
\Crefname{table}{Table}{Tables}
\crefname{enumi}{Item}{Items}
\Crefname{enumi}{Item}{Items}

%   exams:  % Shared cleveref setup for qe-review.
% Load AFTER: hyperref, and AFTER all \newtheorem / amsthm environments.
% Shared theorem counters are aliased in qe-amsthm-envs.tex (aliascnt),
% otherwise \cref on a lemma/proposition prints ``Theorem''.
%
% Usage (path relative to the main .tex being compiled):
%   notes:  \input{../../common/qe-cleveref.tex}
%   exams:  \input{../../../common/qe-cleveref.tex}
%
% Prefer \cref / \Cref over \ref. Use \Cref at sentence start.
% Unnumbered sectioning (e.g. \paragraph with default secnumdepth)
% produces an empty cleveref type; map that to "Paragraph".

\usepackage[nameinlink,noabbrev]{cleveref}

% Fallback for unnumbered \paragraph / \subsubsection labels
\crefname{}{Paragraph}{Paragraphs}
\Crefname{}{Paragraph}{Paragraphs}

% Numbered theorem-like (plain)
\crefname{theorem}{Theorem}{Theorems}
\Crefname{theorem}{Theorem}{Theorems}
\crefname{lemma}{Lemma}{Lemmas}
\Crefname{lemma}{Lemma}{Lemmas}
\crefname{proposition}{Proposition}{Propositions}
\Crefname{proposition}{Proposition}{Propositions}
\crefname{corollary}{Corollary}{Corollaries}
\Crefname{corollary}{Corollary}{Corollaries}
\crefname{claim}{Claim}{Claims}
\Crefname{claim}{Claim}{Claims}
\crefname{fact}{Fact}{Facts}
\Crefname{fact}{Fact}{Facts}

% Definition-like
\crefname{definition}{Definition}{Definitions}
\Crefname{definition}{Definition}{Definitions}
\crefname{example}{Example}{Examples}
\Crefname{example}{Example}{Examples}
\crefname{exercise}{Exercise}{Exercises}
\Crefname{exercise}{Exercise}{Exercises}
\crefname{assumption}{Assumption}{Assumptions}
\Crefname{assumption}{Assumption}{Assumptions}
\crefname{notation}{Notation}{Notations}
\Crefname{notation}{Notation}{Notations}
\crefname{construction}{Construction}{Constructions}
\Crefname{construction}{Construction}{Constructions}

% Remark-like
\crefname{remark}{Remark}{Remarks}
\Crefname{remark}{Remark}{Remarks}
\crefname{heuristic}{Heuristic}{Heuristics}
\Crefname{heuristic}{Heuristic}{Heuristics}
\crefname{convention}{Convention}{Conventions}
\Crefname{convention}{Convention}{Conventions}
\crefname{warning}{Warning}{Warnings}
\Crefname{warning}{Warning}{Warnings}
\crefname{proofsketch}{Proof sketch}{Proof sketches}
\Crefname{proofsketch}{Proof sketch}{Proof sketches}
\crefname{checkpoint}{Checkpoint}{Checkpoints}
\Crefname{checkpoint}{Checkpoint}{Checkpoints}

% Document-wide notes / exam problems
\crefname{note}{Note}{Notes}
\Crefname{note}{Note}{Notes}
\crefname{examnote}{Exam note}{Exam notes}
\Crefname{examnote}{Exam note}{Exam notes}
\crefname{study}{Study note}{Study notes}
\Crefname{study}{Study note}{Study notes}
\crefname{summary}{Summary}{Summaries}
\Crefname{summary}{Summary}{Summaries}
\crefname{examproblem}{Problem}{Problems}
\Crefname{examproblem}{Problem}{Problems}

% Sectioning / floats (explicit, noabbrev-friendly)
\crefname{section}{Section}{Sections}
\Crefname{section}{Section}{Sections}
\crefname{subsection}{Subsection}{Subsections}
\Crefname{subsection}{Subsection}{Subsections}
\crefname{subsubsection}{Subsubsection}{Subsubsections}
\Crefname{subsubsection}{Subsubsection}{Subsubsections}
\crefname{paragraph}{Paragraph}{Paragraphs}
\Crefname{paragraph}{Paragraph}{Paragraphs}
\crefname{equation}{Equation}{Equations}
\Crefname{equation}{Equation}{Equations}
\crefname{figure}{Figure}{Figures}
\Crefname{figure}{Figure}{Figures}
\crefname{table}{Table}{Tables}
\Crefname{table}{Table}{Tables}
\crefname{enumi}{Item}{Items}
\Crefname{enumi}{Item}{Items}

%
% Prefer \cref / \Cref over \ref. Use \Cref at sentence start.
% Unnumbered sectioning (e.g. \paragraph with default secnumdepth)
% produces an empty cleveref type; map that to "Paragraph".

\usepackage[nameinlink,noabbrev]{cleveref}

% Fallback for unnumbered \paragraph / \subsubsection labels
\crefname{}{Paragraph}{Paragraphs}
\Crefname{}{Paragraph}{Paragraphs}

% Numbered theorem-like (plain)
\crefname{theorem}{Theorem}{Theorems}
\Crefname{theorem}{Theorem}{Theorems}
\crefname{lemma}{Lemma}{Lemmas}
\Crefname{lemma}{Lemma}{Lemmas}
\crefname{proposition}{Proposition}{Propositions}
\Crefname{proposition}{Proposition}{Propositions}
\crefname{corollary}{Corollary}{Corollaries}
\Crefname{corollary}{Corollary}{Corollaries}
\crefname{claim}{Claim}{Claims}
\Crefname{claim}{Claim}{Claims}
\crefname{fact}{Fact}{Facts}
\Crefname{fact}{Fact}{Facts}

% Definition-like
\crefname{definition}{Definition}{Definitions}
\Crefname{definition}{Definition}{Definitions}
\crefname{example}{Example}{Examples}
\Crefname{example}{Example}{Examples}
\crefname{exercise}{Exercise}{Exercises}
\Crefname{exercise}{Exercise}{Exercises}
\crefname{assumption}{Assumption}{Assumptions}
\Crefname{assumption}{Assumption}{Assumptions}
\crefname{notation}{Notation}{Notations}
\Crefname{notation}{Notation}{Notations}
\crefname{construction}{Construction}{Constructions}
\Crefname{construction}{Construction}{Constructions}

% Remark-like
\crefname{remark}{Remark}{Remarks}
\Crefname{remark}{Remark}{Remarks}
\crefname{heuristic}{Heuristic}{Heuristics}
\Crefname{heuristic}{Heuristic}{Heuristics}
\crefname{convention}{Convention}{Conventions}
\Crefname{convention}{Convention}{Conventions}
\crefname{warning}{Warning}{Warnings}
\Crefname{warning}{Warning}{Warnings}
\crefname{proofsketch}{Proof sketch}{Proof sketches}
\Crefname{proofsketch}{Proof sketch}{Proof sketches}
\crefname{checkpoint}{Checkpoint}{Checkpoints}
\Crefname{checkpoint}{Checkpoint}{Checkpoints}

% Document-wide notes / exam problems
\crefname{note}{Note}{Notes}
\Crefname{note}{Note}{Notes}
\crefname{examnote}{Exam note}{Exam notes}
\Crefname{examnote}{Exam note}{Exam notes}
\crefname{study}{Study note}{Study notes}
\Crefname{study}{Study note}{Study notes}
\crefname{summary}{Summary}{Summaries}
\Crefname{summary}{Summary}{Summaries}
\crefname{examproblem}{Problem}{Problems}
\Crefname{examproblem}{Problem}{Problems}

% Sectioning / floats (explicit, noabbrev-friendly)
\crefname{section}{Section}{Sections}
\Crefname{section}{Section}{Sections}
\crefname{subsection}{Subsection}{Subsections}
\Crefname{subsection}{Subsection}{Subsections}
\crefname{subsubsection}{Subsubsection}{Subsubsections}
\Crefname{subsubsection}{Subsubsection}{Subsubsections}
\crefname{paragraph}{Paragraph}{Paragraphs}
\Crefname{paragraph}{Paragraph}{Paragraphs}
\crefname{equation}{Equation}{Equations}
\Crefname{equation}{Equation}{Equations}
\crefname{figure}{Figure}{Figures}
\Crefname{figure}{Figure}{Figures}
\crefname{table}{Table}{Tables}
\Crefname{table}{Table}{Tables}
\crefname{enumi}{Item}{Items}
\Crefname{enumi}{Item}{Items}

%
% Prefer \cref / \Cref over \ref. Use \Cref at sentence start.
% Unnumbered sectioning (e.g. \paragraph with default secnumdepth)
% produces an empty cleveref type; map that to "Paragraph".

\usepackage[nameinlink,noabbrev]{cleveref}

% Fallback for unnumbered \paragraph / \subsubsection labels
\crefname{}{Paragraph}{Paragraphs}
\Crefname{}{Paragraph}{Paragraphs}

% Numbered theorem-like (plain)
\crefname{theorem}{Theorem}{Theorems}
\Crefname{theorem}{Theorem}{Theorems}
\crefname{lemma}{Lemma}{Lemmas}
\Crefname{lemma}{Lemma}{Lemmas}
\crefname{proposition}{Proposition}{Propositions}
\Crefname{proposition}{Proposition}{Propositions}
\crefname{corollary}{Corollary}{Corollaries}
\Crefname{corollary}{Corollary}{Corollaries}
\crefname{claim}{Claim}{Claims}
\Crefname{claim}{Claim}{Claims}
\crefname{fact}{Fact}{Facts}
\Crefname{fact}{Fact}{Facts}

% Definition-like
\crefname{definition}{Definition}{Definitions}
\Crefname{definition}{Definition}{Definitions}
\crefname{example}{Example}{Examples}
\Crefname{example}{Example}{Examples}
\crefname{exercise}{Exercise}{Exercises}
\Crefname{exercise}{Exercise}{Exercises}
\crefname{assumption}{Assumption}{Assumptions}
\Crefname{assumption}{Assumption}{Assumptions}
\crefname{notation}{Notation}{Notations}
\Crefname{notation}{Notation}{Notations}
\crefname{construction}{Construction}{Constructions}
\Crefname{construction}{Construction}{Constructions}

% Remark-like
\crefname{remark}{Remark}{Remarks}
\Crefname{remark}{Remark}{Remarks}
\crefname{heuristic}{Heuristic}{Heuristics}
\Crefname{heuristic}{Heuristic}{Heuristics}
\crefname{convention}{Convention}{Conventions}
\Crefname{convention}{Convention}{Conventions}
\crefname{warning}{Warning}{Warnings}
\Crefname{warning}{Warning}{Warnings}
\crefname{proofsketch}{Proof sketch}{Proof sketches}
\Crefname{proofsketch}{Proof sketch}{Proof sketches}
\crefname{checkpoint}{Checkpoint}{Checkpoints}
\Crefname{checkpoint}{Checkpoint}{Checkpoints}

% Document-wide notes / exam problems
\crefname{note}{Note}{Notes}
\Crefname{note}{Note}{Notes}
\crefname{examnote}{Exam note}{Exam notes}
\Crefname{examnote}{Exam note}{Exam notes}
\crefname{study}{Study note}{Study notes}
\Crefname{study}{Study note}{Study notes}
\crefname{summary}{Summary}{Summaries}
\Crefname{summary}{Summary}{Summaries}
\crefname{examproblem}{Problem}{Problems}
\Crefname{examproblem}{Problem}{Problems}

% Sectioning / floats (explicit, noabbrev-friendly)
\crefname{section}{Section}{Sections}
\Crefname{section}{Section}{Sections}
\crefname{subsection}{Subsection}{Subsections}
\Crefname{subsection}{Subsection}{Subsections}
\crefname{subsubsection}{Subsubsection}{Subsubsections}
\Crefname{subsubsection}{Subsubsection}{Subsubsections}
\crefname{paragraph}{Paragraph}{Paragraphs}
\Crefname{paragraph}{Paragraph}{Paragraphs}
\crefname{equation}{Equation}{Equations}
\Crefname{equation}{Equation}{Equations}
\crefname{figure}{Figure}{Figures}
\Crefname{figure}{Figure}{Figures}
\crefname{table}{Table}{Tables}
\Crefname{table}{Table}{Tables}
\crefname{enumi}{Item}{Items}
\Crefname{enumi}{Item}{Items}

% PDF sidebar outline + printed TOC for QE exam transcripts.
% Load in the preamble AFTER hyperref and AFTER \newcommand{\examstatement}.
\hypersetup{
  unicode=true,
  bookmarks=true,
  bookmarksopen=true,
  bookmarksopenlevel=3,
  bookmarksnumbered=false,
  bookmarksdepth=3
}
\IfFileExists{bookmark.sty}{%
  \usepackage{bookmark}%
  \bookmarksetup{open,openlevel=3,numbered=false,depth=3}%
}{}
% Unnumbered in print, but still written to the TOC / PDF outline
% down to \subsubsection. Starred headings create neither.
\setcounter{secnumdepth}{-1}
\setcounter{tocdepth}{3}
\makeatletter
\@ifundefined{examstatement}{}{%
  \let\qe@origexamstatement\examstatement
  \renewcommand{\examstatement}[1]{%
    \par\phantomsection
    \addcontentsline{toc}{subsection}{#1}%
    \qe@origexamstatement{#1}%
  }%
}
\makeatother


\title{\textbf{QUALIFYING EXAM: PROBABILITY \& STATISTICS}\\[0.5em]\Large SPRING 2025}
\date{}
\author{}

\begin{document}

\maketitle
\vspace{-3em}

\tableofcontents

\section{Instructions}
\begin{itemize}
    \item There are 11 problems in this exam (3 pages). You need to choose 8 of them to solve. If you select more than 8, only the first 8 that you have worked on will be graded. Note that 4 of the problems are worth 15 points each and the rest 10 points each.
    \item You must follow all the rules of exam taking. Misconducts will be subject to proper disciplinary actions by the Center.
    \item You must provide all necessary details for full credits. A final answer with no or little explanation/derivation, even if correct, receives a minimal credit.
\end{itemize}

\section{Statement of the problems}

\examstatement{Problem 1. [10 pts.]}
Let \( \{U_n\} \) be i.i.d.\ with uniform distribution on \( [0, 1] \). For \( n \ge 2 \), define

\[
X_n := \min\{U_1, \dots, U_n\}, \qquad Y_n = \max\{U_1, \dots, U_n\}.
\]

\begin{enumerate}[label=(\arabic*)]
    \item Derive the joint distribution of \( (X_n, Y_n) \).
    \item Calculate \( \mathbb{E}[X_n \mid \sigma(Y_n)] \).
\end{enumerate}

\examstatement{Problem 2. [10 pts.]}
Suppose \( (X, Y, Z) \) is three-dimensional mean-zero Gaussian vector with covariance matrix given by

\[
\begin{pmatrix}
1 & \rho & \tau \\
\rho & 1 & 0 \\
\tau & 0 & 1
\end{pmatrix}.
\]

\begin{enumerate}[label=(\arabic*)]
    \item Show that \( \rho^2 + \tau^2 \le 1 \).
    \item Show that there exist real numbers \( a, b, c \) and a standard Gaussian random variable \( W \) independent of \( (Y, Z) \) such that \( X = aY + bZ + cW \) and derive \( a, b, c \).
\end{enumerate}

\examstatement{Problem 3. [15 pts.]}
Suppose \( \{\xi_i\} \) are i.i.d.\ Bernoulli random variables with \( \mathbb{P}[\xi_1 = 1] = \mathbb{P}[\xi_1 = -1] = 1/2 \). Set \( S_n = \sum_{j=1}^n \xi_j \) and define \( \tau = \min\{n : S_n = -a \text{ or } b\} \) for \( a, b \in \mathbb{Z}_{>0} \).
\begin{enumerate}[label=(\arabic*)]
    \item Calculate \( \mathbb{E}[e^{\lambda S_n}] \) for \( \lambda \in \mathbb{R} \).
    \item Calculate \( \mathbb{E}[s^\tau] \) for \( s \in (0, 1) \).
\end{enumerate}

\examstatement{Problem 4. [10 pts.]}
Suppose \( \{\xi_i\} \) are i.i.d.\ Bernoulli random variables with \( \mathbb{P}[\xi_1 = 1] = \mathbb{P}[\xi_1 = -1] = 1/2 \). Set

\[
X_n = \sum_{j=1}^n 2^{j-1} \xi_j.
\]

\begin{enumerate}[label=(\arabic*)]
    \item Define \( T = \min\{j \ge 1 : \xi_j = +1\} \). Derive the law of \( T \) and calculate \( \mathbb{E}[T] \).
    \item Show that \( \{X_n\} \) is a martingale and calculate \( \mathbb{E}[X_T] \).
\end{enumerate}

\examstatement{Problem 5. [10 pts.]}
Let \( \{B_t\}_{t \ge 0} \) be a one-dimensional Brownian motion starting from the origin (i.e, \( B_0 = 0 \)). For \( a \in (0, \infty) \), define \( T_a := \inf\{t \ge 0 : |B_t| = a\} \).
\begin{enumerate}[label=(\arabic*)]
    \item Prove that \( T_a < \infty \) a.s.
    \item Prove that \( T_a \) and \( \mathbf{1}_{\{B_{T_a} = a\}} \) are independent.
\end{enumerate}

\examstatement{Problem 6. [15 pts.]}
Suppose \( X \) and \( Y \) are i.i.d with \( \mathbb{E}[X] = 0 \) and \( \mathrm{var}(X) = 1 \). Suppose further that

\[
\mathbb{E}[(X - Y)^2 \mid \sigma(X + Y)] = 2.
\]

Show that \( X \) and \( Y \) are independent standard Gaussian random variables.

\examstatement{Problem 7. [10 pts.]}
Let \( X_1, \dots, X_n \) (\( n \ge 2 \)) be iid uniform \( U(0, \theta) \) where \( 0 < \theta < \infty \).
\begin{enumerate}[label=(\alph*)]
    \item Under squared error loss \( l(\theta, \delta(X)) = (\delta(X) - \theta)^2 \), derive the generalized Bayes rule for the improper prior

    \[
    \pi(\theta) = 1, \qquad 0 < \theta < +\infty.
    \]
    \item Is the Bayes rule derived in (a) consistent for \( \theta \)?
\end{enumerate}

\examstatement{Problem 8. [10 pts.]}
Let \( X_1, \dots, X_n \) be iid Gamma\( (2025, \theta) \) with \( \theta > 0 \) unknown. Recall that the probability density function of a Gamma\( (\alpha, \theta) \) is \( f(x \mid \alpha, \theta) = \frac{1}{\Gamma(\alpha) \theta^\alpha} x^{\alpha-1} e^{-x/\theta} \), \( x > 0 \), where \( \Gamma(\alpha) = (\alpha - 1)! \) for a positve integer \( \alpha \).
\begin{enumerate}[label=(\alph*)]
    \item Consider the hypothesis testing problem \( H_0 : \theta = 1 \) vs.\ \( H_1 : \theta \neq 1 \). Derive the UMP level \( \alpha = 0.05 \) test if it exists. If it doesn't exist, please argue why it does not exist.
    \item Derive the MLE for \( \frac{1}{\theta} \), what is the asymptotic distribution of this MLE for \( \frac{1}{\theta} \)?
\end{enumerate}

\examstatement{Problem 9. [15 pts.]}
Let \( (X_1, Y_1), \dots, (X_n, Y_n) \) be iid random bivariate vectors. Let \( \rho \) be the correlation coefficient between \( X_i \) and \( Y_i \). Let

\[
\hat{\rho} = \frac{1}{(n-1)\sqrt{S_X^2 S_Y^2}} \sum_{i=1}^n (X_i - \overline{X})(Y_i - \overline{Y}),
\]

where \( \overline{X} = \frac{\sum_{i=1}^n X_i}{n} \), \( \overline{Y} = \frac{\sum_{i=1}^n Y_i}{n} \), \( S_X^2 = \frac{\sum_{i=1}^n (X_i - \overline{X})^2}{n-1} \), \( S_Y^2 = \frac{\sum_{i=1}^n (Y_i - \overline{Y})^2}{n-1} \).
\begin{enumerate}[label=(\alph*)]
    \item Assume that \( \mathbb{E}|X_i|^4 < \infty \) and \( \mathbb{E}|Y_i|^4 < \infty \), prove that \( \sqrt{n}(\hat{\rho} - \rho) \longrightarrow N(0, c) \) for some value \( c \). Please identify \( c \) in terms of moments of \( (X_i, Y_i) \), however, there is no need to simplify your calculation on \( c \).
    \item Further, assume that \( (X_i, Y_i) \) follows a bivariate normal distribution that

    \[
    \begin{pmatrix} X_i \\ Y_i \end{pmatrix}
    \sim
    N\left(
    \begin{pmatrix} 0 \\ 0 \end{pmatrix},
    \begin{pmatrix} 1 & \rho \\ \rho & 1 \end{pmatrix}
    \right).
    \]

    The probability density function of the above bivariate normal distribution is given by

    \[
    f(x, y \mid \rho) = \frac{1}{2\pi \sqrt{1 - \rho^2}} \exp\left\{ -\frac{1}{2(1 - \rho^2)} [x^2 - 2\rho xy + y^2] \right\}.
    \]

    Find a minimal sufficient statistic for \( \rho \). Is this minimal sufficient statistic complete or not? Please explain.
\end{enumerate}

\examstatement{Problem 10. [10 pts.]}
Let \( X_1, \dots, X_n \) be iid from \( N(0, \sigma^2) \) with \( \sigma > 0 \) unknown. Let \( W = \frac{1}{n} \sum_{i=1}^n |X_i| \). Recall \( \chi^2_k \) has probability density function

\[
\frac{1}{2^{k/2} \Gamma(k/2)} x^{k/2 - 1} e^{-x/2}, \qquad x \ge 0.
\]

\begin{enumerate}[label=(\alph*)]
    \item Find an unbiased estimator of \( \sigma \) based on \( W \) and identify its limiting distribution.
    \item Is it admissible under the square error loss? Justify your answer.
\end{enumerate}

\examstatement{Problem 11. [15 pts.]}
Let \( X_1, \dots, X_n \) be iid from exponential distribution with density \( f(x; \lambda, \alpha) = \lambda e^{-\lambda(x-\alpha)} I_{\{x \ge \alpha\}} \), where \( \lambda > 0 \) and \( -\infty < \alpha < \infty \) are unknown parameters.
\begin{enumerate}[label=(\alph*)]
    \item Find the MLE of \( (\lambda, \alpha) \), say \( (\hat{\lambda}, \hat{\alpha}) \).
    \item What distribution does \( \frac{\sqrt{n}(\hat{\lambda} - \lambda)}{n(\hat{\alpha} - \alpha)} \) converge to? You will receive partial credits if you can identify the distribution(s) of the numerator and/or denominator.
\end{enumerate}

\noindent\rule{\linewidth}{0.4pt}

\section{Statement and solutions}

\subsection{Problem 1. [10 pts.]}
Let \( \{U_n\} \) be i.i.d.\ with uniform distribution on \( [0, 1] \). For \( n \ge 2 \), define

\[
X_n := \min\{U_1, \dots, U_n\}, \qquad Y_n = \max\{U_1, \dots, U_n\}.
\]

\begin{enumerate}[label=(\arabic*)]
    \item Derive the joint distribution of \( (X_n, Y_n) \).
    \item Calculate \( \mathbb{E}[X_n \mid \sigma(Y_n)] \).
\end{enumerate}

\setcounter{section}{1}

\subsubsection{Definitions used in Problem 1}

Always \( X_n\le Y_n \), so the law of \( (X_n,Y_n) \) is supported on the triangle \( 0\le x\le y\le 1 \). The two objects below are the only analytic input: the mixed partial that turns a joint CDF into a joint density, and the ratio that turns that joint density into a conditional expectation given \( Y_n \).

\begin{definition}[Joint CDF and joint density]
\label{def:joint-cdf-pdf}
Let \( (X,Y) \) be a random vector in \( \mathbb{R}^{2} \). Its joint CDF is \( F(x,y)\coloneqq\Pbb(X\le x,\,Y\le y) \). If \( F \) is \( C^{2} \) on an open set \( U\subset\mathbb{R}^{2} \), the mixed partial

\[
f(x,y)
\coloneqq
\partial_{x}\partial_{y} F(x,y)
\]

is a joint density of \( (X,Y) \) on \( U \) provided \( f\ge 0 \) and \( \iint_{\mathbb{R}^{2}}f=1 \). Equivalently, if one starts from a joint survival-type function \( G(x,y)\coloneqq\Pbb(X>x,\,Y\le y) \) on a region \( \{x<y\} \), then \( f=-\partial_{x}\partial_{y}G \) on that region. In either case the law is determined by \( f \).
\end{definition}

\begin{definition}[Conditional expectation given \( \sigma(Y) \)]
\label{def:cond-exp-sigma}
If \( X \) is integrable and \( Y \) is real, there is an integrable \( \sigma(Y) \)-measurable random variable \( \E[X\mid\sigma(Y)] \), unique a.s., characterised by

\[
\E\bigl[X\,\ind_{A}\bigr]
=
\E\bigl[\E[X\mid\sigma(Y)]\,\ind_{A}\bigr]
\qquad
\text{for every }A\in\sigma(Y).
\]

Every \( \sigma(Y) \)-measurable random variable is a Borel function of \( Y \), so \( \E[X\mid\sigma(Y)]=g(Y) \) a.s.\ for some Borel \( g \). If \( (X,Y) \) admits a joint density \( f \) and \( Y \) has marginal density \( f_{Y}(y)=\int f(x,y)\,\mathrm{d}x>0 \), a version of \( g \) is

\[
g(y)
=
\E[X\mid Y=y]
=
\int x\,f_{X\mid Y}(x\mid y)\,\mathrm{d}x
=
\frac{\displaystyle\int x\,f(x,y)\,\mathrm{d}x}{f_{Y}(y)}.
\]

The denominator cannot be omitted: \( \int x\,f(x,y)\,\mathrm{d}x \) is the numerator of the conditional mean, not the conditional mean. See \citep[\S4.1]{durrett2010}.
\end{definition}

\begin{proposition}[{Minima and maxima of i.i.d.\ uniforms on \([0,1]\)}]
\label{prop:unif-minmax}
Let \( U_{1},\dots,U_{n} \) be i.i.d.\ Uniform\( [0,1] \), \( n\ge 2 \), and write \( X_{n}=\min_{i}U_{i} \), \( Y_{n}=\max_{i}U_{i} \). Then, for \( 0\le a\le b\le 1 \),

\[
\Pbb(X_{n}\ge a,\,Y_{n}\le b)
=
(b-a)^{n}.
\]

Consequently, on the triangle \( 0\le x\le y\le 1 \),

\[
F_{X_{n},Y_{n}}(x,y)
=
y^{n}-(y-x)^{n},
\]

and \( (X_{n},Y_{n}) \) is absolutely continuous on \( \{0<x<y<1\} \) with joint density

\[
f_{X_{n},Y_{n}}(x,y)
=
n(n-1)(y-x)^{n-2}.
\]

The marginal of the maximum is \( f_{Y_{n}}(y)=n y^{n-1} \) on \( (0,1) \). This is the \( (Y_{1},Y_{n}) \) specialisation of the order-statistic density in \citep[\S4.4]{hogg2019}: in general \( n(n-1)\bigl(F(y)-F(x)\bigr)^{n-2}f(x)f(y) \), which collapses to the display above when \( F(t)=t \) on \( [0,1] \).
\end{proposition}

\begin{proof}
The event \( \{X_{n}\ge a,\,Y_{n}\le b\} \) is \( \{U_{1},\dots,U_{n}\in[a,b]\} \). Independence and the Uniform\( [0,1] \) law give \( (b-a)^{n} \). Taking \( a=0 \) yields \( \Pbb(Y_{n}\le y)=y^{n} \) for \( y\in[0,1] \), hence the marginal density \( n y^{n-1} \). On \( 0\le x\le y\le 1 \),

\[
F_{X_{n},Y_{n}}(x,y)
=
\Pbb(Y_{n}\le y)-\Pbb(X_{n}>x,\,Y_{n}\le y)
=
y^{n}-(y-x)^{n}.
\]

(If \( x\ge y \) one has \( F_{X_{n},Y_{n}}(x,y)=\Pbb(Y_{n}\le y) \) instead, since \( X_{n}\le Y_{n} \) always.) Differentiating the triangle formula as in \cref{def:joint-cdf-pdf},

\[
\partial_{y}F
=
n y^{n-1}-n(y-x)^{n-1},
\qquad
\partial_{x}\partial_{y}F
=
n(n-1)(y-x)^{n-2}.
\]

The mixed partial vanishes off \( \{x<y\} \), and integrates to \( 1 \) over the triangle (substitute \( u=y-x \), or quote \( \Pbb(0\le X_{n}\le Y_{n}\le 1)=1 \)).
\end{proof}

\begin{lemma}[The elementary integral \( I_{k} \)]
\label{lem:Ik}
For \( k\ge 0 \) and \( y>0 \),

\[
I_{k}(y)
\coloneqq
\int_{0}^{y} x(y-x)^{k}\,\mathrm{d}x
=
\frac{y^{k+2}}{(k+1)(k+2)}.
\]
\end{lemma}

\begin{proof}
Substitute \( u=y-x \), so \( x=y-u \) and \( \mathrm{d}x=-\mathrm{d}u \). Then

\[
I_{k}(y)
=
\int_{0}^{y}(y-u)u^{k}\,\mathrm{d}u
=
y\cdot\frac{y^{k+1}}{k+1}-\frac{y^{k+2}}{k+2}
=
\frac{y^{k+2}}{(k+1)(k+2)}.
\]
\end{proof}

The same substitution is the one started in the notes (\( x=y-(y-x) \)). It already closes; no recursion on \( k \) is required.

\subsubsection{Manuscript proof idea}

The calculation below follows the handwritten notes, in the same order. \mserr{Red} marks a false claim or a missing justification. The opening idea --- write the joint CDF on the triangle, differentiate to a joint density, then compute a function of \( y \) --- is the right exam skeleton. The three errors are: both factors in the CDF, the missing marginal in the conditional mean, and the evaluation of \( I_{k} \).

The notes work on \( 0\le x\le y\le 1 \), which is the correct support, and factor

\[
\Pbb(X_{n}\le x,\,Y_{n}\le y)
=
\Pbb(Y_{n}\le y)\,\Pbb(X_{n}\le x\mid Y_{n}\le y)
=
\Pbb(Y_{n}\le y)\bigl(1-\Pbb(X_{n}>x\mid Y_{n}\le y)\bigr).
\]

The factorisation itself is tautological whenever \( \Pbb(Y_{n}\le y)>0 \). The notes then write

\[
\Pbb(X_{n}\le x,\,Y_{n}\le y)
=
\mserr{y\cdot\bigl[1-(y-x)^{n}\bigr]}\,\mathbb{I}(0\le x\le y\le 1).
\]

\mserr{False: \(\Pbb(Y_{n}\le y)=y^{n}\), not \( y \). The formula \( y \) is the CDF of a single Uniform\( [0,1] \), i.e.\ of any one \( U_{i} \), not of the maximum of \( n \) copies.} Independently, \mserr{\(\Pbb(X_{n}>x\mid Y_{n}\le y)=(y-x)^{n}\) is false}: the right-hand side is the \emph{unconditional} probability \( \Pbb(X_{n}>x,\,Y_{n}\le y)=(y-x)^{n} \), not the conditional. Given \( Y_{n}\le y \), the sample is i.i.d.\ Uniform\( [0,y] \), so

\[
\Pbb(X_{n}>x\mid Y_{n}\le y)
=
\Bigl(\frac{y-x}{y}\Bigr)^{n},
\qquad
0\le x\le y\le 1.
\]

The product of the two correct factors is \( y^{n}\bigl(1-((y-x)/y)^{n}\bigr)=y^{n}-(y-x)^{n} \), which is \cref{prop:unif-minmax}. The notes mixed the \( n=1 \) version of the first factor with the unconditional version of the second. (Conditioning on the event \( \{Y_{n}\le y\} \) is also not the same as conditioning on \( \sigma(Y_{n}) \); it is only a device for writing the joint CDF. The object in (2) needs the regular conditional law given \( Y_{n}=y \), as in \cref{def:cond-exp-sigma}.)

The notes differentiate their (incorrect) CDF as

\[
p(x,y)
=
\partial_{x}\partial_{y}\bigl[y-y(y-x)^{n}\bigr]
=
n(y-x)^{n-1}+n(n-1)y(y-x)^{n-2}.
\]

The mixed partial of the displayed expression is algebraically correct. \mserr{The output is not the joint density}: the extra summand \( n(y-x)^{n-1} \) is the \( x \)-derivative of the spurious linear factor \( y \) in the first term of the CDF, and the second summand still carries an extra factor \( y \). The legal mixed partial of \( y^{n}-(y-x)^{n} \) is \( n(n-1)(y-x)^{n-2} \), as in \cref{prop:unif-minmax}. Equivalently, start from \( G(x,y)=(y-x)^{n}=\Pbb(X_{n}>x,\,Y_{n}\le y) \) and compute \( f=-\partial_{x}\partial_{y}G \); that bypasses the factorisation entirely.

For (2) the notes write

\[
\E(X_{n}\mid Y_{n}=y)
=
\mserr{\int_{0}^{y} x\,p(x,y)\,\mathrm{d}x}
=
\int_{0}^{y}\bigl[nx(y-x)^{n-1}+n(n-1)yx(y-x)^{n-2}\bigr]\,\mathrm{d}x.
\]

\mserr{False: this is the numerator of \cref{def:cond-exp-sigma}, not the conditional mean.} Even with the correct joint density one must divide by \( f_{Y_{n}}(y)=ny^{n-1} \). The displayed integral against the notes' \( p \) is not \( \E[X_{n}\mid Y_{n}=y] \) for a second, independent reason: \( p \) itself is not a joint density.

The notes set \( I_{k}(y)\coloneqq\int_{0}^{y}x(y-x)^{k}\,\mathrm{d}x \) and rewrite the (already incorrect) integral as \( n I_{n-1}+n(n-1)y I_{n-2} \). That algebra is consistent with their \( p \). They then claim the recursion

\[
\mserr{\frac{I_{k}(y)}{k!}
=
\frac{y^{k+1}}{(k+1)!}+\frac{I_{k-1}(y)}{(k-1)!}}
\qquad\Longrightarrow\qquad
\mserr{\frac{I_{k}(y)}{k!}
=
\sum_{j=0}^{k}\frac{y^{j+1}}{(j+1)!}}.
\]

(The notes used the letter \( n \) as the summation index; it has been renamed \( j \) to avoid clashing with the sample size.) \mserr{False, already at the base: \( I_{0}(y)=\int_{0}^{y}x\,\mathrm{d}x=y^{2}/2 \), so \( I_{0}/0!=y^{2}/2 \), not \( y \).} The two sides of the claimed recursion are not homogeneous in \( y \): \( I_{k} \) scales as \( y^{k+2} \), while \( y^{k+1} \) scales as \( y^{k+1} \). The substitution \( x=y-(y-x) \) started in the notes does not produce a recursion on \( k \); it produces the closed form in \cref{lem:Ik} in one line.

Substituting the (false) series into \( n I_{n-1}+n(n-1)y I_{n-2} \) and rearranging, the notes reach an expression of the shape

\[
n!\sum_{k=1}^{n-1}\Bigl(\mserr{\frac{1}{(k-1)!}}-\frac{1}{k!}\Bigr)y^{k+1}+n!\,y,
\]

and conclude

\[
\E\bigl[X_{n}\mid\sigma(Y_{n})\bigr]
=
\mserr{n!\sum_{k=0}^{n-1}\frac{1}{(k-1)!}\,Y_{n}^{k+1}}.
\]

\mserr{The \( k=0 \) term contains \( (-1)! \), which is undefined.} Independently the claimed sum is a degree-\( n \) polynomial in \( Y_{n} \), whereas the conditional mean is linear. Completing the notes' own integral with the legal \( I_{k} \) from \cref{lem:Ik}, without yet repairing the density or inserting the marginal, gives

\[
n I_{n-1}(y)+n(n-1)y I_{n-2}(y)
=
\frac{y^{n+1}}{n+1}+y^{n+1}
=
\frac{n+2}{n+1}\,y^{n+1},
\]

which is \( \int_{0}^{y}x\,p(x,y)\,\mathrm{d}x \) for the notes' \( p \), not a version of \( \E[X_{n}\mid Y_{n}=y] \). Dividing that quantity by \( \int_{0}^{y}p(x,y)\,\mathrm{d}x=(n+1)y^{n} \) would produce \( \frac{n+2}{(n+1)^{2}}y \), still not \( y/n \). The CDF error and the missing marginal are not cancelled by evaluating \( I_{k} \) correctly.

\errpanel{%
\textbf{\mserr{The error, in one box.}}

The factorisation of the joint CDF is the right start, but both factors were evaluated as if \( n=1 \) or as if a joint probability were a conditional probability. The legal CDF on the triangle is \( y^{n}-(y-x)^{n} \), and the legal density is \( n(n-1)(y-x)^{n-2} \).

The map \( y\mapsto\int x\,p(x,y)\,\mathrm{d}x \) is not \( \E[X_{n}\mid Y_{n}=y] \). One must divide by \( f_{Y_{n}}(y)=ny^{n-1} \).

The integral \( I_{k} \) is elementary: \( I_{k}(y)=y^{k+2}/[(k+1)(k+2)] \). The factorial series, the term \( (-1)! \), and the final polynomial in \( Y_{n} \) are not needed and are not correct. The answer is \( \E[X_{n}\mid\sigma(Y_{n})]=Y_{n}/n \).
}

\subsubsection{Solution of Problem 1}

\paragraph{Answer.}
On \( \{0<x<y<1\} \), the pair \( (X_{n},Y_{n}) \) has joint density \( n(n-1)(y-x)^{n-2} \). A version of the conditional expectation is \( \E[X_{n}\mid\sigma(Y_{n})]=Y_{n}/n \).

\begin{proof}[Correct proof]
\textbf{(1).}
For \( 0\le a\le b\le 1 \), the event \( \{X_{n}\ge a,\,Y_{n}\le b\} \) is \( \{U_{1},\dots,U_{n}\in[a,b]\} \), hence has probability \( (b-a)^{n} \). In particular \( \Pbb(Y_{n}\le y)=y^{n} \) for \( y\in[0,1] \). Therefore, on the triangle \( 0\le x\le y\le 1 \),

\[
F_{X_{n},Y_{n}}(x,y)
=
\Pbb(X_{n}\le x,\,Y_{n}\le y)
=
y^{n}-(y-x)^{n}.
\]

(The formula is the notes' factorisation, with \( \Pbb(Y_{n}\le y)=y^{n} \) and \( \Pbb(X_{n}>x\mid Y_{n}\le y)=((y-x)/y)^{n} \).) Differentiating as in \cref{def:joint-cdf-pdf},

\[
\partial_{y}F_{X_{n},Y_{n}}(x,y)
=
n y^{n-1}-n(y-x)^{n-1},
\qquad
f_{X_{n},Y_{n}}(x,y)
=
\partial_{x}\partial_{y}F_{X_{n},Y_{n}}(x,y)
=
n(n-1)(y-x)^{n-2},
\]

for \( 0<x<y<1 \), and \( f_{X_{n},Y_{n}}=0 \) off that triangle. This is the joint distribution: the law is absolutely continuous on \( \{0<x<y<1\} \) with the displayed density. (Shorter still: \( G(x,y)\coloneqq\Pbb(X_{n}>x,\,Y_{n}\le y)=(y-x)^{n} \), and \( f=-\partial_{x}\partial_{y}G \).)

\textbf{(2).}
The marginal density of \( Y_{n} \) is \( f_{Y_{n}}(y)=n y^{n-1} \) on \( (0,1) \). By \cref{def:cond-exp-sigma}, for \( y\in(0,1) \),

\[
\E[X_{n}\mid Y_{n}=y]
=
\frac{1}{n y^{n-1}}\int_{0}^{y} x\cdot n(n-1)(y-x)^{n-2}\,\mathrm{d}x
=
\frac{n-1}{y^{n-1}}\,I_{n-2}(y).
\]

\Cref{lem:Ik} gives \( I_{n-2}(y)=y^{n}/[(n-1)n] \), hence

\[
\E[X_{n}\mid Y_{n}=y]
=
\frac{n-1}{y^{n-1}}\cdot\frac{y^{n}}{(n-1)n}
=
\frac{y}{n}.
\]

Thus \( \E[X_{n}\mid\sigma(Y_{n})]=Y_{n}/n \).
\end{proof}

\begin{remark}[The representation that makes the answer visible]
\label{rem:unif-given-max}
Given \( Y_{n}=y \), one coordinate equals \( y \) and the remaining \( n-1 \) coordinates are i.i.d.\ Uniform\( [0,y] \) (this is the conditional law of i.i.d.\ uniforms given their maximum; equivalently, it is the statement that \( f_{X_{n}\mid Y_{n}}(x\mid y)=(n-1)(y-x)^{n-2}/y^{n-1} \) on \( (0,y) \)). Then \( X_{n} \) is the minimum of those \( n-1 \) copies, whose mean is \( y/n \). This is a legitimate exam shortcut once the conditional law is named; the proof above does not need it, because the ratio of densities already produces \( y/n \).
\end{remark}

\subsection{Problem 2. [10 pts.]}
Suppose \( (X, Y, Z) \) is three-dimensional mean-zero Gaussian vector with covariance matrix given by

\[
\begin{pmatrix}
1 & \rho & \tau \\
\rho & 1 & 0 \\
\tau & 0 & 1
\end{pmatrix}.
\]

\begin{enumerate}[label=(\arabic*)]
    \item Show that \( \rho^2 + \tau^2 \le 1 \).
    \item Show that there exist real numbers \( a, b, c \) and a standard Gaussian random variable \( W \) independent of \( (Y, Z) \) such that \( X = aY + bZ + cW \) and derive \( a, b, c \).
\end{enumerate}

\setcounter{section}{2}

\subsubsection{Hint for Problem 2(1)}

Gaussianity is not used. Any random vector with this covariance matrix \(\Sigma\) already forces the inequality, because \(\Sigma\) is a Gram matrix: for every \( v\in\mathbb{R}^{3} \),

\[
v^{T}\Sigma v
=
\Var\bigl(v^{T}(X,Y,Z)^{T}\bigr)
\ge 0.
\]

The claim is this quadratic form at one vector. The natural choice is the residual of \( X \) after subtracting its covariances against \( Y \) and \( Z \), namely \( v=(1,-\rho,-\tau)^{T} \), or equivalently \( X-\rho Y-\tau Z \). Expand the variance using the given entries, in particular \(\operatorname{Cov}(Y,Z)=0\):

\[
\begin{aligned}
\Var(X-\rho Y-\tau Z)
&=
\Var(X)+\rho^{2}\Var(Y)+\tau^{2}\Var(Z)\\
&\qquad
-2\rho\operatorname{Cov}(X,Y)
-2\tau\operatorname{Cov}(X,Z)
+2\rho\tau\operatorname{Cov}(Y,Z)\\
&=
1+\rho^{2}+\tau^{2}-2\rho^{2}-2\tau^{2}
=
1-\rho^{2}-\tau^{2}.
\end{aligned}
\]

Nonnegativity is \( \rho^{2}+\tau^{2}\le 1 \). The same quantity is \( \det\Sigma \), and is the \( c^{2} \) of part~(2). (Sylvester's criterion on principal minors also works: the three \( 1\times 1 \) minors are \( 1 \), the \( 2\times 2 \) minors are \( 1-\rho^{2} \) and \( 1-\tau^{2} \), and the full determinant is \( 1-\rho^{2}-\tau^{2} \); the last is the strongest.)

\subsubsection{Definitions used in Problem 2}

The covariance matrix already records that \(\operatorname{Cov}(Y,Z)=0\). The slogan ``uncorrelated Gaussians are independent'' is legal here, but only because \((X,Y,Z)\) was given as a Gaussian \emph{vector}. Characteristic functions prove the slogan; they do not let one drop the joint-Gaussian hypothesis. That hypothesis is precisely the second-order fact that a Gaussian law is determined by its mean and covariance.

\begin{definition}[Gaussian random vector]
\label{def:gaussian-vector}
A random vector \( V\in\mathbb{R}^{d} \) is \emph{Gaussian} (possibly degenerate) if \( a^{T}V \) is univariate Gaussian for every \( a\in\mathbb{R}^{d} \). Equivalently, there exist \( \bm{\mu}\in\mathbb{R}^{d} \) and a symmetric positive semidefinite \( \Sigma\in\mathbb{R}^{d\times d} \) such that

\[
\E\bigl[e^{it^{T}V}\bigr]
=
\exp\Bigl(it^{T}\bm{\mu}-\tfrac12 t^{T}\Sigma t\Bigr),
\qquad
t\in\mathbb{R}^{d}.
\]

Then \( \E[V]=\bm{\mu} \) and \( \operatorname{Cov}(V)=\Sigma \). Linear images of Gaussian vectors are Gaussian: if \( A \) is \( k\times d \), then \( AV\sim\mathcal{N}(A\bm{\mu},\,A\Sigma A^{T}) \). In particular every subvector is Gaussian. If \( \Sigma \) is positive definite, \( V \) has Lebesgue density

\[
(2\pi)^{-d/2}(\det\Sigma)^{-1/2}
\exp\Bigl(-\tfrac12(v-\bm{\mu})^{T}\Sigma^{-1}(v-\bm{\mu})\Bigr)
\]

on \( \mathbb{R}^{d} \); if \( \Sigma \) is only semidefinite the law is supported on an affine subspace and need not be absolutely continuous. See \citep[\S3.3, \S3.9]{durrett2010}.
\end{definition}

\begin{proposition}[Independence via characteristic functions]
\label{prop:chf-indep}
Real random variables \( Y,Z \) are independent if and only if

\[
\E\bigl[e^{isY+itZ}\bigr]
=
\E\bigl[e^{isY}\bigr]\,\E\bigl[e^{itZ}\bigr]
\qquad
\text{for all }s,t\in\mathbb{R}.
\]

More generally, the coordinates of a random vector \( V \) are mutually independent if and only if the joint characteristic function factors as a product of the one-dimensional characteristic functions. This is the Fourier form of \( \mathrm{Law}(V)=\bigotimes_{j}\mathrm{Law}(V_{j}) \). See \citep[Theorem~3.3.2]{durrett2010} and \citep[Theorem~8.4.4]{ren2024}.
\end{proposition}

\begin{theorem}[Uncorrelated Gaussian coordinates are independent]
\label{thm:gauss-uncorr-indep}
Let \( V=(V_{1},\dots,V_{d})^{T} \) be a Gaussian vector, in the sense of \cref{def:gaussian-vector}. The following are equivalent:
\begin{enumerate}[label=\textup{(\roman*)}]
    \item
    \( \operatorname{Cov}(V_{j},V_{k})=0 \) for all \( j\neq k \), i.e.\ \( \operatorname{Cov}(V) \) is diagonal;
    \item
    \( V_{1},\dots,V_{d} \) are mutually independent.
\end{enumerate}
In particular, if \( (Y,Z) \) is a two-dimensional Gaussian vector, then \( Y \) and \( Z \) are independent if and only if \( \operatorname{Cov}(Y,Z)=0 \). This is \citep[Theorem~6.2.1]{ren2024}, stated for vectors rather than for a bivariate density.
\end{theorem}

\begin{proof}
Independence always implies uncorrelatedness whenever second moments exist, so (ii)\(\Rightarrow\)(i) needs no Gaussianity. For the converse, centre \( V \) (independence is unaffected by deterministic shifts) and write \( \Sigma=\operatorname{diag}(\sigma_{1}^{2},\dots,\sigma_{d}^{2}) \). \Cref{def:gaussian-vector} gives, for every \( t=(t_{1},\dots,t_{d})^{T} \),

\[
\E\bigl[e^{it^{T}V}\bigr]
=
\exp\Bigl(-\tfrac12 t^{T}\Sigma t\Bigr)
=
\exp\Bigl(-\tfrac12\sum_{j=1}^{d}\sigma_{j}^{2}t_{j}^{2}\Bigr)
=
\prod_{j=1}^{d}\exp\Bigl(-\tfrac12\sigma_{j}^{2}t_{j}^{2}\Bigr)
=
\prod_{j=1}^{d}\E\bigl[e^{it_{j}V_{j}}\bigr].
\]

\Cref{prop:chf-indep} yields (ii). The two-dimensional case is \( d=2 \).
\end{proof}

The characteristic-function argument does not avoid the ``second-order Gaussian'' input: it \emph{is} that input. The joint characteristic function of a Gaussian vector is an exponential of a quadratic form in the mean and covariance, and nothing else. Setting the off-diagonal of \( \Sigma \) to zero makes the quadratic form a sum, hence the exponential a product. One cannot run the same computation from the marginal characteristic functions \( \E[e^{isY}] \) and \( \E[e^{itZ}] \) together with the scalar \( \operatorname{Cov}(Y,Z)=0 \): those data do not determine \( \E[e^{isY+itZ}] \).

The density proof in \citep[Theorem~6.2.1]{ren2024} is the same theorem under a nondegeneracy assumption. If \( (Y,Z)\sim\mathcal{N}(0,\Sigma) \) with \( \det\Sigma>0 \), the joint density is proportional to \( \exp\bigl(-\tfrac12(y,z)\Sigma^{-1}(y,z)^{T}\bigr) \). When \( \operatorname{Cov}(Y,Z)=0 \) one has \( \Sigma^{-1} \) diagonal, so the density factors as \( f_{Y}(y)f_{Z}(z) \). That argument uses the bivariate normal density, therefore requires \( \Sigma \) invertible, and fails on the boundary \( \rho^{2}+\tau^{2}=1 \) of part~(1), where the three-dimensional law of \( (X,Y,Z) \) is supported on a plane. The characteristic-function proof covers the degenerate case, which is why it is the one to write on the exam.

\begin{warning}[Marginal Gaussians plus a vanishing covariance are not enough]
\label{warn:marginal-gauss-not-joint}
The hypothesis in \cref{thm:gauss-uncorr-indep} is that \( (Y,Z) \) is a Gaussian \emph{vector}, not that \( Y \) and \( Z \) are separately Gaussian. Two independent Gaussians automatically form a Gaussian vector (the joint characteristic function is the product of the one-dimensional Gaussians, hence of the form in \cref{def:gaussian-vector}). Two uncorrelated Gaussians need not: \citep[after Theorem~6.2.1]{ren2024}.
\end{warning}

\begin{example}[Sign-flip: uncorrelated standard Gaussians that are not independent]
\label{ex:signflip}
Let \( Y\sim\mathcal{N}(0,1) \) and let \( \varepsilon \) be independent of \( Y \) with \( \Pbb(\varepsilon=\pm 1)=\tfrac12 \). Set \( Z\coloneqq\varepsilon Y \). Then \( Z\sim\mathcal{N}(0,1) \) by symmetry, and

\[
\operatorname{Cov}(Y,Z)
=
\E[\varepsilon Y^{2}]
=
\E[\varepsilon]\,\E[Y^{2}]
=
0.
\]

They are not independent: \( |Y|=|Z| \) almost surely. The pair is not a Gaussian vector: \( Y+Z=Y(1+\varepsilon) \) equals \( 0 \) or \( 2Y \) each with probability \( \tfrac12 \), which is not Gaussian. The joint characteristic function can be computed directly and does not factor:

\[
\E\bigl[e^{isY+itZ}\bigr]
=
\E\bigl[\exp\bigl(-\tfrac12(s+t\varepsilon)^{2}\bigr)\bigr]
=
\exp\bigl(-\tfrac12(s^{2}+t^{2})\bigr)\cosh(st).
\]

The product of the marginals is \( \exp\bigl(-\tfrac12(s^{2}+t^{2})\bigr) \). These agree for all \( (s,t) \) if and only if \( \cosh(st)=1 \) for all \( (s,t) \), which is false. Characteristic functions detect the failure; they do not repair it. A density-level variant is \citep[Example~3, \S6.2]{ren2024}: if \( \varphi \) is the standard Gaussian density and \( g \) is a continuous even function with \( |g|\le\varphi \) and \( \int g=0 \), then \( p(x,y)=\varphi(x)\varphi(y)+g(x)g(y) \) has Gaussian margins and vanishing covariance, but \( p \) does not factor.
\end{example}

In the present problem \( (X,Y,Z) \) is a Gaussian vector by hypothesis, so the subvector \( (Y,Z) \) is a Gaussian vector by \cref{def:gaussian-vector}, and the given covariance matrix has \( \operatorname{Cov}(Y,Z)=0 \). \Cref{thm:gauss-uncorr-indep} therefore yields \( Y\perp Z \). The same theorem is the independence statement in part~(2): once \( W \) is constructed as a linear image of \( (X,Y,Z) \), the triple \( (W,Y,Z) \) is Gaussian, and \( W\perp(Y,Z) \) as soon as \( \operatorname{Cov}(W,Y)=\operatorname{Cov}(W,Z)=0 \). Pairwise vanishing covariances among the three coordinates would give mutual independence of \( W,Y,Z \); independence of \( W \) from the pair \( (Y,Z) \) is the slightly weaker block-diagonal statement \( \operatorname{Cov}(W,(Y,Z))=0 \), which is what the exam asks.

\subsubsection{Manuscript proof idea}

The calculation below follows the handwritten notes, in the same order. The notes label the two items (a) and (b); these are the exam's (1) and (2). \mserr{Red} marks a false claim or a missing justification. The residual written in (b) is the right object; the Cauchy--Schwarz attempts in (a) do not reach \( \rho^{2}+\tau^{2}\le 1 \), and (b) stops before checking that the residual is standard Gaussian and uncorrelated with \( (Y,Z) \).

The notes open correctly:
\[
\Var(X)=\Var(Y)=\Var(Z)=1
\qquad\Longrightarrow\qquad
\E[X^{2}]=\E[Y^{2}]=\E[Z^{2}]=1,
\]
together with \( \operatorname{Cov}(X,Y)=\rho \), \( \operatorname{Cov}(X,Z)=\tau \), and \( \operatorname{Cov}(Y,Z)=0 \). They then write
\[
\operatorname{Cov}(Y,Z)=0
\qquad\Longrightarrow\qquad
\mserr{Y\perp Z}.
\]
The implication is true in this problem, by \cref{thm:gauss-uncorr-indep}, because \( (Y,Z) \) is a subvector of a Gaussian vector. \mserr{Gap: that joint-Gaussian hypothesis was not named.} The slogan is illegal from the margins alone; see \cref{warn:marginal-gauss-not-joint,ex:signflip}. (Independence of \( Y \) and \( Z \) is convenient for part~(2), because it kills the cross term in \( \Var(\rho Y+\tau Z) \). It is not needed for part~(1).)

The notes then try to bound \( \rho^{2}+\tau^{2} \) by Cauchy--Schwarz on each pair:
\[
\bigl(\E[XY]\bigr)^{2}
\le
\E[X^{2}]\,\E[Y^{2}]
=
1,
\]
and likewise \( \tau^{2}\le 1 \). Both estimates are correct. \mserr{They do not imply \( \rho^{2}+\tau^{2}\le 1 \).} A crossed-out comparison of the shape \( \mserr{a^{2}+b^{2}\le \E(a+b)^{2}} \) is false in general (take \( a=1 \), \( b=-1 \)).

A second rewriting on the notes is correct as an identity:
\[
\rho^{2}+\tau^{2}
=
\tfrac12\bigl(\E[X(Y+Z)]\bigr)^{2}
+
\tfrac12\bigl(\E[X(Y-Z)]\bigr)^{2},
\]
since \( \E[X(Y+Z)]=\rho+\tau \) and \( \E[X(Y-Z)]=\rho-\tau \). Applying Cauchy--Schwarz to each summand produces
\[
\bigl(\E[X(Y+Z)]\bigr)^{2}
\le
\E[(Y+Z)^{2}]
=
2,
\]
hence each half is at most \( 1 \), and the sum is at most \( 2 \). \mserr{The bound is \( \rho^{2}+\tau^{2}\le 2 \), not \( \le 1 \).} The \( Y\pm Z \) splitting does not close part~(1).

The notes then draw \( X \) in the plane of \( Y,Z \) plus an orthogonal direction \( W \), and write
\[
\mserr{\Var(X)=a^{2}+b^{2}+c^{2}=1},
\qquad
\operatorname{Cov}(X,Y)=a=\rho,
\qquad
\operatorname{Cov}(X,Z)=b=\tau,
\]
and \( c=\sqrt{1-\rho^{2}-\tau^{2}} \).
The coefficients \( a=\rho \), \( b=\tau \) are the right ones, and the picture is the geometry of part~(2). \mserr{The Pythagoras \( a^{2}+b^{2}+c^{2}=1 \) assumes an \( L^{2} \)-orthonormal triple \( (Y,Z,W) \) already in hand, which is the conclusion of (2) and uses \( c^{2}\ge 0 \), i.e.\ part~(1).} Using the representation to prove the inequality is circular. The legal order is the hint above: compute \( \Var(X-\rho Y-\tau Z)=1-\rho^{2}-\tau^{2}\ge 0 \) first, then normalise.

For (b) the notes define
\[
W
\coloneqq
\frac{X-\rho Y-\tau Z}{\sqrt{1-\rho^{2}-\tau^{2}}}
\]
and conclude \( \mserr{W\perp(Y,Z)} \). The formula is the right residual, once part~(1) has put a nonnegative number under the square root. \mserr{Three gaps.} First, if \( \rho^{2}+\tau^{2}=1 \) the denominator vanishes; then \( X=\rho Y+\tau Z \) almost surely and one takes \( c=0 \) with any standard Gaussian \( W \) independent of \( (Y,Z) \). Second, \( W\sim\mathcal{N}(0,1) \) and \( \operatorname{Cov}(W,Y)=\operatorname{Cov}(W,Z)=0 \) were not checked. Third, uncorrelatedness of \( W \) with \( Y \) and with \( Z \) yields \( W\perp(Y,Z) \) only after naming that \( (W,Y,Z) \) is a Gaussian vector (\cref{def:gaussian-vector,thm:gauss-uncorr-indep}).

\errpanel{%
\textbf{\mserr{The error, in one box.}}

Cauchy--Schwarz on \( (X,Y) \) and \( (X,Z) \) separately, or on \( (X,Y\pm Z) \), does not give \( \rho^{2}+\tau^{2}\le 1 \). The inequality is \( \Var(X-\rho Y-\tau Z)\ge 0 \).

The orthonormal expansion \( X=\rho Y+\tau Z+cW \) is the answer to (2), not a proof of (1).

The displayed \( W \) is correct for \( c>0 \), but one must still compute \( \Var(W)=1 \), \( \operatorname{Cov}(W,Y)=\operatorname{Cov}(W,Z)=0 \), invoke joint Gaussianity of \( (W,Y,Z) \), and treat \( c=0 \) separately. The coefficients are \( a=\rho \), \( b=\tau \), \( c=\sqrt{1-\rho^{2}-\tau^{2}} \).
}

\subsubsection{Solution of Problem 2}

\paragraph{Answer.}
\( \rho^{2}+\tau^{2}\le 1 \). The representation holds with \( a=\rho \), \( b=\tau \), \( c=\sqrt{1-\rho^{2}-\tau^{2}} \), and
\[
W
\coloneqq
\begin{cases}
\dfrac{X-\rho Y-\tau Z}{c}, & c>0,\\[0.6em]
\text{any }\mathcal{N}(0,1)\text{ independent of }(Y,Z), & c=0.
\end{cases}
\]

\begin{proof}[Correct proof]
\textbf{(1).}
As in the hint: \( \Var(X-\rho Y-\tau Z)=1-\rho^{2}-\tau^{2}\ge 0 \).

\textbf{(2).}
Write \( R\coloneqq X-\rho Y-\tau Z \), so \( \Var(R)=c^{2} \) with \( c=\sqrt{1-\rho^{2}-\tau^{2}} \) by (1). The subvector \( (Y,Z) \) is Gaussian and \( \operatorname{Cov}(Y,Z)=0 \), hence \( Y\perp Z \) by \cref{thm:gauss-uncorr-indep}. In particular
\[
\operatorname{Cov}(R,Y)
=
\rho-\rho\cdot 1-\tau\cdot 0
=
0,
\qquad
\operatorname{Cov}(R,Z)
=
\tau-\rho\cdot 0-\tau\cdot 1
=
0.
\]

If \( c=0 \), then \( R=0 \) a.s., so \( X=\rho Y+\tau Z \). Take any \( W\sim\mathcal{N}(0,1) \) independent of \( (Y,Z) \). Then \( X=aY+bZ+cW \) holds with \( a=\rho \), \( b=\tau \), \( c=0 \).

If \( c>0 \), set \( W\coloneqq R/c \). Then \( W \) is a linear image of the Gaussian vector \( (X,Y,Z) \), hence Gaussian, with mean \( 0 \) and variance \( 1 \). The triple \( (W,Y,Z) \) is likewise a linear image, hence a Gaussian vector, and
\[
\operatorname{Cov}(W,Y)
=
\frac{\operatorname{Cov}(R,Y)}{c}
=
0,
\qquad
\operatorname{Cov}(W,Z)
=
0.
\]
\Cref{thm:gauss-uncorr-indep} (or the block-diagonal form of the joint characteristic function) gives \( W\perp(Y,Z) \). Rearranging the definition of \( W \) yields \( X=\rho Y+\tau Z+cW \).
\end{proof}

\subsection{Problem 3. [15 pts.]}
Suppose \( \{\xi_i\} \) are i.i.d.\ Bernoulli random variables with \( \mathbb{P}[\xi_1 = 1] = \mathbb{P}[\xi_1 = -1] = 1/2 \). Set \( S_n = \sum_{j=1}^n \xi_j \) and define \( \tau = \min\{n : S_n = -a \text{ or } b\} \) for \( a, b \in \mathbb{Z}_{>0} \).
\begin{enumerate}[label=(\arabic*)]
    \item Calculate \( \mathbb{E}[e^{\lambda S_n}] \) for \( \lambda \in \mathbb{R} \).
    \item Calculate \( \mathbb{E}[s^\tau] \) for \( s \in (0, 1) \).
\end{enumerate}

\setcounter{section}{3}

\subsubsection{Definitions used in Problem 3}

The walk \( S_n=\xi_1+\cdots+\xi_n \) is simple symmetric random walk on \( \mathbb{Z} \), started at \( 0 \). Part~(1) is the moment generating function of a fixed-time position. Part~(2) is the probability generating function of a \emph{random} time: the first exit from \( (-a,b) \).

The readable order for (2) is: say what \( \E[s^{\tau}] \) is; enlarge to a function \( u(x) \) of the starting point; write the one-step relation (the \( s \) is the step just used); solve the recurrence; read \( u(0) \). Optional stopping is the same linear system, obtained from the martingale of part~(1) instead of from the recurrence. Do not begin by listing paths of length \( n \).

\begin{definition}[Discrete stopping time]
\label{def:stopping-time-discrete}
Let \( \mathcal{F}_n=\sigma(\xi_1,\dots,\xi_n) \) (and \( \mathcal{F}_0 \) trivial). A random variable \( \tau \) taking values in \( \{0,1,2,\dots\}\cup\{\infty\} \) is a \emph{stopping time} if \( \{\tau=n\}\in\mathcal{F}_n \) for every finite \( n \): the decision to stop at time \( n \) uses only the path up to \( n \). Equivalently \( \{\tau\le n\}\in\mathcal{F}_n \). First hitting and first exit times of subsets of \( \mathbb{Z} \) are stopping times. See \citep[\S5.2]{durrett2010}.
\end{definition}

\begin{definition}[Two-barrier exit]
\label{def:two-barrier-exit}
For \( a,b\in\mathbb{Z}_{>0} \),
\[
\tau
\coloneqq
\min\{n\ge 0: S_n\in\{-a,b\}\},
\]
with the convention \( \min\emptyset=\infty \). This is a stopping time. On \( \{\tau<\infty\} \) one has \( S_{\tau}\in\{-a,b\} \). Started from \( 0 \), necessarily \( \tau\ge\min\{a,b\} \). If one starts already at a barrier, \( \tau=0 \).
\end{definition}

\begin{lemma}[\( \tau<\infty \) almost surely]
\label{lem:tau-finite}
The exit time of \cref{def:two-barrier-exit} is almost surely finite, and \( \E[\tau]<\infty \). In particular \( s^{\tau} \) is well-defined and bounded by \( 1 \) for \( s\in(0,1) \).
\end{lemma}

\begin{proof}
The interval \( \{-a+1,\dots,b-1\} \) is finite. From any interior point, a block of \( a+b \) consecutive \( +1 \) (or of \( a+b \) consecutive \( -1 \)) forces an exit, and has probability \( 2^{-(a+b)}>0 \). Thus the number of such blocks before exit is stochastically dominated by a geometric random variable, so \( \tau \) has an exponential tail.
\end{proof}

\begin{proposition}[Exponential martingale of the walk]
\label{prop:exp-mg-srw}
For each \( \lambda\in\mathbb{R} \),
\[
M_n^{\lambda}
\coloneqq
\frac{e^{\lambda S_n}}{(\cosh\lambda)^{n}}
=
e^{\lambda S_n}\,s_{\lambda}^{n},
\qquad
s_{\lambda}
\coloneqq
\frac{1}{\cosh\lambda},
\]
is a martingale with respect to \( (\mathcal{F}_n) \). This is the output of part~(1), rewritten as a multiplicative increment of mean \( 1 \).
\end{proposition}

\begin{proof}
Independence of the steps and part~(1) give \( \E[e^{\lambda\xi_{n+1}}\mid\mathcal{F}_n]=\cosh\lambda \), hence \( \E[M_{n+1}^{\lambda}\mid\mathcal{F}_n]=M_n^{\lambda} \).
\end{proof}

\begin{theorem}[Optional stopping theorem (OST), discrete case]
\label{thm:ost-discrete}
Let \( (M_n)_{n\ge 0} \) be a martingale and \( \tau \) a stopping time. Then \( \E[M_{\tau}]=\E[M_0] \) in each of the following situations \citep[Theorems~5.7.3--5.7.6]{durrett2010}:
\begin{enumerate}[label=\textup{(\roman*)}]
    \item
    \( \tau \) is bounded;
    \item
    \( \tau<\infty \) a.s.\ and \( \{M_{\tau\wedge n}:n\ge 0\} \) is uniformly integrable (in particular if \( |M_{\tau\wedge n}|\le K \) for a deterministic \( K \));
    \item
    \( \E[\tau]<\infty \) and the increments of \( M \) are uniformly bounded.
\end{enumerate}
The standing method is to apply (i) to \( \tau\wedge n \) and pass to the limit by (ii). Item~(iii) is the usual justification of Wald identities \( \E[S_{\tau}]=0 \) and \( \E[\tau]=\E[S_{\tau}^{2}] \). It is \emph{not} automatic: Problem~4 of this exam is the standard counter-example.
\end{theorem}

\begin{remark}[Why \( \tau\wedge n \)]
\label{rem:tau-wedge-n}
The symbol is the lattice meet, not a hat and not an estimator:
\[
x\wedge y
\coloneqq
\min(x,y),
\qquad
x\vee y
\coloneqq
\max(x,y).
\]
Thus \( \tau\wedge n \) is the earlier of ``the walk has already exited'' and ``the clock has reached the deterministic time \( n \)''. It is itself a stopping time, and it is bounded: \( \tau\wedge n\le n \). Pathwise,
\[
S_{\tau\wedge n}
=
\begin{cases}
S_{\tau}\in\{-a,b\},
&
\text{if }\tau\le n,\\
S_n\in\{-a+1,\dots,b-1\},
&
\text{if }\tau>n.
\end{cases}
\]
The same splitting holds for any process evaluated at \( \tau\wedge n \).

Why this cutoff appears in every optional-stopping computation. The theorem is not ``\( M \) martingale and \( \tau \) a stopping time \( \Rightarrow\E[M_{\tau}]=\E[M_0] \)''. Item~(i) requires a \emph{bounded} stopping time. The hitting time \( \tau \) of the present problem is finite a.s.\ (\cref{lem:tau-finite}) but not bounded: for every \( N \) one has \( \Pbb(\tau>N)>0 \). So one cannot quote (i) at \( \tau \) itself. The workaround is mechanical:

\begin{enumerate}[label=\textup{(\alph*)}]
    \item
    stop at the bounded time \( \tau\wedge n \), which (i) always allows, and obtain \( \E[M_{\tau\wedge n}]=\E[M_0] \);
    \item
    send \( n\to\infty \). On \( \{\tau<\infty\} \) one has \( \tau\wedge n\to\tau \) and \( M_{\tau\wedge n}\to M_{\tau} \);
    \item
    pass the expectation through the limit. This is the only step that can fail: it needs uniform integrability of \( \{M_{\tau\wedge n}:n\ge 0\} \), or a deterministic bound \( |M_{\tau\wedge n}|\le K \), or dominated convergence.
\end{enumerate}

In this problem the walk cannot leave \( [-a,b] \) before \( \tau \), so \( |S_{\tau\wedge n}|\le\max(a,b) \) and \( \{M_{\tau\wedge n}^{\theta}\} \) is bounded. That is why (b)--(c) are free here, and why one may write \( \E[M_{\tau}^{\theta}]=1 \) after a one-line appeal to (ii).

The drills after the solution reuse the same cutoff, with different answers to (c):

\begin{itemize}
    \item
    \cref{ex:stop-s2n} (and Problem~5): \( S_{\tau\wedge n}^{2} \) or \( B_{T_a\wedge t}^{2} \) stays bounded by the walls, so \( \E[\tau\wedge n]=\E[S_{\tau\wedge n}^{2}] \) upgrades to \( \E[\tau]=\E[S_{\tau}^{2}] \);
    \item
    \cref{ex:stop-asymmetric,ex:stop-bm-laplace}: one barrier, so position is not two-sided-bounded; one uses a one-sided domination (\( M(t)\ge 1 \) and \( S\le k \), or \( u>0 \) and \( B\le a \)) and then dominated convergence;
    \item
    \cref{ex:stop-ost-fails} / Problem~4: \( X_{T\wedge n}=1 \) on \( \{T\le n\} \) and \( X_{T\wedge n}=1-2^{n} \) on \( \{T>n\} \). Then \( \E[X_{T\wedge n}]=0 \) for every \( n \), but \( X_{T\wedge n}\to 1 \) a.s.\ and \( \{X_{T\wedge n}\} \) is not UI. The \( L^{1} \) limit and the a.s.\ limit disagree, which is exactly (c) failing.
\end{itemize}

Continuous time is identical, with \( T\wedge t \) in place of \( \tau\wedge n \). First-step recurrences never introduce this cutoff, because they never invoke optional stopping.
\end{remark}

\subsubsection{Manuscript proof idea}

The calculation below follows the handwritten notes. \mserr{Red} marks a false claim or a missing justification. Part~(1) is finished and correct. Part~(2) sets up the event \( \{\tau=n\} \) correctly and then stops; the last-step observation is the right local picture, but summing \( s^{n}\Pbb(\tau=n) \) by enumerating paths is not the exam computation.

\textbf{(1).}
The notes write
\[
\E\bigl[e^{\lambda S_n}\bigr]
=
\prod_{i=1}^{n}\E\bigl[e^{\lambda\xi_i}\bigr]
=
\prod_{i=1}^{n}\tfrac12\bigl(e^{\lambda}+e^{-\lambda}\bigr)
=
(\cosh\lambda)^{n}.
\]
This is correct: independence turns the exponential of a sum into a product, and each factor is \( \cosh\lambda \). That identity is \cref{prop:exp-mg-srw}.

\textbf{(2).}
The notes write \( \E[s^{\tau}]=\sum_{n\ge 0}s^{n}\Pbb(\tau=n) \), record \( \Pbb(\tau=n)=0 \) for \( n<\min\{a,b\} \), and identify
\[
\{\tau=n\}
=
\{S_k\in(-a,b)\text{ for all }k<n\}
\cap
\{S_n\in\{-a,b\}\}.
\]
All of this is correct. The last line, ``consider the \( (n-1) \)-st step \( \Rightarrow -a+1 \) or \( b-1 \)'', is also correct: the only legal positions one step before exit are the two adjacent sites. \mserr{Gap: this describes the last step of a path of length \( n \), which is the input to a reflection-principle count of \( \Pbb(\tau=n) \). With two barriers the method of images is an infinite sum and is not what the exam wants.} The same adjacency, used from a \emph{variable} starting point rather than from a fixed time \( n \), is the first-step equation of the correct proof.

\errpanel{%
\textbf{\mserr{The error, in one box.}}

Part~(1) is done. For part~(2) do not try to sum \( \Pbb(\tau=n) \). Introduce \( u(x)=\E_x[s^{\tau}] \) and write the first-step relation \( u(x)=\frac{s}{2}(u(x+1)+u(x-1)) \) on the interior, with \( u(-a)=u(b)=1 \). The value wanted is \( u(0) \). Equivalently, stop the exponential martingale of part~(1) at \( \tau \).
}

\subsubsection{Solution of Problem 3}

\paragraph{Answer.}
\( \E[e^{\lambda S_n}]=(\cosh\lambda)^{n} \). For \( s\in(0,1) \), writing \( \mu=(1+\sqrt{1-s^{2}})/s>1 \),
\[
\E[s^{\tau}]
=
\frac{\mu^{a}+\mu^{b}}{1+\mu^{a+b}}.
\]
Equivalently, if \( \theta>0 \) solves \( \cosh\theta=1/s \), then
\[
\E[s^{\tau}]
=
\frac{\cosh\bigl(\tfrac{a-b}{2}\,\theta\bigr)}{\cosh\bigl(\tfrac{a+b}{2}\,\theta\bigr)}.
\]

\begin{proof}[Correct proof]
\textbf{(1).}
Independence:
\[
\E\bigl[e^{\lambda S_n}\bigr]
=
\prod_{i=1}^{n}\E\bigl[e^{\lambda\xi_i}\bigr]
=
(\cosh\lambda)^{n}.
\]

\textbf{(2), what is being computed.}
The object \( \E[s^{\tau}] \) is the probability generating function of the first time the walk hits \( -a \) or \( b \). Directly summing \( \sum s^{n}\Pbb(\tau=n) \) means counting two-barrier paths of every length; that is the method of images, and with two walls it becomes an infinite series. The usable move is to enlarge the problem: compute the same generating function from \emph{every} starting point in \( \{-a,\dots,b\} \), then specialise to the exam start \( 0 \).

\textbf{(2), enlarge.}
For each lattice point \( x\in\{-a,\dots,b\} \), write \( \E_x \) for the law of the walk started at \( S_0=x \), and set
\[
u(x)
\coloneqq
\E_x[s^{\tau}].
\]
The exam wants \( u(0) \). If the walk is already at a wall, one stops immediately: \( \tau=0 \) and \( s^{0}=1 \), so
\[
u(-a)
=
u(b)
=
1.
\]

\textbf{(2), one step.}
Now take an interior point \( x\in\{-a+1,\dots,b-1\} \). The first step is forced: \( \xi_1=\pm 1 \) with probability \( 1/2 \) each, the clock advances by \( 1 \), and the remaining time to hit a wall, started from \( x\pm 1 \), has the same law as a fresh copy of \( \tau \). In symbols \( \tau=1+\tau' \), where \( \tau' \) is the exit time of the shifted walk. Therefore
\[
s^{\tau}
=
s\cdot s^{\tau'},
\]
and taking \( \E_x \) gives the first-step equation
\[
u(x)
=
\frac{s}{2}\,u(x+1)
+
\frac{s}{2}\,u(x-1).
\]
The factor \( s \) is exactly the ``one step has already been used''. The two copies of \( u \) are the two places one can be after that step.

The smallest case makes the same bookkeeping visible. If \( a=b=1 \), the only interior point is \( 0 \), the two neighbours are already the walls, and the equation collapses to \( u(0)=s\cdot 1 \). Indeed \( \tau=1 \) identically, so \( \E[s^{\tau}]=s \).

\textbf{(2), solve the recurrence.}
Rearrange the interior equation as
\[
s\,u(x+1)
-
2u(x)
+
s\,u(x-1)
=
0.
\]
This is a linear homogeneous recurrence with constant coefficients. Try \( u(x)=r^{x} \):
\[
s r^{x+1}+s r^{x-1}
=
2r^{x}
\qquad\Longleftrightarrow\qquad
s r^{2}-2r+s
=
0.
\]
The discriminant is \( 4-4s^{2}=4(1-s^{2})>0 \). The two roots are
\[
\mu
\coloneqq
\frac{1+\sqrt{1-s^{2}}}{s}
>
1
\qquad\text{and}\qquad
\mu^{-1}
=
\frac{1-\sqrt{1-s^{2}}}{s}
\in
(0,1),
\]
the second being the reciprocal because the product of the roots is \( 1 \). The general interior solution is therefore
\[
u(x)
=
A\,\mu^{x}
+
B\,\mu^{-x}.
\]
(The same function automatically extends to the two walls, and we now fit it there.)

\textbf{(2), fit the two walls.}
The conditions \( u(-a)=1 \) and \( u(b)=1 \) read
\[
A\mu^{-a}+B\mu^{a}
=
1,
\qquad
A\mu^{b}+B\mu^{-b}
=
1.
\]
The two walls are being asked to take the same value, so try the symmetric ansatz \( A=c\mu^{a} \), \( B=c\mu^{b} \). Either boundary becomes
\[
c+c\mu^{a+b}
=
1
\qquad\Longrightarrow\qquad
c
=
\frac{1}{1+\mu^{a+b}}.
\]
The other boundary gives the same identity. Hence
\[
u(x)
=
\frac{\mu^{a+x}+\mu^{b-x}}{1+\mu^{a+b}}.
\]
The exam starts at \( 0 \):
\[
\E[s^{\tau}]
=
u(0)
=
\frac{\mu^{a}+\mu^{b}}{1+\mu^{a+b}}.
\]
If one prefers hyperbolic functions, set \( \mu=e^{\theta} \) with \( \theta>0 \) and \( \cosh\theta=1/s \) (this is \( \mu+\mu^{-1}=2/s \)). Factor \( e^{(a+b)\theta/2} \) from numerator and denominator to get
\[
u(0)
=
\frac{\cosh\bigl(\tfrac{a-b}{2}\,\theta\bigr)}{\cosh\bigl(\tfrac{a+b}{2}\,\theta\bigr)}.
\]

\textbf{(2), the same computation via part~(1).}
The first-step route never used part~(1). Optional stopping is the same linear algebra, read off the exponential martingale instead of off a recurrence.

Why this object at all. The target is \( \E[s^{\tau}] \). Directly summing \( \sum_n s^{n}\Pbb(\tau=n) \) is the wrong computation: two-barrier hitting probabilities are ugly. The usable fact is a mismatch of information at the exit,
\[
S_{\tau}\in\{-a,b\}
\quad\text{is two-valued},\qquad
\tau
\quad\text{is not}.
\]
Part~(1) says that \( e^{\theta S_n} \) grows by exactly \( \cosh\theta \) per step in expectation. Choose the tilt \( \theta \) so that this growth is \( 1/s \). Then the discounted process \( s^{n}e^{\theta S_n} \) is a martingale, and optional stopping produces
\[
\E\bigl[s^{\tau}e^{\theta S_{\tau}}\bigr]
=
1.
\]
The factor \( e^{\theta S_{\tau}} \) is no longer a function of a complicated path: it is \( \mu^{b} \) or \( \mu^{-a} \) according to which wall was hit. That is one linear relation between the two partial generating functions \( \alpha=\E[s^{\tau}\ind_{\{S_{\tau}=b\}}] \) and \( \beta=\E[s^{\tau}\ind_{\{S_{\tau}=-a\}}] \). The opposite tilt \( -\theta \) (same \( \cosh \), opposite wall weights) gives the second relation. Their sum is \( \E[s^{\tau}] \).

By \cref{prop:exp-mg-srw}, \( M_n^{\theta}=e^{\theta S_n}/(\cosh\theta)^{n} \) is a martingale. Choose the particular \( \theta>0 \) with \( \cosh\theta=1/s \), i.e.\ the same \( \mu=e^{\theta} \) as above. Then
\[
M_n^{\theta}
=
s^{n}\,e^{\theta S_n}.
\]
Before exit the walk is trapped in \( \{-a+1,\dots,b-1\} \), so \( \{M_{\tau\wedge n}^{\theta}\} \) is bounded (\cref{rem:tau-wedge-n}). \Cref{thm:ost-discrete} and \cref{lem:tau-finite} give \( \E[M_{\tau}^{\theta}]=M_0^{\theta}=1 \), that is
\[
\E\bigl[s^{\tau}e^{\theta S_{\tau}}\bigr]
=
1.
\]
The point of the two barriers: \( S_{\tau} \) takes only the two values \( b \) and \( -a \). Split the expectation according to which wall is hit,
\[
\alpha
\coloneqq
\E\bigl[s^{\tau}\ind_{\{S_{\tau}=b\}}\bigr],
\qquad
\beta
\coloneqq
\E\bigl[s^{\tau}\ind_{\{S_{\tau}=-a\}}\bigr],
\]
and the identity becomes \( \alpha\mu^{b}+\beta\mu^{-a}=1 \). The same argument with \( \theta\mapsto-\theta \) (the walk is symmetric) gives \( \alpha\mu^{-b}+\beta\mu^{a}=1 \). These are the same two equations as the boundary conditions above, with \( (\alpha,\beta) \) in place of \( (A,B) \). In particular \( \alpha+\beta=\E[s^{\tau}] \), and the same ansatz \( \alpha=c\mu^{a} \), \( \beta=c\mu^{b} \) produces
\[
\E[s^{\tau}]
=
\alpha+\beta
=
\frac{\mu^{a}+\mu^{b}}{1+\mu^{a+b}}.
\]
So part~(1) is not a separate computation: it manufactures the martingale whose value at \( \tau \) is \( s^{\tau} \) times a function of the exit wall, and optional stopping turns that into the boundary system for \( u \).
\end{proof}

\begin{remark}[Sanity checks]
\label{rem:p3-checks}
If \( a=b=1 \) then \( \tau=1 \) identically, so \( \E[s^{\tau}]=s \). Here \( \mu+1/\mu=2/s \) and \( 2\mu/(1+\mu^{2})=s \). As \( s\uparrow 1 \), one has \( \mu\downarrow 1 \) and the formula tends to \( 1 \), which matches \( \tau<\infty \) a.s. The same \( s\uparrow 1 \) limit of \( (1-s)^{-1}(1-\E[s^{\tau}]) \) recovers \( \E[\tau]=ab \), the classical gambler's-ruin mean (also \( \E[S_{\tau}^{2}]=\E[\tau] \) by optional stopping of \( S_n^{2}-n \)).
\end{remark}

\begin{remark}[The same stopping-time machine fails in Problem~4]
\label{rem:p3-vs-p4-ost}
Problem~4 uses the same steps \( \xi_i \) and another a.s.\ finite stopping time \( T \). Optional stopping still applies to the bounded cutoff \( T\wedge n \), and gives \( \E[X_{T\wedge n}]=0 \). It does \emph{not} give \( \E[X_T]=0 \): the family \( \{X_{T\wedge n}\} \) is not uniformly integrable, because the stake \( 2^{n} \) explodes on \( \{T>n\} \). The two-barrier trap of the present problem is exactly what prevents that explosion. The full contrast is \cref{rem:p4-ost-fails}.
\end{remark}

\subsubsection{Stopping-time examples}

The exam never asks for the definition of a stopping time by itself. It asks for a number attached to one: a PGF, a Laplace transform, a mean, a hitting probability, or a check that optional stopping fails. The machine is almost always one of four.

\begin{center}
\renewcommand{\arraystretch}{1.2}
\small
\begin{tabular}{@{}p{0.28\textwidth}p{0.34\textwidth}p{0.28\textwidth}@{}}
\toprule
Want & Do this & Do not \\
\midrule
\( \E[s^{\tau}] \) or \( \E[e^{-\lambda\tau}] \)
  & write the exp.\ martingale; OST if bounded; else first-step
  & enumerate two-barrier paths \\
\( \E[\tau] \)
  & stop \( S_n^{2}-n \) or \( B_t^{2}-t \); or differentiate the PGF
  & compute \( \Pbb(\tau=n) \) first \\
\( \Pbb(\tau<\infty) \)
  & send \( s\uparrow 1 \), or first-step for ruin
  & assume recurrence \\
\( \E[M_{\tau}]\stackrel{?}{=}\E[M_0] \)
  & check UI / bounded increments
  & quote OST blindly \\
\bottomrule
\end{tabular}
\end{center}

Decision, in order: (i)~write the exponential martingale from the increment MGF; (ii)~ask whether the hitting time keeps the stopped process bounded (two compact walls: yes; doubling stakes or a one-sided walk drifting away: no); (iii)~if yes, optional stopping; if no, first-step for the ruin probability, or exhibit the failure. Finite chains with no increment structure skip (i) and go straight to a linear system. One Brownian barrier with a path event (``not hit \( 0 \) by time \( t \)'') is reflection, not optional stopping.

The nine drills below are the QE corpus items, rewritten as self-contained exercises. Problem~3 itself is the two-barrier model case; it is not repeated. Whenever a solution writes \( \tau\wedge n \) or \( T\wedge t \), it is executing \cref{rem:tau-wedge-n}: stop at a bounded time first, then pass to the limit.

\begin{example}[Two-barrier mean via \( S_n^{2}-n \)]
\label{ex:stop-s2n}
In the setting of Problem~3, compute \( \E[\tau] \) without generating functions. (Same computation as the \( s\uparrow 1 \) limit in \cref{rem:p3-checks}.)
\end{example}

\begin{proof}[Solution]
Why this martingale. The linear martingale \( S \) only prices the exit side: OST gives \( \E[S_{\tau}]=0 \), hence the ruin probabilities, but not the clock. The identity \( \E[S_n^{2}]=n \) (i.e.\ \( \mathrm{Var}(S_n)=n \)) is the unconditional statement that time is paid for by quadratic variation. The process \( M_n=S_n^{2}-n \) is that identity made pathwise, so stopping it converts \( \E[\tau] \) into \( \E[S_{\tau}^{2}] \).

The filtration is the same as in Problem~3: \( \mathcal{F}_n=\sigma(\xi_1,\dots,\xi_n) \), with \( \mathcal{F}_0 \) trivial. Then \( S_n \) and \( M_n \) are \( \mathcal{F}_n \)-measurable, and \( |M_n|\le n+n^{2} \) so \( M_n\in L^{1} \). Write \( S_{n+1}=S_n+\xi_{n+1} \). Expanding,
\[
M_{n+1}
=
S_{n+1}^{2}-(n+1)
=
M_n+2S_n\xi_{n+1}+(\xi_{n+1}^{2}-1).
\]
The new step \( \xi_{n+1} \) is independent of \( \mathcal{F}_n \), with \( \E[\xi_{n+1}]=0 \) and \( \xi_{n+1}^{2}=1 \) a.s., so
\[
\E[S_n\xi_{n+1}\mid\mathcal{F}_n]
=
S_n\E[\xi_{n+1}]
=
0,
\qquad
\E[\xi_{n+1}^{2}-1\mid\mathcal{F}_n]
=
0.
\]
Thus \( \E[M_{n+1}\mid\mathcal{F}_n]=M_n \). (If the increments were only mean-zero with variance \( \sigma^{2} \), the same computation would produce the martingale \( S_n^{2}-\sigma^{2}n \).)

Stop at \( \tau\wedge n \): the walk has not left \( [-a,b] \), so \( |S_{\tau\wedge n}|\le\max(a,b) \) and the optional stopping theorem (\cref{thm:ost-discrete}; OST below) applies,
\[
\E[S_{\tau\wedge n}^{2}]
=
\E[\tau\wedge n].
\]
Let \( n\to\infty \). Bounded convergence on the left and monotone convergence on the right give \( \E[S_{\tau}^{2}]=\E[\tau] \).

The exit probabilities are a separate, easier optional stopping. The walk \( S \) itself is a martingale (\( \E[\xi_{n+1}\mid\mathcal{F}_n]=0 \)). Up to \( \tau \) it is bounded by \( \max(a,b) \), so OST gives \( \E[S_{\tau}]=S_0=0 \). But \( S_{\tau} \) takes only two values: write \( p\coloneqq\Pbb(S_{\tau}=b) \). Then
\[
\E[S_{\tau}]
=
p\,b+(1-p)(-a)
=
0
\qquad\Longrightarrow\qquad
p=\frac{a}{a+b},
\qquad
1-p=\frac{b}{a+b}.
\]
(The larger the lower wall \( a \), the more likely to hit \( +b \) first: you start closer to the top.) The second moment is then elementary,
\[
\E[S_{\tau}^{2}]
=
p\,b^{2}+(1-p)\,a^{2}
=
\frac{a}{a+b}\,b^{2}+\frac{b}{a+b}\,a^{2}
=
\frac{ab(a+b)}{a+b}
=
ab.
\]
Thus \( \E[\tau]=\E[S_{\tau}^{2}]=ab \). Technique: when the increment is mean-zero and variance-one, the quadratic martingale converts a hitting time into an exit-position second moment. The continuous analogue is \( B_t^{2}-t \) at the two-sided exit of \( (-a,a) \), which yields \( \E[T_a]=a^{2} \) (Problem~5).
\end{proof}

\begin{example}[One-barrier PGF, unbounded domain]
\label{ex:stop-one-barrier-pgf}
Let \( (S_n) \) be simple symmetric random walk on \( \mathbb{Z} \), \( a\in\mathbb{N} \), and \( T_a=\inf\{n\ge 0:S_n=a\} \). For \( 0<s<1 \), compute \( \E[s^{T_a}] \).
\end{example}

\begin{proof}[Solution]
Set \( \varphi(x)=\E_x[s^{T_a}] \) for \( x\le a \). Then \( \varphi(a)=1 \), and the first-step equation is the same recurrence as in Problem~3,
\[
\varphi(x)
=
\frac{s}{2}\bigl(\varphi(x+1)+\varphi(x-1)\bigr),
\qquad
x<a.
\]
This is a linear homogeneous recurrence with constant coefficients. Rearrange as
\[
\varphi(x+1)-\frac{2}{s}\,\varphi(x)+\varphi(x-1)=0
\]
and try \( \varphi(x)=r^{x} \). The characteristic equation is
\[
r+\frac{1}{r}=\frac{2}{s}
\qquad\Longleftrightarrow\qquad
r^{2}-\frac{2}{s}\,r+1=0.
\]
The product of the roots is \( 1 \), so they are reciprocals. The discriminant is \( 4(1-s^{2})/s^{2} \), hence
\[
r
=
\frac{1\pm\sqrt{1-s^{2}}}{s}.
\]
Write \( \mu\coloneqq(1+\sqrt{1-s^{2}})/s \). Then \( \mu>1 \) (because \( 1+\sqrt{1-s^{2}}>s \)) and the other root is
\[
\mu^{-1}
=
\frac{1-\sqrt{1-s^{2}}}{s}
\in
(0,1).
\]
The general solution on \( \{\dots,a-2,a-1,a\} \) is therefore \( \varphi(x)=A\mu^{x}+B\mu^{-x} \).

The generating function is bounded: \( 0\le\varphi(x)\le 1 \). As \( x\to-\infty \), the term \( \mu^{x}\to 0 \), while \( \mu^{-x}=\mu^{|x|}\to\infty \). The only way to stay bounded is \( B=0 \). The remaining condition is the wall: \( \varphi(a)=1 \) forces \( A\mu^{a}=1 \), so \( A=\mu^{-a} \) and \( \varphi(x)=\mu^{x-a} \). Starting at \( 0 \),
\[
\E[s^{T_a}]
=
\varphi(0)
=
\mu^{-a}
=
\bigl(\mu^{-1}\bigr)^{a}
=
\Bigl(\frac{1-\sqrt{1-s^{2}}}{s}\Bigr)^{a}.
\]
The last line is not a new formula: it is the smaller characteristic root, raised to the distance from the wall. For \( a=1 \) this is already the generating function of the first return / first hit of \( +1 \) from \( 0 \); the spatial Markov property multiplies one more copy of \( \mu^{-1} \) for each extra unit of distance.

Technique: one compact wall plus an unbounded side --- keep the root that stays bounded at infinity. This is Problem~3 with the lower barrier sent to \( -\infty \). Path-counting (ballot / reflection) also works here, because there is only one wall; it becomes unusable the moment a second wall is present.
\end{proof}

\begin{example}[OST fails: doubling martingale]
\label{ex:stop-ost-fails}
Same \( \xi_i \) as Problem~3. Let \( X_n=\sum_{j=1}^{n}2^{j-1}\xi_j \) and \( T=\min\{j\ge 1:\xi_j=+1\} \). Show that \( (X_n) \) is a martingale, identify the law of \( T \), and compute \( \E[X_T] \). Why does optional stopping fail? (2025~Spring~P4, next on this paper.)
\end{example}

\begin{proof}[Solution]
Independence and \( \E[\xi_{n+1}]=0 \) give \( \E[X_{n+1}\mid\mathcal{F}_n]=X_n \). The waiting time is Geometric\( (1/2) \) on \( \{1,2,\dots\} \), so \( \E[T]=2 \). On \( \{T=k\} \) one has \( \xi_1=\cdots=\xi_{k-1}=-1 \) and \( \xi_k=+1 \), hence
\[
X_T
=
-\sum_{j=1}^{k-1}2^{j-1}+2^{k-1}
=
1
\]
almost surely, and \( \E[X_T]=1\neq 0=\E[X_0] \). The increments grow like \( 2^{n} \): \( \{X_{T\wedge n}\} \) is not uniformly integrable. Technique: OST is a theorem with hypotheses, not an identity \( \E[M_{\tau}]=\E[M_0] \). Contrast with Problem~3, where \( |S_{\tau\wedge n}|\le\max(a,b) \) keeps the exponential martingale bounded, so OST applies.
\end{proof}

\begin{example}[Asymmetric one-barrier: OST, then \( s\uparrow 1 \)]
\label{ex:stop-asymmetric}
Let \( \Pbb(X_i=+1)=p=1-\Pbb(X_i=-1) \), \( S_n=X_1+\cdots+X_n \), \( T_k=\inf\{n:S_n=k\} \) for \( k\in\mathbb{N} \), and \( M(t)=\E[e^{tX_1}]=pe^{t}+(1-p)e^{-t} \). Show that \( e^{tS_n}M(t)^{-n} \) is a martingale. For \( p\ge 1/2 \) and \( t>0 \), deduce \( e^{tk}\E[M(t)^{-T_k}]=1 \). For \( p<1/2 \), compute \( \Pbb(T_k<\infty) \). (2024~Spring~P4.)
\end{example}

\begin{proof}[Solution]
The increment MGF is \( M(t) \), so the exponential martingale is the same construction as Problem~3(1). If \( p\ge 1/2 \) then \( T_k<\infty \) a.s.\ (nonnegative drift). For \( t>0 \) one has \( M(t)\ge 1 \), and on \( \{n\le T_k\} \) the walk has not reached \( k \), so \( e^{tS_n}M(t)^{-n}\le e^{tk} \). Stop at \( T_k\wedge n \) and pass to the limit by dominated convergence:
\[
e^{tk}\E\bigl[M(t)^{-T_k}\bigr]
=
1.
\]
If \( p<1/2 \) the walk is transient to \( -\infty \). The hitting probabilities \( u(x)=\Pbb_x(T_k<\infty) \) solve \( u(x)=p\,u(x+1)+(1-p)\,u(x-1) \) with \( u(k)=1 \) and \( u(-\infty)=0 \), hence \( u(0)=(p/(1-p))^{k} \). Technique: the same exponential martingale as part~(1); the new ingredient is a domination that uses \( M(t)\ge 1 \) and a one-sided barrier. When domination fails (negative drift), switch to the first-step equation for \( \Pbb(\tau<\infty) \). Sending the PGF parameter to \( 1 \) is the discrete analogue of that switch.
\end{proof}

\begin{example}[Brownian hitting-time Laplace transform]
\label{ex:stop-bm-laplace}
Let \( B \) be standard Brownian motion, \( a>0 \), and \( T_a=\inf\{t>0:B_t\ge a\} \). Compute \( \E[e^{-\lambda T_a}] \) for \( \lambda>0 \). (2024~Spring~P6.)
\end{example}

\begin{proof}[Solution]
The process \( e^{uB_t-u^{2}t/2} \) is a martingale for every \( u\in\mathbb{R} \). Stop at \( T_a\wedge t \): on this event \( B\le a \), so if \( u>0 \) the exponential is bounded by \( e^{ua} \). Let \( t\to\infty \). Continuity of paths gives \( B_{T_a}=a \) on \( \{T_a<\infty\} \), and Fatou plus the bound produce
\[
\E\bigl[e^{ua-u^{2}T_a/2}\mathbf{1}_{\{T_a<\infty\}}\bigr]
=
1.
\]
Choose \( u=\sqrt{2\lambda} \). Then \( \E[e^{-\lambda T_a}]=e^{-a\sqrt{2\lambda}} \). Sending \( \lambda\downarrow 0 \) yields \( T_a<\infty \) a.s. Technique: continuous-time copy of Problem~3(1)--(2). The free parameter \( u \) is chosen so that the spatial factor becomes a constant and the time factor becomes the Laplace kernel. For the two-sided exit \( \inf\{t:|B_t|=a\} \) (Problem~5 on this paper), the same martingale works, but symmetry of \( \pm B \) is shorter for independence of time and exit side.
\end{proof}

\begin{example}[Write the exponential martingale, then stop]
\label{ex:stop-write-mgf}
Let \( (X_k) \) be i.i.d.\ standard Laplace (density \( \tfrac12 e^{-|x|} \)). Find the deterministic \( A_n(\theta) \) making \( \exp(\theta\sum_{k=1}^{n}X_k-A_n(\theta)) \) a martingale, for \( |\theta|<1 \). (2023~Fall~P2.)
\end{example}

\begin{proof}[Solution]
Independence forces \( A_n(\theta)=n\log\E[e^{\theta X_1}] \). The Laplace MGF is \( \E[e^{\theta X_1}]=1/(1-\theta^{2}) \) for \( |\theta|<1 \), so \( A_n(\theta)=n\log\bigl(1/(1-\theta^{2})\bigr) \). That is Problem~3(1) with \( \cosh\lambda \) replaced by the increment MGF. Technique: this is always the first line of a hitting-time Laplace-transform problem. After \( A_n \) is known, every subsequent computation is optional stopping (or first-step, which solves the same linear system).
\end{proof}

\begin{example}[Extinction time of a branching process]
\label{ex:stop-branching}
Let \( (Z_n) \) be Galton--Watson with offspring PGF \( f \), \( Z_0=1 \), and \( \tau=\inf\{n:Z_n=0\} \). Express \( \Pbb(\tau\le n) \) and \( \Pbb(\tau<\infty) \). (2026~Spring~P5.)
\end{example}

\begin{proof}[Solution]
The extinction time is a stopping time of \( (Z_n) \). One generation of offspring, then \( Z_1 \) independent copies: \( \Pbb(\tau\le n)=f^{\circ n}(0) \), and \( \Pbb(\tau<\infty) \) is the smallest root in \( [0,1] \) of \( f(s)=s \). Technique: first-step, but the state is the population size, not a spatial position. There is no useful exponential martingale on \( \mathbb{Z} \) here; the generating function \emph{is} the computation.
\end{proof}

\begin{example}[Mean hitting time on a finite chain]
\label{ex:stop-finite-chain}
Three fair dice; at each step one may reroll any subset. Let \( T \) be the waiting time until the face \( 666 \) is held. Under the optimal policy (keep every current six, reroll the rest), write the system for \( e_k=\E_k[T] \), where \( k \) is the number of sixes already held. (2023~Fall~P4.)
\end{example}

\begin{proof}[Solution]
The state space is \( \{0,1,2,3\} \), absorbing at \( 3 \). First-step:
\[
e_3=0,
\qquad
e_k
=
1+\sum_{j=0}^{3}p_{kj}e_j,
\qquad
k=0,1,2,
\]
where \( p_{kj} \) is the binomial probability of obtaining \( j-k \) new sixes among the \( 3-k \) dice rerolled (and \( j<k \) is impossible). The chain is finite, so the linear system is the whole computation: no exponential martingale, no uniform-integrability check. Technique: same first-step machine as \( u(x) \) in Problem~3, with the unknown changed from a generating function to a mean. Use this whenever the state space is already finite and the increment has no MGF structure worth exploiting.
\end{proof}

\begin{example}[Reflection, not optional stopping]
\label{ex:stop-reflection}
Let \( X \) be Brownian motion started at \( x>0 \), and \( \tau_0=\inf\{t:X_t=0\} \). Identify the joint law of \( (X_t,\{\tau_0>t\}) \). (2024~Fall~P6(a).)
\end{example}

\begin{proof}[Solution]
This is a path event at a fixed time, not a transform of a hitting time. The reflection principle gives the transition density of Brownian motion killed at \( 0 \):
\[
\Pbb_x(X_t\in dy,\,\tau_0>t)
=
\bigl(g_t(y-x)-g_t(y+x)\bigr)\,dy,
\qquad
y>0,
\]
where \( g_t \) is the \( \mathcal{N}(0,t) \) density. Technique: one absorbing Brownian barrier plus a question about paths up to a deterministic time --- reflection. The same path-count is what the handwritten notes attempted for Problem~3; it is the wrong default as soon as there are two barriers (or as soon as the question is a generating function rather than a path probability).
\end{proof}

\subsection{Problem 4. [10 pts.]}
Suppose \( \{\xi_i\} \) are i.i.d.\ Bernoulli random variables with \( \mathbb{P}[\xi_1 = 1] = \mathbb{P}[\xi_1 = -1] = 1/2 \). Set

\[
X_n = \sum_{j=1}^n 2^{j-1} \xi_j.
\]

\begin{enumerate}[label=(\arabic*)]
    \item Define \( T = \min\{j \ge 1 : \xi_j = +1\} \). Derive the law of \( T \) and calculate \( \mathbb{E}[T] \).
    \item Show that \( \{X_n\} \) is a martingale and calculate \( \mathbb{E}[X_T] \).
\end{enumerate}

\setcounter{section}{4}

\subsubsection{Definitions used in Problem 4}

The process in part~(2) is not the exponential martingale of Problem~3. It is the doubling martingale on the same steps \( \xi_i \): the stake at time \( n+1 \) is \( 2^{n} \), while the increment still has conditional mean zero. The waiting time \( T \) of part~(1) is a stopping time in the sense of \cref{def:stopping-time-discrete}.

\begin{definition}[Discrete martingale]
\label{def:discrete-martingale}
Let \( (\mathcal{F}_n)_{n\ge 1} \) be a filtration (an increasing sequence of \(\sigma\)-algebras). An integrable process \( (X_n)_{n\ge 1} \) is a \emph{martingale} with respect to \( (\mathcal{F}_n) \) if \( X_n \) is \(\mathcal{F}_n\)-measurable for every \( n \) and
\[
\E[X_{n+1}\mid\mathcal{F}_n]=X_n
\qquad\text{a.s.}
\]
Equivalently, \( \E[X_{n+1}-X_n\mid\mathcal{F}_n]=0 \). If the equality is replaced by \( \le \) (resp.\ \( \ge \)), the process is a supermartingale (resp.\ submartingale). See \citep[\S5.2]{durrett2010}.
\end{definition}

\begin{definition}[Doubling martingale]
\label{def:doubling-martingale}
Let \( \{\xi_i\} \) be as in the problem, and let \( \mathcal{F}_n=\sigma(\xi_1,\dots,\xi_n) \) (with \( \mathcal{F}_0 \) trivial). The \emph{doubling martingale} of Problem~4 is
\[
X_n
\coloneqq
\sum_{j=1}^{n}2^{j-1}\xi_j,
\qquad
n\ge 1,
\]
together with the convention \( X_0=0 \). Integrability is automatic: \( |X_n|\le 2^{n}-1 \). The one-step decomposition is \( X_{n+1}=X_n+2^{n}\xi_{n+1} \). Since \( \xi_{n+1} \) is independent of \( \mathcal{F}_n \) and \( \E[\xi_{n+1}]=0 \),
\[
\E[X_{n+1}\mid\mathcal{F}_n]
=
X_n+2^{n}\E[\xi_{n+1}]
=
X_n,
\]
so \( (X_n)_{n\ge 1} \) is a martingale in the sense of \cref{def:discrete-martingale}. This is the discrete St.\ Petersburg / doubling process: the stake doubles at every step, but the sign remains mean-zero.
\end{definition}

\begin{definition}[First plus-one time]
\label{def:first-plus-one}
On the same filtration,
\[
T
\coloneqq
\min\{j\ge 1:\xi_j=+1\},
\]
with the convention \( \min\emptyset=\infty \). This is a stopping time: \( \{T=n\}=\{\xi_1=\cdots=\xi_{n-1}=-1,\,\xi_n=+1\}\in\mathcal{F}_n \). On the present i.i.d.\ Rademacher steps one has \( T<\infty \) a.s.\ and \( T \) is Geometric\( (1/2) \) on \( \{1,2,\dots\} \).
\end{definition}

\subsubsection{Manuscript proof idea}

The calculation below follows the handwritten notes. The notes label the two items (a) and (b); these are the exam's (1) and (2). \mserr{Red} marks a false claim or a missing justification. Part~(1) and the evaluation of \( \E[X_T] \) are finished and correct. Part~(2) states the definition of a martingale and then replaces the conditional-expectation check by the weaker identity \( \E[X_n]=0 \).

\textbf{(1).}
The notes identify
\[
\{T=k\}
=
\{\xi_j=-1\text{ for all }j<k\}
\cap
\{\xi_k=+1\}
\]
and compute
\[
\Pbb(T=k)
=
\bigl(\tfrac12\bigr)^{k-1}\cdot\tfrac12
=
\bigl(\tfrac12\bigr)^{k}.
\]
This is the law of \( T \): Geometric\( (1/2) \) on \( \{1,2,\dots\} \). They then write \( \E[T]=\sum_{k\ge 1}k\,(1/2)^{k} \). After a crossed-out derivative they recover the standard identity
\[
\sum_{k\ge 1}k\,x^{k-1}
=
\Bigl(\frac{x}{1-x}\Bigr)'
=
\frac{1}{1-x}+\frac{x}{(1-x)^{2}}
=
\frac{1}{(1-x)^{2}},
\]
evaluate at \( x=1/2 \) to get \( 4 \), and conclude \( \E[T]=(1/2)\cdot 4=2 \). The arithmetic is correct. (Equivalently \( \sum k x^{k}=x/(1-x)^{2} \) at \( x=1/2 \) is \( 2 \), with no extra factor.)

\textbf{(2), martingale.}
The notes list the three items of \cref{def:discrete-martingale}: integrability, adaptedness, and \( \E[X_{n+1}\mid\mathcal{F}_n]=X_n \). Integrability and adaptedness are immediate from \cref{def:doubling-martingale}. They then compute only
\[
\E[X_n]
=
\sum_{j=1}^{n}2^{j-1}\E[\xi_j]
=
0.
\]
This is true, and it is what a martingale implies for the unconditional means, but \mserr{it is not the martingale property.} A crossed-out line on the notes begins the right computation \( \E[X_n+2^{n}\xi_{n+1}\mid\mathcal{F}_n] \) and then stops. The missing step is \( \E[\xi_{n+1}\mid\mathcal{F}_n]=\E[\xi_{n+1}]=0 \), hence \( \E[X_{n+1}\mid\mathcal{F}_n]=X_n \).

\textbf{(2), \( \E[X_T] \).}
A crossed-out first attempt replaces \( X_T \) by the unconditional mean:
\[
\mserr{\E[X_T]=\sum_{n}\E[X_n]\,\Pbb(T=n)=0}.
\]
That identity would require pulling \( X_n \) out of the event \( \{T=n\} \), which is illegal: on \( \{T=n\} \) the path is deterministic. The notes correctly abandon this and compute
\[
X_n\bigm|_{T=n}
=
-\sum_{j=1}^{n-1}2^{j-1}+2^{n-1}
=
-(2^{n-1}-1)+2^{n-1}
=
1,
\]
hence \( \E[X_T\mid T=n]=1 \) and \( \E[X_T]=\sum_{n\ge 1}1\cdot(1/2)^{n}=1 \). This is correct, and in fact stronger: \( X_T=1 \) almost surely, so the average is unnecessary.

\errpanel{%
\textbf{\mserr{The error, in one box.}}

The law of \( T \) and \( \E[T]=2 \) are done. \( \E[X_T]=1 \) is done.

\( \E[X_n]=0 \) does not show that \( \{X_n\} \) is a martingale. Write \( X_{n+1}=X_n+2^{n}\xi_{n+1} \) and use independence of the next sign.

Do not write \( \E[X_T]=\sum\E[X_n]\Pbb(T=n) \). On \( \{T=n\} \) one has \( X_T=1 \), not \( 0 \). Optional stopping fails: \( \E[X_T]=1\neq 0=\E[X_n] \), because \( \{X_{T\wedge n}\} \) is not uniformly integrable.
}

\subsubsection{Solution of Problem 4}

\paragraph{Answer.}
\( T \) is Geometric\( (1/2) \) on \( \{1,2,\dots\} \), and \( \E[T]=2 \). The process \( \{X_n\} \) is a martingale with respect to \( \mathcal{F}_n=\sigma(\xi_1,\dots,\xi_n) \), and \( \E[X_T]=1 \).

\begin{proof}[Correct proof]
\textbf{(1).}
Independence of the steps gives
\[
\Pbb(T=k)
=
\Pbb(\xi_1=\cdots=\xi_{k-1}=-1,\,\xi_k=+1)
=
2^{-k},
\qquad
k\ge 1.
\]
Thus \( T\sim\mathrm{Geometric}(1/2) \) on \( \{1,2,\dots\} \). In particular \( T<\infty \) a.s.\ and
\[
\E[T]
=
\sum_{k\ge 1}k\,2^{-k}
=
\frac{x}{(1-x)^{2}}\Bigm|_{x=1/2}
=
2.
\]

\textbf{(2), martingale.}
Let \( \mathcal{F}_n=\sigma(\xi_1,\dots,\xi_n) \). Then \( X_n \) is \(\mathcal{F}_n\)-measurable and \( |X_n|\le 2^{n}-1 \), so \( \E|X_n|<\infty \). The increment is \( X_{n+1}-X_n=2^{n}\xi_{n+1} \), and \( \xi_{n+1} \) is independent of \( \mathcal{F}_n \) with mean zero, hence
\[
\E[X_{n+1}\mid\mathcal{F}_n]
=
X_n+2^{n}\E[\xi_{n+1}]
=
X_n.
\]
This is \cref{def:doubling-martingale}.

\textbf{(2), \( \E[X_T] \).}
On the event \( \{T=n\} \) one has \( \xi_1=\cdots=\xi_{n-1}=-1 \) and \( \xi_n=+1 \), so
\[
X_T
=
X_n
=
-\sum_{j=1}^{n-1}2^{j-1}+2^{n-1}
=
1.
\]
Since \( T<\infty \) a.s., this holds on a set of probability one: \( X_T=1 \) a.s., and therefore \( \E[X_T]=1 \).
\end{proof}

\begin{remark}[Why OST applies in Problem~3 but not here]
\label{rem:p4-ost-fails}
The stopping time is not the obstruction. Both \( \tau \) of Problem~3 and \( T \) of Problem~4 are stopping times, both are a.s.\ finite, and both processes are martingales. The naive line
\[
\text{``\( X \) is a martingale and \( T \) is a stopping time, so \( \E[X_T]=\E[X_n]=0 \)''}
\]
is a misreading of \cref{thm:ost-discrete}: the theorem is not ``martingale + stopping time \( \Rightarrow\E[M_{\tau}]=\E[M_0] \)''. It requires one of (i)--(iii). The same checklist, applied to the two problems:

\begin{center}
\renewcommand{\arraystretch}{1.2}
\small
\begin{tabular}{@{}lll@{}}
\toprule
 & Problem~3, \( M_n^{\theta} \) or \( S_n^{2}-n \) & Problem~4, \( X_n \) \\
\midrule
stopping time a.s.\ finite & yes (\cref{lem:tau-finite}) & yes (\( T \) geometric) \\
bounded? (i) & no (\( \tau \) unbounded) & no (\( T \) unbounded) \\
increments uniformly bounded? (iii) & yes (\( \pm 1 \)) & no (stake \( 2^{n} \)) \\
\( \{M_{\tau\wedge n}\} \) bounded / UI? (ii)
  & yes: trapped in \( [-a,b] \)
  & no: see below \\
conclusion \( \E[M_{\tau}]=\E[M_0] \) & legal & illegal \\
\bottomrule
\end{tabular}
\end{center}

What one is always allowed to do, in both problems, is stop at the \emph{bounded} time \( \tau\wedge n \) or \( T\wedge n \). Item~(i) gives, with no extra hypothesis,
\[
\E\bigl[M_{\tau\wedge n}^{\theta}\bigr]
=
1,
\qquad
\E[X_{T\wedge n}]
=
0.
\]
The second identity is true for every \( n \), and is exactly the ``equals \( \E[X_n]=0 \)'' the naive argument wanted --- but for \( T\wedge n \), not for \( T \). The only additional step is \( n\to\infty \). That is \cref{rem:tau-wedge-n}(c), and it is where the two problems split.

In Problem~3 the walk cannot leave \( [-a,b] \) before \( \tau \), so \( |S_{\tau\wedge n}|\le\max(a,b) \) and \( \{M_{\tau\wedge n}^{\theta}\} \) is deterministically bounded. Bounded convergence upgrades \( \E[M_{\tau\wedge n}^{\theta}]=1 \) to \( \E[M_{\tau}^{\theta}]=1 \).

Here the stake doubles. Pathwise,
\[
X_{T\wedge n}
=
\begin{cases}
1,
&
\text{if }T\le n,\\
1-2^{n},
&
\text{if }T>n.
\end{cases}
\]
The event \( \{T>n\} \) has probability \( 2^{-n} \), but on that event the process sits at \( 1-2^{n} \). The rare path still contributes
\[
(1-2^{n})\cdot 2^{-n}
=
2^{-n}-1
\]
to the mean, which exactly cancels the \( +1 \) accumulated on \( \{T\le n\} \):
\[
\E[X_{T\wedge n}]
=
1\cdot\bigl(1-2^{-n}\bigr)+(1-2^{n})\,2^{-n}
=
0.
\]
As \( n\to\infty \) one has \( T<\infty \) a.s., so \( X_{T\wedge n}\to X_T=1 \) almost surely. The expectations stay at \( 0 \) and never catch the a.s.\ limit. That is the definition of a failure of uniform integrability: the \( L^{1} \) limit and the a.s.\ limit disagree. Hence \( \E[X_T]=1\neq 0=\E[X_0] \).

The exam writes Problem~4 immediately after Problem~3 so that this contrast is the point. Same \( \xi_i \), same filtration, a perfectly legal stopping time in both cases; the difference is whether the stopped process is prevented from exploding on the event \( \{\tau>n\} \).
\end{remark}

\subsection{Problem 5. [10 pts.]}
Let \( \{B_t\}_{t \ge 0} \) be a one-dimensional Brownian motion starting from the origin (i.e, \( B_0 = 0 \)). For \( a \in (0, \infty) \), define \( T_a := \inf\{t \ge 0 : |B_t| = a\} \).
\begin{enumerate}[label=(\arabic*)]
    \item Prove that \( T_a < \infty \) a.s.
    \item Prove that \( T_a \) and \( \mathbf{1}_{\{B_{T_a} = a\}} \) are independent.
\end{enumerate}

\setcounter{section}{5}

\subsubsection{Hint for Problem 5}

Do not compute \( \Pbb(T_a=t) \) or the density of \( T_a \). The time is continuous, so every singleton has probability zero; the density exists (reflection / theta function) but is not what the exam wants. The same warning as in Problem~3: path-counting the exit is the wrong default once there are two barriers.

The object is two-sided. Write \( \varepsilon\coloneqq\ind_{\{B_{T_a}=a\}} \). Then \( T_a \) is the first exit of \( (-a,a) \), and on \( \{T_a<\infty\} \) one has \( B_{T_a}\in\{-a,a\} \). This is not the one-sided hitting time \( \inf\{t:B_t\ge a\} \) of 2024 Spring Problem~6; that paper uses the exponential martingale \( e^{uB_t-u^{2}t/2} \). Here the two walls are symmetric about the origin, so \( \pm B \) is the shorter machine.

\textbf{(1).}
A nonnegative random variable is almost surely finite if and only if its tail vanishes: \( \Pbb(T_a>t)\to 0 \). The identity \( \E[T_a]=\int_0^{\infty}\Pbb(T_a>t)\,dt \) is true and stronger, but the exam only asks for \( T_a<\infty \) a.s. The integral route also forces a tail estimate. The crudest bound
\[
\Pbb(T_a>t)
\le
\Pbb(|B_t|<a)
\asymp
\frac{a}{\sqrt{t}}
\]
has a divergent integral, so it does \emph{not} prove \( \E[T_a]<\infty \). It does prove the tail vanishes, because
\[
\{T_a>t\}
\subset
\{|B_t|<a\}
\]
and \( \Pbb(|B_t|<a)\to 0 \). That inclusion is the short write-up of (1).

If one wants the expectation as well (optional, and the continuous analogue of \( \E[\tau]=ab \) in Problem~3), stop the martingale \( B_t^{2}-t \) at \( T_a\wedge t \). Then \( |B_{T_a\wedge t}|\le a \), hence \( \E[T_a\wedge t]=\E[B_{T_a\wedge t}^{2}]\le a^{2} \). Let \( t\to\infty \) to get \( \E[T_a]\le a^{2} \), which is more than the exam asks.

\textbf{(2).}
Independence of \( T_a \) and \( \varepsilon \) is the statement that \( \E[f(T_a)\,\varepsilon]=\E[f(T_a)]\,\E[\varepsilon] \) for every bounded Borel \( f \). The process \( -B \) is again a standard Brownian motion, and the two walls are symmetric, so
\[
T_a(B)=T_a(-B),
\qquad
\varepsilon(-B)=1-\varepsilon(B).
\]
Thus \( (T_a,\varepsilon)\stackrel{d}{=}(T_a,1-\varepsilon) \). This forces \( \Pbb(\varepsilon=1)=1/2 \), and then \( \E[f(T_a)\,\varepsilon]=\frac12\E[f(T_a)] \). Do not try to factor a joint density of \( (T_a,B_{T_a}) \).

The dictionary against the neighbouring papers:

\begin{center}
\renewcommand{\arraystretch}{1.25}
\small
\begin{tabular}{@{}lll@{}}
\toprule
& Discrete two-barrier (P3) & Continuous two-barrier (P5) \\
\midrule
Finiteness & finite state, geometric blocks & \( \{T_a>t\}\subset\{|B_t|<a\} \), or \( B_t^{2}-t \) \\
Which side & \( \E[s^{\tau}e^{\theta S_{\tau}}]=1 \) & \( \pm B \) swaps the two walls \\
Do not compute & \( \Pbb(\tau=n) \) by reflection & density of \( T_a \) \\
\bottomrule
\end{tabular}
\end{center}

\begin{definition}[Two-sided Brownian exit]
\label{def:bm-two-sided-exit}
Let \( (B_t)_{t\ge 0} \) be standard one-dimensional Brownian motion with \( B_0=0 \), and let \( a>0 \). The \emph{two-sided exit time} is
\[
T_a
\coloneqq
\inf\{t\ge 0:|B_t|=a\},
\]
with the convention \( \inf\emptyset=\infty \). This is a stopping time of the Brownian filtration: \( \{T_a\le t\}=\{\sup_{s\le t}|B_s|\ge a\} \). On \( \{T_a<\infty\} \) the exit location is \( B_{T_a}\in\{-a,a\} \). Write \( \varepsilon\coloneqq\ind_{\{B_{T_a}=a\}} \) for the upper-wall indicator. See \citep[\S8.1, \S8.2]{durrett2010}.
\end{definition}

\subsubsection{Solution of Problem 5}

\paragraph{Answer.}
\( T_a<\infty \) almost surely, because \( \{T_a>t\}\subset\{|B_t|<a\} \) and the right-hand side has probability tending to zero. Writing \( \varepsilon\coloneqq\ind_{\{B_{T_a}=a\}} \), the pair \( (T_a,\varepsilon) \) has the same law as \( (T_a,1-\varepsilon) \) by \( B\mapsto -B \), hence \( \varepsilon\sim\mathrm{Bernoulli}(1/2) \) and \( T_a\perp\varepsilon \).

\begin{proof}[Correct proof]
\textbf{(1).}
Paths of \( B \) are continuous, so \( T_a \) is the first hitting time of the closed set \( \{\pm a\} \) and is a stopping time, as in \cref{def:bm-two-sided-exit}. For every \( t>0 \),
\[
\{T_a>t\}
=
\Bigl\{\sup_{s\le t}|B_s|<a\Bigr\}
\subset
\{|B_t|<a\}.
\]
Since \( B_t\sim\mathcal{N}(0,t) \),
\[
\Pbb(|B_t|<a)
=
\Pbb\bigl(|Z|<a/\sqrt{t}\bigr)
\to 0
\qquad
(t\to\infty),
\]
where \( Z\sim\mathcal{N}(0,1) \). Therefore \( \Pbb(T_a>t)\to 0 \), which is equivalent to \( T_a<\infty \) almost surely. (In particular \( \Pbb(T_a=\infty)\le\Pbb(|B_t|<a) \) for every \( t \), so the same limit kills the event \( \{T_a=\infty\} \).)

\textbf{(2).}
By (1) the exit location \( B_{T_a} \) is almost surely well-defined and takes values in \( \{\pm a\} \). Let \( \varepsilon\coloneqq\ind_{\{B_{T_a}=a\}} \). The process \( \widehat{B}_t\coloneqq -B_t \) is again a standard Brownian motion started at \( 0 \). The two walls are symmetric about the origin, so
\[
T_a(\widehat{B})
=
\inf\{t\ge 0:|\widehat{B}_t|=a\}
=
\inf\{t\ge 0:|B_t|=a\}
=
T_a(B),
\]
while \( \widehat{B}_{T_a}=-B_{T_a} \), hence
\[
\ind_{\{\widehat{B}_{T_a}=a\}}
=
\ind_{\{B_{T_a}=-a\}}
=
1-\varepsilon.
\]
The law of \( B \) is invariant under \( B\mapsto\widehat{B} \), therefore
\[
(T_a,\varepsilon)
\;\stackrel{d}{=}\;
(T_a,1-\varepsilon).
\]
For every Borel set \( A\subset[0,\infty) \),
\[
\Pbb(T_a\in A,\,\varepsilon=1)
=
\Pbb(T_a\in A,\,\varepsilon=0).
\]
Adding the two sides gives \( 2\Pbb(T_a\in A,\,\varepsilon=1)=\Pbb(T_a\in A) \). Taking \( A=[0,\infty) \) yields \( \Pbb(\varepsilon=1)=1/2 \). Thus
\[
\Pbb(T_a\in A,\,\varepsilon=1)
=
\Pbb(T_a\in A)\,\Pbb(\varepsilon=1)
\]
for every Borel \( A \), which is \( T_a\perp\varepsilon \).
\end{proof}

\begin{remark}[{Optional: \( \E[T_a]=a^{2} \)}]
\label{rem:p5-mean-exit}
The exam does not ask for the mean. The martingale \( M_t=B_t^{2}-t \) and optional stopping at the bounded time \( T_a\wedge t \) give \( \E[B_{T_a\wedge t}^{2}]=\E[T_a\wedge t] \). The left side is at most \( a^{2} \), so \( \E[T_a]\le a^{2} \) after \( t\to\infty \), which is another proof of (1). Continuity of paths and (1) give \( B_{T_a\wedge t}^{2}\to a^{2} \) almost surely, and bounded convergence upgrades the inequality to \( \E[T_a]=a^{2} \). This is the continuous counterpart of \( \E[\tau]=ab \) in Problem~3 (here the two distances to the walls are both \( a \)).
\end{remark}

\subsection{Problem 6. [15 pts.]}
Suppose \( X \) and \( Y \) are i.i.d with \( \mathbb{E}[X] = 0 \) and \( \mathrm{var}(X) = 1 \). Suppose further that

\[
\mathbb{E}[(X - Y)^2 \mid \sigma(X + Y)] = 2.
\]

Show that \( X \) and \( Y \) are independent standard Gaussian random variables.

\setcounter{section}{6}

\subsubsection{Hint for Problem 6}

Independence of \( X \) and \( Y \) is already in the hypotheses (they are i.i.d.). The claim is that the common law is \( \mathcal{N}(0,1) \). Gaussianity is not assumed: it is the output of a differential equation for the characteristic function \( \varphi(t)=\E[e^{itX}] \).

Write \( S=X+Y \) and \( D=X-Y \). The given condition is \( \E[D^{2}\mid\sigma(S)]=2 \). Unconditionally, independence and the two moments already give
\[
\E[D^{2}]
=
\E[X^{2}]+\E[Y^{2}]-2\E[X]\E[Y]
=
2,
\]
so the assumption is that the conditional second moment of \( D \) given \( S \) is constant, equal to its unconditional mean. If \( X,Y \) are i.i.d.\ \( \mathcal{N}(0,1) \), then \( S\perp D \) (both \( \mathcal{N}(0,2) \)) and the condition holds; that is the easy direction, and is worth one sentence on the exam.

Do not try to prove the stronger Bernstein statement \( S\perp D \) from the second-moment condition alone. The exam path is: multiply the identity \( \E[D^{2}\mid S]=2 \) by \( e^{i\lambda S} \) and take expectations, then factor by independence. Finite second moment makes \( \varphi \) twice continuously differentiable, with
\[
\varphi'(t)=i\E[Xe^{itX}],
\qquad
\varphi''(t)=-\E[X^{2}e^{itX}].
\]
The resulting identity is an ODE for \( \varphi \). In logarithmic form it is
\[
(\log\varphi)''=-1
\]
on the open set where \( \varphi\neq 0 \). The initial conditions \( \varphi(0)=1 \) and \( \varphi'(0)=i\E[X]=0 \) force \( \varphi(t)=e^{-t^{2}/2} \) near the origin. To pass from a neighbourhood of \( 0 \) to all of \( \mathbb{R} \): \( \varphi \) cannot have a first zero, because the Gaussian formula would extend continuously up to that point and remain nonzero. See \cref{prop:chf-indep} for the characteristic-function calculus; uniqueness of characteristic functions then yields \( X\sim\mathcal{N}(0,1) \), and likewise for \( Y \).

Three things not to do: (i) compute the joint density of \( (S,D) \); (ii) quote ``uncorrelated Gaussians are independent'' (there is no Gaussian hypothesis yet; that slogan is Problem~2); (iii) treat \( \E[D^{2}\mid S]=2 \) as if it already said \( S\perp D \).

\subsubsection{How to prove something is Gaussian, on this QE}

The word ``Gaussian'' on the Qiuzhen paper is three different tasks. Mixing them is the usual source of illegal slogans (especially quoting Problem~2 inside Problem~6).

\begin{center}
\renewcommand{\arraystretch}{1.25}
\small
\begin{tabular}{@{}lll@{}}
\toprule
Task & Input & Machine \\
\midrule
Already Gaussian; compute & joint ch.f.\ is \( e^{it^{T}\mu-\frac12 t^{T}\Sigma t} \) & Problem~2; 2026~S~P1 \\
Prove a law \emph{is} Gaussian & ODE / independence of linear forms & this problem \\
Prove a limit \emph{is} Gaussian & \( \varphi_n(t)\to e^{-t^{2}/2} \), or Stein & 2024~S~P1, 2026~S~P2, 2023~F~P6 \\
\bottomrule
\end{tabular}
\end{center}

\paragraph{1. Characteristic function (the default).}

A real random variable is \( \mathcal{N}(\mu,\sigma^{2}) \) if and only if \( \varphi(t)=\exp(it\mu-\sigma^{2}t^{2}/2) \). A vector is Gaussian (possibly degenerate) if and only if the joint characteristic function is \( \exp(it^{T}\mu-\frac12 t^{T}\Sigma t) \); that is \cref{def:gaussian-vector}. The three exam moves are the next three propositions. The density \( (2\pi)^{-d/2}(\det\Sigma)^{-1/2}\exp(-\frac12(x-\mu)^{T}\Sigma^{-1}(x-\mu)) \) is equivalent only when \( \Sigma \) is positive definite: Problem~2 is the warning that on \( \rho^{2}+\tau^{2}=1 \) there is no density, while the characteristic-function definition still works.

\begin{proposition}[Match a closed-form characteristic function]
\label{prop:chf-match-gaussian}
Characteristic functions determine laws: if \( \varphi_X=\varphi_Z \) on \( \mathbb{R} \), then \( X\stackrel{d}{=}Z \) \citep[Theorem~3.3.1]{durrett2010}. Consequently:
\begin{enumerate}[label=\textup{(\roman*)}]
    \item
    \( X\sim\mathcal{N}(\mu,\sigma^{2}) \) if and only if \( \varphi_X(t)=\exp(it\mu-\sigma^{2}t^{2}/2) \).
    \item
    If \( X\perp Y \), then \( \varphi_{X+Y}=\varphi_X\varphi_Y \). In particular a sum of independent Gaussians is Gaussian, with means and variances added.
    \item
    If \( V \) is a Gaussian vector with mean \( \mu \) and covariance \( \Sigma \), and \( A \) is a deterministic matrix, then \( AV \) is Gaussian with mean \( A\mu \) and covariance \( A\Sigma A^{T} \). In particular every linear image \( a^{T}V \) is univariate Gaussian.
\end{enumerate}
This is the machine of Problem~2 and of 2026 Spring Problem~1 (after an exponential tilt, the joint characteristic function is still quadratic, with mean \( \mu+V\alpha \) and the same covariance).
\end{proposition}

\begin{proof}
Uniqueness is \citep[Theorem~3.3.1]{durrett2010}: the inversion formula recovers \( \mathrm{Law}(X) \) from \( \varphi_X \), so two random variables with the same characteristic function are equal in law. Item~(i) is that statement at the Gaussian characteristic function. Item~(ii) is Fubini plus independence: \( \E[e^{it(X+Y)}]=\E[e^{itX}]\E[e^{itY}] \). Item~(iii) is the substitution \( t\mapsto A^{T}t \) in \cref{def:gaussian-vector}:
\[
\E\bigl[e^{it^{T}(AV)}\bigr]
=
\E\bigl[e^{i(A^{T}t)^{T}V}\bigr]
=
\exp\bigl(it^{T}A\mu-\tfrac12 t^{T}A\Sigma A^{T}t\bigr).
\]
The last display is the characteristic function of \( \mathcal{N}(A\mu,A\Sigma A^{T}) \).
\end{proof}

\begin{proposition}[The ODE \( (\log\varphi)''=-\sigma^{2} \)]
\label{prop:chf-ode-gaussian}
Let \( X \) be real with \( \E[X]=0 \), \( \E[X^{2}]=\sigma^{2}\in(0,\infty) \), and write \( \varphi(t)=\E[e^{itX}] \). Then \( \varphi \) is \( C^{2} \), with \( \varphi(0)=1 \), \( \varphi'(0)=0 \), and \( \varphi''(0)=-\sigma^{2} \). If moreover \( \varphi \) never vanishes and \( (\log\varphi)''=-\sigma^{2} \) on \( \mathbb{R} \), then \( \varphi(t)=\exp(-\sigma^{2}t^{2}/2) \), hence \( X\sim\mathcal{N}(0,\sigma^{2}) \). The same conclusion holds if the ODE is known only on a neighbourhood of \( 0 \): \( \varphi \) cannot have a first zero, because the Gaussian formula extends continuously up to that point and remains nonzero.

In the present problem one reaches the ODE from the conditional-moment identity, as follows. Let \( X,Y \) be i.i.d.\ with \( \E[X]=0 \), \( \E[X^{2}]=1 \), and \( \E[(X-Y)^{2}\mid X+Y]=2 \). Then \( \varphi \) satisfies \( (\log\varphi)''=-1 \) wherever \( \varphi\neq 0 \), so \( X\sim\mathcal{N}(0,1) \).
\end{proposition}

\begin{proof}
Finite second moment gives \( \varphi\in C^{2} \) and the three jet identities by differentiating under the expectation. On \( \{\varphi\neq 0\} \) set \( \psi=\log\varphi \) (the continuous branch with \( \psi(0)=0 \)). The ODE \( \psi''=-\sigma^{2} \) integrates to \( \psi(t)=-\sigma^{2}t^{2}/2+at \), and \( \varphi'(0)=0 \) forces \( a=0 \). Thus \( \varphi(t)=\exp(-\sigma^{2}t^{2}/2) \) on the connected component of \( \{\varphi\neq 0\} \) containing \( 0 \). That component is an open interval. If it were not all of \( \mathbb{R} \), there would be a first zero \( t_0 \); the Gaussian formula would then give \( \varphi(t_0)=\exp(-\sigma^{2}t_0^{2}/2)\neq 0 \), a contradiction. \Cref{prop:chf-match-gaussian} identifies the law.

For the present exam: write \( S=X+Y \), \( D=X-Y \), and \( \varphi(t)=\E[e^{itX}] \). The assumption plus \( \E[D^{2}]=2 \) yields \( \E[D^{2}e^{i\lambda S}]=2\varphi(\lambda)^{2} \). Independence expands the left side as
\[
2\varphi(\lambda)\,\E[X^{2}e^{i\lambda X}]-2\bigl(\E[Xe^{i\lambda X}]\bigr)^{2}
=
-2\varphi(\lambda)\varphi''(\lambda)+2\bigl(\varphi'(\lambda)\bigr)^{2},
\]
because \( \E[Xe^{itX}]=-i\varphi'(t) \) and \( \E[X^{2}e^{itX}]=-\varphi''(t) \), and \( (-i\varphi')^{2}=-(\varphi')^{2} \). Therefore \( (\varphi')^{2}-\varphi\varphi''=\varphi^{2} \), or \( \varphi\varphi''-(\varphi')^{2}=-\varphi^{2} \), which is \( (\log\varphi)''=-1 \) on \( \{\varphi\neq 0\} \). The first paragraph gives \( X\sim\mathcal{N}(0,1) \).
\end{proof}

\begin{proposition}[Lévy continuity, Gaussian target]
\label{prop:levy-gaussian}
Let \( (Z_n) \) be real random variables with characteristic functions \( \varphi_n \). If \( \varphi_n(t)\to\exp(-t^{2}/2) \) for every \( t\in\mathbb{R} \), then \( Z_n\Rightarrow\mathcal{N}(0,1) \). This is the special case of Lévy's continuity theorem \citep[Theorem~3.3.6]{durrett2010} used in 2024 Spring Problem~1 and 2026 Spring Problem~2. It identifies a \emph{limit} as Gaussian; it does not identify a single pair \( (X,Y) \).
\end{proposition}

\begin{proof}
The standard tightness estimate from the characteristic function \citep[\S3.3]{durrett2010} reads, for every \( u>0 \),
\[
\Pbb\bigl(|Z_n|>2/u\bigr)
\le
\frac{C}{u}\int_{0}^{u}\bigl(1-\mathrm{Re}\,\varphi_n(t)\bigr)\,dt
\]
with a universal \( C \). Pointwise convergence to \( e^{-t^{2}/2} \) and bounded convergence on \( [0,u] \) send the right-hand side to the same integral against \( 1-\mathrm{Re}\,e^{-t^{2}/2} \), which can be made arbitrarily small by taking \( u \) small (since \( 1-\mathrm{Re}\,e^{-t^{2}/2}\sim t^{2}/2 \)). Thus \( (Z_n) \) is tight. Every subsequence has a further subsequence converging in distribution to some \( Z \). Along that subsequence, Lévy's theorem gives \( \varphi_Z(t)=e^{-t^{2}/2} \), so \( Z\sim\mathcal{N}(0,1) \) by \cref{prop:chf-match-gaussian}. All subsequential limits coincide, hence \( Z_n\Rightarrow\mathcal{N}(0,1) \).
\end{proof}

\paragraph{2. Independence of linear forms (exact characterizations).}

These are the two-variable statements that look like the present problem but are strictly stronger. The exam gives only \( \E[D^{2}\mid S]=2 \), which is the second-moment shadow of \cref{prop:bernstein}, not Bernstein itself; that is why one goes through \cref{prop:chf-ode-gaussian} rather than quoting the next proposition.

\begin{proposition}[Bernstein]
\label{prop:bernstein}
Let \( X,Y \) be independent with \( \E[X^{2}]+\E[Y^{2}]<\infty \). If \( X+Y \) is independent of \( X-Y \), then \( X \) and \( Y \) are Gaussian, and necessarily \( \Var(X)=\Var(Y) \). (The theorem holds without second moments; the proof below is the \( C^{2} \) argument one writes on the exam.)
\end{proposition}

\begin{proof}
Write \( \varphi(t)=\E[e^{itX}] \), \( \psi(t)=\E[e^{itY}] \), and \( S=X+Y \), \( D=X-Y \). Independence of the pair \( (X,Y) \) gives
\[
\E\bigl[e^{isS+itD}\bigr]
=
\varphi(s+t)\psi(s-t).
\]
Independence of \( S \) and \( D \) factors the same characteristic function as
\[
\varphi_S(s)\varphi_D(t)
=
\varphi(s)\psi(s)\,\varphi(t)\psi(-t).
\]
Thus, near the origin (where \( \varphi \) and \( \psi \) have no zeros, since both equal \( 1 \) at \( 0 \)), the logarithms \( f=\log\varphi \) and \( g=\log\psi \) satisfy
\[
f(s+t)+g(s-t)
=
f(s)+g(s)+f(t)+g(-t).
\]
Second moments make \( f,g \) of class \( C^{2} \). Differentiate in \( t \) and then in \( s \):
\[
f'(s+t)-g'(s-t)
=
f'(t)-g'(-t),
\qquad
f''(s+t)-g''(s-t)
=
0.
\]
Hence \( f''(s+t)=g''(s-t) \) for all small \( s,t \), so both second derivatives are equal and constant. Therefore \( f(t)=-\frac12\sigma^{2}t^{2}+i\mu_X t \) and \( g(t)=-\frac12\sigma^{2}t^{2}+i\mu_Y t \) near \( 0 \), after matching the first derivatives to the means. The same no-first-zero argument as in \cref{prop:chf-ode-gaussian} extends the identities to \( \mathbb{R} \). \Cref{prop:chf-match-gaussian} yields that both laws are Gaussian with the same variance. (The variances agree for an a priori reason as well: \( \operatorname{Cov}(S,D)=\Var(X)-\Var(Y) \), and independence of \( S \) and \( D \) forces the covariance to vanish.)
\end{proof}

\begin{proposition}[Cramér]
\label{prop:cramer}
If \( X\perp Y \) and \( X+Y \) is Gaussian, then \( X \) and \( Y \) are each Gaussian. The present problem does not give that the sum is Gaussian.
\end{proposition}

\begin{proof}
Write \( X+Y\sim\mathcal{N}(\mu,\sigma^{2}) \). The moment generating function of a Gaussian exists on all of \( \mathbb{R} \), so \( \E[e^{\lambda(X+Y)}]<\infty \) for every real \( \lambda \). Independence and \( e^{\lambda X}e^{\lambda Y}=e^{\lambda(X+Y)} \) give
\[
\E[e^{\lambda X}]\,\E[e^{\lambda Y}]
=
\E[e^{\lambda(X+Y)}]
<\infty.
\]
Thus both moment generating functions exist on \( \mathbb{R} \). Their logarithms \( \Lambda_X,\Lambda_Y \) are convex and \( C^{\infty} \) on \( \mathbb{R} \), and
\[
\Lambda_X(\lambda)+\Lambda_Y(\lambda)
=
\mu\lambda+\tfrac12\sigma^{2}\lambda^{2}.
\]
Differentiate twice: \( \Lambda_X''+\Lambda_Y''=\sigma^{2} \). Each second derivative is nonnegative (convexity of a log-mgf), and the two sum to a constant, so each is itself a constant. Therefore each \( \Lambda \) is quadratic, hence each law is Gaussian.
\end{proof}

\begin{proposition}[Maxwell]
\label{prop:maxwell}
If \( X,Y \) are i.i.d.\ and the law of \( (X,Y) \) is invariant under rotations of \( \mathbb{R}^{2} \), then the common law is centred Gaussian. Equivalently: the direction \( (X,Y)/|(X,Y)| \) is uniform on the circle and independent of the radius (on \( \{|(X,Y)|>0\} \)). Rare on this QE.
\end{proposition}

\begin{proof}
Rotational invariance plus independence give \( \varphi(s)\varphi(t)=\varphi(\sqrt{s^{2}+t^{2}}) \) for the common characteristic function \( \varphi \), because the left side is \( \E[e^{isX+itY}] \) and the right side is the same expectation after rotating \( (s,t) \) onto \( (\sqrt{s^{2}+t^{2}},0) \), using \( \varphi(0)=1 \). In particular \( \varphi \) is even and real, so \( X \) is symmetric. On a neighbourhood of \( 0 \) one has \( \varphi\neq 0 \); set \( f=\log\varphi \) and \( h(u)=f(\sqrt{u}) \) for \( u\ge 0 \) small. The identity becomes Cauchy's equation \( h(s^{2})+h(t^{2})=h(s^{2}+t^{2}) \). Continuity of \( \varphi \) forces \( h(u)=\kappa u \), hence \( f(s)=\kappa s^{2} \) and \( \varphi(s)=e^{\kappa s^{2}} \). The bound \( |\varphi|\le 1 \) gives \( \kappa\le 0 \). The no-first-zero argument of \cref{prop:chf-ode-gaussian} extends the formula to \( \mathbb{R} \), and \cref{prop:chf-match-gaussian} identifies a centred Gaussian. (If \( \E[X^{2}]<\infty \), one can also differentiate the ODE \( (\log\varphi)''=\mathrm{const} \) directly.)
\end{proof}

\paragraph{3. Already a Gaussian vector (Problem~2 calculus).}
Here Gaussianity is an \emph{input}. The legal consequences, all from the quadratic characteristic function, are:
\begin{itemize}
    \item
    linear images are Gaussian;
    \item
    uncorrelated coordinates are independent (\cref{thm:gauss-uncorr-indep}) --- illegal from margins alone, cf.\ \cref{warn:marginal-gauss-not-joint,ex:signflip};
    \item
    the conditional law of one block given another is Gaussian, with affine conditional mean and constant conditional covariance (the regression formula used to build \( W \) in Problem~2);
    \item
    a Gaussian process is a process whose finite-dimensional laws are Gaussian vectors (Brownian motion in Problem~5; Brownian bridge vs.\ GFF in 2024 Fall Problem~6, identified by matching covariances).
\end{itemize}
None of these proves that a non-Gaussian-looking pair is Gaussian. They compute inside a Gaussian world.

\paragraph{4. Stein's identity (limits, not exact laws).}
A random variable \( Z \) is standard Gaussian if and only if \( \E[f'(Z)-Zf(Z)]=0 \) for every compactly supported \( C^{1} \) test function \( f \). On the QE this is used for \emph{convergence}: if \( \E[f'(Z_n)-Z_nf(Z_n)]\to 0 \) for all such \( f \), then \( Z_n\Rightarrow\mathcal{N}(0,1) \). That is the hint of 2023 Fall Problem~6. It is the wrong tool for the exact identification in the present problem (one would have to check the identity for all \( f \), which is heavier than solving the ODE for \( \varphi \)).

\paragraph{5. Moment generating functions.}
If \( \E[e^{tX}]=\exp(t\mu+\sigma^{2}t^{2}/2) \) on a neighbourhood of \( 0 \), then \( X\sim\mathcal{N}(\mu,\sigma^{2}) \). This is the real-variable twin of the characteristic-function definition, and is the clean way to run 2026 Spring Problem~1 (stay on the real axis; do not analytically continue the characteristic function by hand). It is unavailable for Cauchy and for the infinite-variance walk of 2024 Spring Problem~1.

The single decision tree: if the object is already called a Gaussian vector, use the quadratic characteristic function and \cref{thm:gauss-uncorr-indep}. If the object is a pair of i.i.d.\ random variables plus a relation between sum and difference, write \( \varphi \) and apply \cref{prop:chf-ode-gaussian}, not \cref{prop:bernstein,prop:cramer}. If the object is a scaled sum \( S_n/b_n \), expand \( \varphi_{S_n/b_n} \) and apply \cref{prop:levy-gaussian}, or run Stein. Do not import a theorem from a stronger hypothesis than the paper gave.

\subsubsection{Solution of Problem 6}

\paragraph{Answer.}
Independence is already in the hypotheses. The common law is \( \mathcal{N}(0,1) \). Equivalently, \( \varphi(t)=\E[e^{itX}] \) satisfies \( (\log\varphi)''=-1 \) with \( \varphi(0)=1 \) and \( \varphi'(0)=0 \), hence \( \varphi(t)=e^{-t^{2}/2} \).

\begin{proof}[Correct proof]
Write \( S=X+Y \), \( D=X-Y \), and \( \varphi(t)=\E[e^{itX}] \). Independence of \( X \) and \( Y \) is given, so it remains to identify the common law.

If \( X,Y \) are i.i.d.\ \( \mathcal{N}(0,1) \), the joint characteristic function of \( (S,D) \) is
\[
\E\bigl[e^{isS+itD}\bigr]
=
\exp\bigl(-\tfrac12(s+t)^{2}-\tfrac12(s-t)^{2}\bigr)
=
\exp(-s^{2}-t^{2})
=
\E[e^{isS}]\,\E[e^{itD}],
\]
so \( S\perp D \) with \( D\sim\mathcal{N}(0,2) \). Therefore \( \E[D^{2}\mid S]=\E[D^{2}]=2 \), and the hypothesis holds. This is the easy direction.

For the converse, the given moments and independence already yield \( \E[D^{2}]=2 \). The assumption is thus \( \E[D^{2}\mid\sigma(S)]=\E[D^{2}] \). Multiply by \( e^{i\lambda S} \) and take expectations:
\[
\E[D^{2}e^{i\lambda S}]
=
2\,\E[e^{i\lambda S}]
=
2\,\varphi(\lambda)^{2}.
\]
Expand \( D^{2}=(X-Y)^{2} \) and use independence of \( X \) and \( Y \):
\[
\E[D^{2}e^{i\lambda S}]
=
2\varphi(\lambda)\,\E[X^{2}e^{i\lambda X}]
-
2\bigl(\E[Xe^{i\lambda X}]\bigr)^{2}.
\]
Finite second moment makes \( \varphi \) of class \( C^{2} \), with \( \varphi'(t)=i\E[Xe^{itX}] \) and \( \varphi''(t)=-\E[X^{2}e^{itX}] \). In particular \( \E[Xe^{itX}]=-i\varphi'(t) \) and \( (-i\varphi')^{2}=-(\varphi')^{2} \), so the previous display equals \( -2\varphi\varphi''+2(\varphi')^{2} \). Therefore
\[
(\varphi')^{2}-\varphi\varphi''
=
\varphi^{2},
\qquad\text{i.e.}\qquad
\varphi\varphi''-(\varphi')^{2}
=
-\varphi^{2},
\]
which is \( (\log\varphi)''=-1 \) on the open set \( \{\varphi\neq 0\} \). The initial conditions are \( \varphi(0)=1 \) and \( \varphi'(0)=i\E[X]=0 \). Integrating the ODE on the connected component of \( \{\varphi\neq 0\} \) containing \( 0 \) gives \( \varphi(t)=e^{-t^{2}/2} \) there. This component is all of \( \mathbb{R} \): otherwise there would be a first zero \( t_0 \), but then \( \varphi(t_0)=e^{-t_0^{2}/2}\neq 0 \). \Cref{prop:chf-match-gaussian} yields \( X\sim\mathcal{N}(0,1) \), and likewise \( Y \).

(The argument does not prove \( S\perp D \). That is Bernstein's hypothesis, \cref{prop:bernstein}, which is strictly stronger than \( \E[D^{2}\mid S]=2 \).)
\end{proof}

\subsection{Problem 7. [10 pts.]}
Let \( X_1, \dots, X_n \) (\( n \ge 2 \)) be iid uniform \( U(0, \theta) \) where \( 0 < \theta < \infty \).
\begin{enumerate}[label=(\alph*)]
    \item Under squared error loss \( l(\theta, \delta(X)) = (\delta(X) - \theta)^2 \), derive the generalized Bayes rule for the improper prior

    \[
    \pi(\theta) = 1, \qquad 0 < \theta < +\infty.
    \]
    \item Is the Bayes rule derived in (a) consistent for \( \theta \)?
\end{enumerate}

\setcounter{section}{7}

\subsubsection{Definitions used in Problem 7}

The new vocabulary is decision-theoretic: a (generalized) Bayes rule minimises posterior expected loss. Part~(b) is a different question. Consistency is a frequentist property of a \emph{sequence} of estimators as \( n\to\infty \); bias is a property of the mean at a \emph{fixed} \( n \). The two are not interchangeable. Squared error forces the minimiser in (a) to be the posterior mean, when that mean is finite.

\begin{definition}[Formal posterior, and generalized Bayes rule]
\label{def:gen-bayes}
Let \( f(x\mid\theta) \) be a likelihood and let \( \pi \) be a \(\sigma\)-finite measure on the parameter space (a probability prior is the special case \( \pi(\Theta)=1 \)). The \emph{formal posterior} is the probability measure
\[
\pi(\mathrm{d}\theta\mid x)
\propto
f(x\mid\theta)\,\pi(\mathrm{d}\theta),
\]
whenever the right-hand side has finite total mass. A decision rule \( \delta \) is a \emph{generalized Bayes rule} for a loss \( \ell(\theta,a) \) if, for almost every \( x \), the action \( a=\delta(x) \) minimises the posterior expected loss
\[
\rho(a\mid x)
\coloneqq
\int \ell(\theta,a)\,\pi(\mathrm{d}\theta\mid x).
\]
When \( \pi \) is a probability this is an ordinary Bayes rule. The adjective ``generalized'' records that \( \pi(\theta)=1 \) on \( (0,\infty) \) is not a probability; the computation is the same once the formal posterior has been normalised. See \citep[\S11.1]{hogg2019}.
\end{definition}

\begin{proposition}[Squared error \(\mapsto\) posterior mean]
\label{prop:l2-posterior-mean}
If \( \ell(\theta,a)=(a-\theta)^{2} \) and the formal posterior has finite second moment, then
\[
\rho(a\mid x)
=
\bigl(a-\E[\theta\mid x]\bigr)^{2}
+
\Var(\theta\mid x),
\]
so any generalized Bayes rule is the posterior mean \( \delta(x)=\E[\theta\mid x] \). If the posterior mean is infinite, no real-valued generalized Bayes rule exists. This is \citep[\S11.1]{hogg2019}: \( \E[(W-b)^{2}] \) is minimised at \( b=\E[W] \) whenever the second moment is finite.
\end{proposition}

\begin{definition}[Convergence in probability]
\label{def:conv-p}
Let \( (Z_n)_{n\ge 1} \) and \( Z \) be random variables on a common probability space. One says that \( Z_n \) \emph{converges in probability} to \( Z \), written \( Z_n\xrightarrow{\Pbb}Z \), if
\[
\lim_{n\to\infty}\Pbb\bigl(|Z_n-Z|\ge\varepsilon\bigr)
=
0
\qquad\text{for every }\varepsilon>0,
\]
equivalently \( \Pbb\bigl(|Z_n-Z|<\varepsilon\bigr)\to 1 \). If \( Z \) is the constant \( a \), one writes \( Z_n\xrightarrow{\Pbb}a \). This is \citep[Definition~5.1.1]{hogg2019}. The statement is about the \emph{law} of \( Z_n-Z \) concentrating at \( 0 \). It does not mention \( \E[Z_n] \).
\end{definition}

\begin{definition}[Consistency]
\label{def:consistency}
Let \( \{f(\,\cdot\mid\theta):\theta\in\Omega\} \) be a parametric family, and let \( X_1,X_2,\dots \) be i.i.d.\ with law \( f(\,\cdot\mid\theta) \). A sequence of statistics \( T_n=T_n(X_1,\dots,X_n) \) is a \emph{consistent estimator} of \( g(\theta) \) if
\[
T_n
\xrightarrow{\Pbb_\theta}
g(\theta)
\qquad\text{for every }\theta\in\Omega
\]
in the sense of \cref{def:conv-p}: for every \( \theta\in\Omega \) and every \( \varepsilon>0 \),
\[
\lim_{n\to\infty}
\Pbb_\theta\bigl(\bigl|T_n-g(\theta)\bigr|\ge\varepsilon\bigr)
=
0.
\]
When \( g \) is the identity one says that \( T_n \) is consistent for \( \theta \). This is \citep[Definition~5.1.2]{hogg2019}. The subscript \( \theta \) on \( \Pbb_\theta \) records that the same sequence must concentrate at the true value, whichever value is true.

Three structural points. First, the object is a \emph{sequence} \( (T_n)_{n\ge n_0} \): one estimator for each sample size. A single function of a fixed-\( n \) sample cannot be called consistent. Second, the mode is convergence in probability, not convergence of expectations. Third, the quantifier is \( \forall\theta\in\Omega \). The exam's part~(b) is exactly this definition with \( T_n=\delta_n^{\pi} \) and \( g(\theta)=\theta \).

If \( T_n\to g(\theta) \) holds \( \Pbb_\theta \)-almost surely for every \( \theta \), one says that \( T_n \) is \emph{strongly consistent}. Strong consistency implies consistency. The exam writes only ``consistent'', which is the in-probability statement.
\end{definition}

\begin{definition}[Bias, and unbiasedness]
\label{def:bias}
Let \( T=T(X_1,\dots,X_n) \) be an estimator of \( g(\theta) \) at a \emph{fixed} sample size \( n \), and assume \( \E_\theta|T|<\infty \). The \emph{bias} of \( T \) at \( \theta \) is
\[
\operatorname{Bias}_\theta(T)
\coloneqq
\E_\theta[T]-g(\theta).
\]
One says that \( T \) is \emph{unbiased} for \( g(\theta) \) if \( \operatorname{Bias}_\theta(T)=0 \) for every \( \theta\in\Omega \), i.e.\ \( \E_\theta[T]=g(\theta) \) identically. See \citep[\S2.8, \S5.1.1]{hogg2019}. This is a statement about the mean of the sampling distribution at that \( n \). It does not mention \( n\to\infty \), and it does not mention \( \Pbb\bigl(|T-g(\theta)|\ge\varepsilon\bigr) \).
\end{definition}

\begin{remark}[The two properties do not imply each other]
\label{rem:consist-vs-bias}
Consistency is \cref{def:consistency}. Bias is \cref{def:bias}. Neither is a rewriting of the other.
\begin{itemize}
    \item
    Unbiased does not imply consistent. The estimator \( T_n=X_1 \) is unbiased for \( \E[X_1] \) at every \( n \), and \( \Pbb\bigl(|T_n-\theta|\ge\varepsilon\bigr) \) does not tend to \( 0 \).
    \item
    Consistent does not imply unbiased. On the present model, \( M_n=X_{(n)} \) satisfies \( \E_\theta[M_n]=\frac{n}{n+1}\theta\neq\theta \), but \( M_n\xrightarrow{\Pbb_\theta}\theta \) by the exact tail in \cref{lem:unif-max}. This is \citep[Example~5.1.2]{hogg2019}. The generalized Bayes rule \( \delta_n=\frac{n-1}{n-2}M_n \) is the same phenomenon: \( \operatorname{Bias}_\theta(\delta_n)=\frac{2\theta}{(n-2)(n+1)}\neq 0 \), while \( \delta_n\xrightarrow{\Pbb_\theta}\theta \).
    \item
    A sufficient (not necessary) criterion that uses both moments is: if \( \E_\theta[T_n]\to g(\theta) \) and \( \Var_\theta(T_n)\to 0 \), then Chebyshev gives consistency. Vanishing bias is then an input, not a replacement for the probability statement. Computing only \( \E[T_n] \) does not decide \cref{def:consistency}.
\end{itemize}
\end{remark}

\begin{lemma}[Sample maximum of \( U(0,\theta) \)]
\label{lem:unif-max}
Let \( X_1,\dots,X_n \) be i.i.d.\ \( U(0,\theta) \) and write \( M_n\coloneqq X_{(n)}=\max_i X_i \). Then \( M_n \) has density \( n y^{n-1}\theta^{-n} \) on \( (0,\theta) \), hence
\[
\E_\theta[M_n]
=
\frac{n}{n+1}\,\theta,
\qquad
\Pbb_\theta(M_n\le t)
=
\bigl(t/\theta\bigr)^{n}
\quad(0<t\le\theta).
\]
In particular \( M_n\xrightarrow{\Pbb_\theta}\theta \). This is \citep[Example~5.1.2]{hogg2019}.
\end{lemma}

\begin{proof}
The cdf of \( M_n \) is \( \Pbb(X_1\le t)^{n}=(t/\theta)^{n} \) on \( (0,\theta] \). Differentiate to get the density. The mean is
\[
\E[M_n]
=
\int_{0}^{\theta} y\cdot n y^{n-1}\theta^{-n}\,\mathrm{d}y
=
\frac{n}{n+1}\,\theta.
\]
For \( 0<\varepsilon<\theta \),
\[
\Pbb_\theta\bigl(|M_n-\theta|\ge\varepsilon\bigr)
=
\Pbb_\theta(M_n\le\theta-\varepsilon)
=
\Bigl(\frac{\theta-\varepsilon}{\theta}\Bigr)^{n}
\to 0.
\]
\end{proof}

\subsubsection{Manuscript proof idea}

The calculation below follows the handwritten notes. \mserr{Red} marks a false claim or a missing justification. The skeleton is the right exam machine: write the formal posterior, minimise posterior expected squared error, then discuss the resulting estimator as \( n\to\infty \). Three concrete errors sit on that skeleton.

\textbf{(a), the objective.}
The notes open with Bayes rule and write
\[
a
=
\argmin_{\delta}
\E\Bigl(\sum_{i=1}^{n}\ell(x_i;\theta)\Bigm|\vec{X}\Bigr)
=
\sum_{i=1}^{n}\int \ell(x_i;\theta)\,\mathrm{d}\pi(\theta\mid\vec{X}).
\]
\mserr{The loss on the exam is \( \ell(\theta,\delta(X))=(\delta(X)-\theta)^{2} \), a single action against a single parameter, not a sum of per-observation losses.} The subsequent identification of \( a \) with the objective (rather than with its minimiser) is the same slip. The later first-order condition in \( \delta \) shows that the intended object was nevertheless posterior expected squared error, which is the right functional.

\textbf{(a), the formal posterior.}
The notes write
\[
\pi(\theta\mid\vec{X})
\propto
f(\vec{X}\mid\theta)\,\pi(\theta)
=
\theta^{-n}\,\ind_{\{X_i\in(0,\theta)\ \forall i\}}.
\]
This is correct: the likelihood of i.i.d.\ \( U(0,\theta) \) is \( \theta^{-n} \) on \( \{0<X_i<\theta\text{ for all }i\} \), i.e.\ on \( \{\theta\ge M\} \) with \( M\coloneqq\max_i X_i \), and the improper prior is Lebesgue. The normalising integral is
\[
\int_{M}^{\infty}\theta^{-n}\,\mathrm{d}\theta
=
\frac{M^{1-n}}{n-1}
\qquad(n>1),
\]
so the formal posterior is the Pareto density
\[
\pi(\theta\mid\vec{X})
=
(n-1)\,M^{n-1}\,\theta^{-n}\,\ind_{\{\theta>M\}}.
\]
The notes instead record
\[
\mserr{\pi(\theta\mid\vec{X})=\frac{n-1}{M^{n-1}}\,\theta^{-n}\,\ind_{\{\theta>M\}}}.
\]
The power of \( M \) is inverted. (That un-normalised function does not integrate to \( 1 \).) The shape \( \theta\mapsto\theta^{-n}\ind_{\{\theta>M\}} \) is right, so a correctly executed first-order condition would still recover the posterior mean.

\textbf{(a), the minimiser.}
The notes differentiate
\[
\int_{M}^{\infty}(\delta-\theta)^{2}\,\frac{n-1}{M^{n-1}}\,\theta^{-n}\,\mathrm{d}\theta
\]
in \( \delta \) and set the derivative to zero. That is the correct calculus for \cref{prop:l2-posterior-mean}: the critical point is the mean of whatever measure is being integrated against. They conclude
\[
\mserr{\delta=\frac{n-1}{n-2}\,\frac{1}{M}}.
\]
The coefficient \( \frac{n-1}{n-2} \) is the right Pareto factor (for \( n>2 \)), but it multiplies \( M \), not \( M^{-1} \). Explicitly, even against the inverted constant,
\[
\frac{\displaystyle\int_{M}^{\infty}\theta\cdot\theta^{-n}\,\mathrm{d}\theta}{\displaystyle\int_{M}^{\infty}\theta^{-n}\,\mathrm{d}\theta}
=
\frac{M^{2-n}/(n-2)}{M^{1-n}/(n-1)}
=
\frac{n-1}{n-2}\,M.
\]
The reciprocal is an extra algebraic slip, not a consequence of the inverted normaliser.

\textbf{(b).}
The notes compute the law of \( M \) correctly, \( \Pbb(M\le y\mid\theta)=(y/\theta)^{n} \) with density \( n y^{n-1}\theta^{-n} \) on \( (0,\theta) \), and then
\[
\E\bigl[M^{-1}\bigr]
=
\int_{0}^{\theta} y^{-1}\cdot n y^{n-1}\theta^{-n}\,\mathrm{d}y
=
\frac{n}{n-1}\,\frac{1}{\theta},
\]
hence \( \E[\delta]=\frac{n}{n-2}\,\theta^{-1} \) for their estimator. The integral is correct for \( M^{-1} \). \mserr{Part~(b) asks whether \( \delta_n\xrightarrow{\Pbb}\theta \), not whether \( \E[\delta_n]=\theta \).} An expectation of order \( \theta^{-1} \) does not answer consistency, and the object whose expectation was computed is not the rule of part~(a).

\errpanel{%
\textbf{\mserr{The error, in one box.}}

The loss is \( (\delta-\theta)^{2} \), not \( \sum_{i}\ell(x_i;\theta) \). Under that loss the generalized Bayes rule is the posterior mean.

The formal posterior is \( (n-1)M^{n-1}\theta^{-n} \) on \( \theta>M \), not \( (n-1)M^{-(n-1)}\theta^{-n} \). Then \( \E[\theta\mid X]=\frac{n-1}{n-2}\,M \) for \( n\ge 3 \).

Part~(b) is \( \delta_n\xrightarrow{\Pbb}\theta \). Use \( M_n\xrightarrow{\Pbb}\theta \) from the exact cdf and \( \frac{n-1}{n-2}\to 1 \). Do not compute \( \E[\delta] \).
}

\subsubsection{Solution of Problem 7}

\paragraph{Answer.}
Write \( M=X_{(n)} \). For \( n\ge 3 \) the generalized Bayes rule is
\[
\delta^{\pi}(X)
=
\frac{n-1}{n-2}\,M.
\]
For \( n=2 \) the formal posterior is proper but has infinite mean, so no real-valued generalized Bayes rule exists. The sequence \( \delta_n=\frac{n-1}{n-2}\,X_{(n)} \) (\( n\ge 3 \)) is consistent for \( \theta \): \( \delta_n\xrightarrow{\Pbb_\theta}\theta \) for every \( \theta>0 \). (In fact \( \delta_n\to\theta \) almost surely.)

\begin{proof}[Correct proof]
\textbf{(a).}
Write \( x=(x_1,\dots,x_n) \) and \( M\coloneqq\max_i x_i \). The joint density of the sample is
\[
f(x\mid\theta)
=
\theta^{-n}\ind_{\{0<x_i<\theta\ \forall i\}}
=
\theta^{-n}\ind_{\{M<\theta\}}\ind_{\{\min_i x_i>0\}}.
\]
The prior \( \pi(\theta)=1 \) on \( (0,\infty) \) is improper. The formal posterior measure of \cref{def:gen-bayes} has un-normalised density \( \theta^{-n} \) on \( \theta>M \). Its total mass is
\[
\int_{M}^{\infty}\theta^{-n}\,\mathrm{d}\theta
=
\Bigl[\frac{\theta^{1-n}}{1-n}\Bigr]_{M}^{\infty}
=
\frac{M^{1-n}}{n-1}.
\]
This is finite for every \( n>1 \), hence for every \( n\ge 2 \) allowed by the exam. The formal posterior is therefore the probability density
\[
\pi(\theta\mid x)
=
(n-1)\,M^{n-1}\,\theta^{-n},
\qquad
\theta>M
\]
(Pareto with scale \( M \) and shape \( n-1 \)). Check:
\[
\int_{M}^{\infty}(n-1)\,M^{n-1}\,\theta^{-n}\,\mathrm{d}\theta
=
(n-1)\,M^{n-1}\cdot\frac{M^{1-n}}{n-1}
=
1.
\]

Squared-error posterior expected loss is
\[
\rho(a\mid x)
=
\int_{M}^{\infty}(a-\theta)^{2}\,(n-1)\,M^{n-1}\,\theta^{-n}\,\mathrm{d}\theta,
\qquad
a\in\mathbb{R}.
\]
The integrand is a quadratic polynomial in \( a \) with nonnegative coefficients on the posterior, so \( a\mapsto\rho(a\mid x) \) is a convex function \( [0,+\infty]\to[0,+\infty] \). Whenever \( \rho(\,\cdot\mid x) \) is finite at even one point it is \( C^{\infty} \) and
\[
\frac{\partial\rho}{\partial a}(a\mid x)
=
2\int_{M}^{\infty}(a-\theta)\,\pi(\theta\mid x)\,\mathrm{d}\theta
=
2\bigl(a-\E[\theta\mid x]\bigr),
\qquad
\frac{\partial^{2}\rho}{\partial a^{2}}
=
2.
\]
The unique critical point is therefore the posterior mean, and it is a minimum. This is \cref{prop:l2-posterior-mean}, written out. It remains to compute \( \E[\theta\mid x] \) and to record when the integrals converge.

The posterior mean is
\[
\E[\theta\mid x]
=
(n-1)\,M^{n-1}\int_{M}^{\infty}\theta^{1-n}\,\mathrm{d}\theta.
\]
For \( n>2 \) the antiderivative is \( \theta^{2-n}/(2-n) \), hence
\[
\int_{M}^{\infty}\theta^{1-n}\,\mathrm{d}\theta
=
\frac{M^{2-n}}{n-2}
\]
and
\[
\E[\theta\mid x]
=
(n-1)\,M^{n-1}\cdot\frac{M^{2-n}}{n-2}
=
\frac{n-1}{n-2}\,M.
\]
Thus \( \delta^{\pi}(x)=\frac{n-1}{n-2}\,X_{(n)} \) for \( n\ge 3 \). For \( n=2 \) the same integral is \( \int_{M}^{\infty}\theta^{-1}\,\mathrm{d}\theta=\infty \), so \( \E[\theta\mid x]=\infty \) and \( \rho(a\mid x)=\infty \) for every real \( a \). No real-valued generalized Bayes rule exists at \( n=2 \).

(The posterior second moment \( \E[\theta^{2}\mid x]=(n-1)M^{n-1}\int_{M}^{\infty}\theta^{2-n}\,\mathrm{d}\theta \) is finite if and only if \( n>3 \). For \( n=3 \) one still has \( \E[\theta\mid x]=2M \), but \( \rho(a\mid x)=\infty \) for every finite \( a \). The exam's formula \( \frac{n-1}{n-2}M \) is the formal posterior mean for all \( n\ge 3 \); uniqueness of the \( L^{2} \) minimiser uses \( n\ge 4 \).)

\textbf{(b).}
Part~(b) is a statement about the \emph{sequence} of rules from (a), indexed by the sample size. For each \( n\ge 3 \) write \( M_n\coloneqq X_{(n)} \) and
\[
\delta_n
\coloneqq
\frac{n-1}{n-2}\,M_n
=
c_n M_n,
\qquad
c_n
\coloneqq
\frac{n-1}{n-2}.
\]
Fix an arbitrary \( \theta\in(0,\infty) \). One must show \( \delta_n\xrightarrow{\Pbb_\theta}\theta \) in the sense of \cref{def:consistency}: \( \Pbb_\theta\bigl(|\delta_n-\theta|\ge\varepsilon\bigr)\to 0 \) for every \( \varepsilon>0 \). Unbiasedness is not asked, and will fail.

\emph{Law of \( M_n \).}
Always \( 0<M_n\le\theta \) almost surely. For \( 0<t\le\theta \),
\[
\Pbb_\theta(M_n\le t)
=
\Pbb_\theta(X_1\le t)^{n}
=
\Bigl(\frac{t}{\theta}\Bigr)^{n},
\]
so \( M_n \) has density \( n y^{n-1}\theta^{-n} \) on \( (0,\theta) \). This is \cref{lem:unif-max}.

\emph{Consistency of \( M_n \).}
The only way \( M_n \) can fail to be close to \( \theta \) is by falling short: for \( 0<\varepsilon<\theta \),
\[
\Pbb_\theta\bigl(|M_n-\theta|\ge\varepsilon\bigr)
=
\Pbb_\theta(M_n\le\theta-\varepsilon)
=
\Bigl(\frac{\theta-\varepsilon}{\theta}\Bigr)^{n}
\to 0,
\]
because \( (\theta-\varepsilon)/\theta\in(0,1) \). If \( \varepsilon\ge\theta \) the same probability is zero for every \( n \). Hence \( M_n\xrightarrow{\Pbb_\theta}\theta \).

\emph{The prefactor.}
\( c_n=1+\frac{1}{n-2}\to 1 \). In particular there is \( N_0 \) such that \( n\ge N_0 \) implies \( 1\le c_n\le 2 \).

\emph{Product.}
The triangle inequality and \( M_n\le\theta \) give
\[
\bigl|\delta_n-\theta\bigr|
=
\bigl|c_n M_n-\theta\bigr|
\le
c_n\bigl|M_n-\theta\bigr|
+
\bigl|c_n-1\bigr|\,\theta.
\]
Fix \( \varepsilon>0 \). Choose \( N_1\ge N_0 \) so that \( n\ge N_1 \) implies \( |c_n-1|\,\theta<\varepsilon/2 \). Then, on \( \{n\ge N_1\} \),
\[
\bigl\{\bigl|\delta_n-\theta\bigr|\ge\varepsilon\bigr\}
\subset
\bigl\{c_n\bigl|M_n-\theta\bigr|\ge\varepsilon/2\bigr\}
\subset
\bigl\{\bigl|M_n-\theta\bigr|\ge\varepsilon/4\bigr\},
\]
the last inclusion using \( c_n\le 2 \). Therefore
\[
\Pbb_\theta\bigl(|\delta_n-\theta|\ge\varepsilon\bigr)
\le
\Pbb_\theta\bigl(|M_n-\theta|\ge\varepsilon/4\bigr)
=
\Bigl(\frac{\theta-\varepsilon/4}{\theta}\Bigr)^{n}
\to 0
\]
as soon as \( \varepsilon<4\theta \); if \( \varepsilon\ge 4\theta \) the same bound is eventually zero because \( |c_n-1|\theta\to 0 \) already forces \( |\delta_n-\theta|<\varepsilon \) once \( M_n \) is within \( \theta/2 \) of \( \theta \). This is \( \delta_n\xrightarrow{\Pbb_\theta}\theta \).

\emph{Second write-up: \( L^{2} \) via mean and variance.}
The same density gives
\[
\E_\theta[M_n]
=
\int_{0}^{\theta} y\cdot n y^{n-1}\theta^{-n}\,\mathrm{d}y
=
\frac{n}{n+1}\,\theta,
\qquad
\E_\theta[M_n^{2}]
=
\int_{0}^{\theta} y^{2}\cdot n y^{n-1}\theta^{-n}\,\mathrm{d}y
=
\frac{n}{n+2}\,\theta^{2},
\]
hence
\[
\Var_\theta(M_n)
=
\frac{n}{n+2}\,\theta^{2}
-
\Bigl(\frac{n}{n+1}\Bigr)^{2}\theta^{2}
=
\frac{n}{(n+1)^{2}(n+2)}\,\theta^{2}.
\]
(The last identity is \( n(n+1)^{2}-n^{2}(n+2)=n \).) Therefore
\[
\E_\theta[\delta_n]
=
c_n\cdot\frac{n}{n+1}\,\theta
=
\frac{n(n-1)}{(n-2)(n+1)}\,\theta
\to
\theta
\]
and
\[
\Var_\theta(\delta_n)
=
c_n^{2}\,\Var_\theta(M_n)
=
\frac{(n-1)^{2}n}{(n-2)^{2}(n+1)^{2}(n+2)}\,\theta^{2}
\to
0.
\]
Chebyshev then yields, for every \( \varepsilon>0 \),
\[
\Pbb_\theta\bigl(|\delta_n-\theta|\ge\varepsilon\bigr)
\le
\frac{\E_\theta[(\delta_n-\theta)^{2}]}{\varepsilon^{2}}
=
\frac{\Var_\theta(\delta_n)+\bigl(\E_\theta[\delta_n]-\theta\bigr)^{2}}{\varepsilon^{2}}
\to 0,
\]
which is the same conclusion. The bias term is
\[
\E_\theta[\delta_n]-\theta
=
\frac{2\theta}{(n-2)(n+1)}
\neq
0,
\]
so \( \delta_n \) is biased for every finite \( n\ge 3 \); bias tending to zero, together with variance tending to zero, is enough for consistency. Computing only \( \E[\delta_n] \) is not.

\emph{Third write-up: almost sure.}
Fix \( 0<\varepsilon<\theta \) and write \( r=(\theta-\varepsilon)/\theta\in(0,1) \). Then \( \sum_{n}\Pbb_\theta(M_n\le\theta-\varepsilon)=\sum_{n}r^{n}<\infty \). Borel--Cantelli I gives \( \Pbb_\theta(M_n\le\theta-\varepsilon\text{ i.o.})=0 \). Intersecting over a countable sequence \( \varepsilon_k\downarrow 0 \) yields \( M_n\to\theta \) almost surely. Combined with \( c_n\to 1 \) one has \( \delta_n\to\theta \) almost surely, which is stronger than consistency.
\end{proof}

\begin{remark}[What (b) is not]
\label{rem:p7-bias}
The three comparisons of \( M_n \) on this model are the MLE \( M_n \), the UMVUE \( \frac{n+1}{n}M_n \), and the generalized Bayes rule \( \frac{n-1}{n-2}M_n \). All three are consistent, by the same tail (or the same \( L^{2} \)) argument: each is \( a_n M_n \) with \( a_n\to 1 \). None of the three except the UMVUE is unbiased. The notes' computation of \( \E[M_n^{-1}] \) answers a different question for a different estimator.

The extreme-value scale is \( n(\theta-M_n)\xrightarrow{(d)}\mathrm{Exp}(1/\theta) \), so \( M_n-\theta=O_{\Pbb}(n^{-1}) \). The Bayes prefactor satisfies \( n(c_n-1)=n/(n-2)\to 1 \), which is the same order; consistency is not a statement about that rate.
\end{remark}

\subsection{Problem 8. [10 pts.]}
Let \( X_1, \dots, X_n \) be iid Gamma\( (2025, \theta) \) with \( \theta > 0 \) unknown. Recall that the probability density function of a Gamma\( (\alpha, \theta) \) is \( f(x \mid \alpha, \theta) = \frac{1}{\Gamma(\alpha) \theta^\alpha} x^{\alpha-1} e^{-x/\theta} \), \( x > 0 \), where \( \Gamma(\alpha) = (\alpha - 1)! \) for a positve integer \( \alpha \).
\begin{enumerate}[label=(\alph*)]
    \item Consider the hypothesis testing problem \( H_0 : \theta = 1 \) vs.\ \( H_1 : \theta \neq 1 \). Derive the UMP level \( \alpha = 0.05 \) test if it exists. If it doesn't exist, please argue why it does not exist.
    \item Derive the MLE for \( \frac{1}{\theta} \), what is the asymptotic distribution of this MLE for \( \frac{1}{\theta} \)?
\end{enumerate}

\setcounter{section}{8}

\subsubsection{Hint for Problem 8(a)}

Do not write an equal-tailed two-sided rejection region and call it UMP. The exam alternative is \( H_1:\theta\neq 1 \), and Karlin--Rubin produces a UMP test only for a \emph{one-sided} alternative. The expected write-up of (a) is: exhibit the monotone likelihood ratio, record the two opposite Neyman--Pearson tails, and conclude that no single size-\( 0.05 \) region is most powerful against every \( \theta\neq 1 \).

The shape parameter \( \alpha=2025 \) is known. The joint density is
\[
L(\theta,\mathbf{x})
=
\bigl(\Gamma(\alpha)\bigr)^{-n}\,\theta^{-n\alpha}\Bigl(\prod_{i=1}^{n}x_i^{\alpha-1}\Bigr)\exp\bigl(-S/\theta\bigr),
\qquad
S
\coloneqq
\sum_{i=1}^{n}X_i,
\]
on \( \{x_i>0\} \). All \( \theta \)-dependence sits in \( \theta^{-n\alpha}e^{-S/\theta} \). For \( \theta_1<\theta_2 \) the likelihood ratio \( L(\theta_1)/L(\theta_2) \) is a monotone function of \( S \) (\cref{def:mlr}). Decide the direction before quoting Karlin--Rubin.

Against a simple alternative \( \theta''>1 \), Neyman--Pearson rejects for large \( S \). Against \( \theta''<1 \), it rejects for small \( S \). Those two critical regions are not equal (up to a null set) at level \( 0.05 \). A UMP test of \( H_0:\theta=1 \) versus \( H_1:\theta\neq 1 \) would have to be a best critical region of size \( 0.05 \) for \emph{every} simple pair \( (1,\theta'') \) with \( \theta''\neq 1 \) (\cref{def:ump-size-alpha}). That is impossible once the most powerful region flips with the sign of \( \theta''-1 \). This is Hogg's two-sided normal-mean example \citep[Example~8.2.3]{hogg2019}, with \( S \) in place of \( \overline{X} \).

Three things not to do: (i) invent a two-sided UMP by taking both tails at once; (ii) replace the question by the one-sided pair \( H_0:\theta\le 1 \) versus \( H_1:\theta>1 \), for which a UMP test \emph{does} exist (\cref{thm:karlin-rubin}); (iii) compute a \( p \)-value or a Wald interval. Size is \cref{def:test-size}; here the null is a singleton, so size is just \( \Pbb_{\theta=1}(\mathbf{X}\in C) \).

The one-sided tests are worth one sentence as contrast: \( S\sim\mathrm{Gamma}(n\alpha,\theta) \), and at \( \theta=1 \) one can mark the \( 0.05 \)-quantile of that Gamma (or of \( 2S\sim\chi^{2}_{2n\alpha} \)). That constructs the Karlin--Rubin test against \( \{\theta>1\} \) or against \( \{\theta<1\} \). It does not produce a UMP test against \( \{\theta\neq 1\} \).

\begin{definition}[Size / significance level]
\label{def:test-size}
Let \( H_0:\theta\in\omega_0 \) versus \( H_1:\theta\in\omega_1 \), with \( \omega_0\cap\omega_1=\emptyset \). The \emph{size} (also called the \emph{significance level}) of a test with critical region \( C \) is
\[
\alpha
\coloneqq
\max_{\theta\in\omega_0}\Pbb_{\theta}(\mathbf{X}\in C)
=
\max_{\theta\in\omega_0}\gamma_C(\theta),
\]
where \( \gamma_C(\theta)=\Pbb_{\theta}(\mathbf{X}\in C) \) is the power function. This is \citep[(8.1.3)]{hogg2019}. When \( \omega_0=\{\theta'\} \) is a singleton the maximum is just \( \Pbb_{\theta'}(\mathbf{X}\in C) \). In part~(a) one has \( \theta'=1 \) and the prescribed size \( 0.05 \).
\end{definition}

\begin{definition}[UMP test of size \( \alpha \)]
\label{def:ump-size-alpha}
A critical region \( C \) is a \emph{uniformly most powerful (UMP) critical region of size \( \alpha \)} for a simple null \( H_0:\theta=\theta' \) against a (possibly composite) alternative \( H_1 \) if \( \Pbb_{\theta'}(\mathbf{X}\in C)=\alpha \) and, for every \( \theta''\in H_1 \), \( C \) maximises \( \Pbb_{\theta''}(\mathbf{X}\in C) \) among all size-\( \alpha \) regions. See \citep[Definition~8.2.1]{hogg2019}. Uniformly means one and the same \( C \) must be Neyman--Pearson most powerful at every point of \( H_1 \).
\end{definition}

\begin{definition}[Monotone likelihood ratio]
\label{def:mlr}
The likelihood \( L(\theta,\mathbf{x}) \) has \emph{monotone likelihood ratio} in the statistic \( y=u(\mathbf{x}) \) if, for every \( \theta_1<\theta_2 \), the ratio \( L(\theta_1,\mathbf{x})/L(\theta_2,\mathbf{x}) \) is a monotone function of \( y \). See \citep[Definition~8.2.2]{hogg2019}. If the ratio is decreasing in \( y \), large \( y \) favours large \( \theta \).
\end{definition}

\begin{theorem}[Karlin--Rubin]
\label{thm:karlin-rubin}
Assume \( L(\theta,\mathbf{x}) \) has monotone decreasing likelihood ratio in \( Y=u(\mathbf{X}) \), in the sense of \cref{def:mlr}. Fix \( \theta' \) and \( \alpha\in(0,1) \), and choose \( c_Y \) so that \( \Pbb_{\theta'}(Y\ge c_Y)=\alpha \).
\begin{enumerate}[label=\textup{(\roman*)}]
    \item
    For \( H_0:\theta\le\theta' \) versus \( H_1:\theta>\theta' \), reject \( H_0 \) if \( Y\ge c_Y \). This test is UMP of level \( \alpha \).
    \item
    For \( H_0:\theta\ge\theta' \) versus \( H_1:\theta<\theta' \), reject \( H_0 \) if \( Y\le c_Y' \), with \( \Pbb_{\theta'}(Y\le c_Y')=\alpha \). This test is UMP of level \( \alpha \).
\end{enumerate}
In both cases the power function is monotone, so the size over the composite null equals the Type~I error at the boundary \( \theta' \). This is \citep[(8.2.3), (8.2.6)]{hogg2019}. The theorem does \emph{not} cover \( H_1:\theta\neq\theta' \).
\end{theorem}

\subsubsection{Hint for Problem 8(b), and the delta method}

Part~(b) is invariance plus a one-dimensional delta method. First maximise \( L(\theta,\mathbf{x}) \) to get the MLE \( \hat{\theta} \) of the scale; the MLE of \( 1/\theta \) is then \( 1/\hat{\theta} \). The CLT supplies \( \sqrt{n}(\hat{\theta}-\theta)\xrightarrow{(d)}\mathcal{N}(0,\sigma^{2}(\theta)) \). The map \( g(u)=1/u \) is \( C^{1} \) on \( (0,\infty) \), so \cref{thm:delta-method}(i) converts that limit into the law of \( \sqrt{n}\bigl(1/\hat{\theta}-1/\theta\bigr) \). Problem~9 uses the same theorem in several dimensions, part~(ii). Do not derive the exact density of \( 1/\overline{X} \).

\begin{theorem}[Delta method]
\label{thm:delta-method}
\label{thm:delta-method-mv}
\begin{enumerate}[label=\textup{(\roman*)}]
    \item
    \emph{One dimension.} Let \( (W_n)_{n\ge 1} \) be real and let \( \theta\in\mathbb{R} \), \( \sigma^{2}\ge 0 \) satisfy
    \[
    \sqrt{n}\,(W_n-\theta)
    \xrightarrow{(d)}
    \mathcal{N}(0,\sigma^{2}).
    \]
    If \( g \) is differentiable at \( \theta \), then
    \[
    \sqrt{n}\,\bigl(g(W_n)-g(\theta)\bigr)
    \xrightarrow{(d)}
    \mathcal{N}\bigl(0,\,\sigma^{2}\bigl(g'(\theta)\bigr)^{2}\bigr).
    \]
    If \( g'(\theta)=0 \) the limiting variance vanishes and one only gets \( \sqrt{n}\bigl(g(W_n)-g(\theta)\bigr)\xrightarrow{\Pbb}0 \); a nondegenerate limit then needs a second-order expansion. This is \citep[Theorem~5.2.9]{hogg2019}.
    \item
    \emph{Several dimensions.} Let \( W_n \) be \(\mathbb{R}^{d}\)-valued with
    \[
    \sqrt{n}\,(W_n-\mu)
    \xrightarrow{(d)}
    \mathcal{N}(0,\Sigma),
    \]
    and let \( g:\mathbb{R}^{d}\to\mathbb{R}^{k} \) be continuously differentiable on a neighbourhood of \( \mu \). Write \( Dg(\mu) \) for the \( k\times d \) Jacobian of \( g \) at \( \mu \). Then
    \[
    \sqrt{n}\,\bigl(g(W_n)-g(\mu)\bigr)
    \xrightarrow{(d)}
    \mathcal{N}\bigl(0,\,Dg(\mu)\,\Sigma\,Dg(\mu)^{T}\bigr).
    \]
    If \( k=1 \), the limiting variance is the scalar \( \nabla g(\mu)^{T}\Sigma\nabla g(\mu) \). This is \citep[Theorem~5.4.6]{hogg2019}. Part~(i) is the case \( d=k=1 \).
\end{enumerate}
\end{theorem}

\begin{proof}
Differentiability of a real function at \( \theta \) is \( g(\theta+h)=g(\theta)+g'(\theta)\,h+r(h) \) with \( r(h)/h\to 0 \) as \( h\to 0 \). Set \( h=W_n-\theta \). Convergence in distribution to a tight limit forces \( W_n\xrightarrow{\Pbb}\theta \), hence \( r(W_n-\theta)=o_{\Pbb}(|W_n-\theta|) \), and
\[
\sqrt{n}\,\bigl(g(W_n)-g(\theta)\bigr)
=
g'(\theta)\cdot\sqrt{n}\,(W_n-\theta)
+
o_{\Pbb}\bigl(\sqrt{n}\,|W_n-\theta|\bigr).
\]
The sequence \( \sqrt{n}\,(W_n-\theta) \) is bounded in probability, so the remainder tends to \( 0 \) in probability. If \( A_n\xrightarrow{(d)}A \) and \( R_n\xrightarrow{\Pbb}0 \), then \( A_n+R_n\xrightarrow{(d)}A \). The display therefore has the same limit as \( g'(\theta)\cdot\sqrt{n}\,(W_n-\theta) \), which is \( \mathcal{N}\bigl(0,\sigma^{2}(g'(\theta))^{2}\bigr) \). This is (i).

For (ii), the same expansion in \( \mathbb{R}^{k} \) reads
\[
g(W_n)-g(\mu)
=
Dg(\mu)\,(W_n-\mu)+o_{\Pbb}(\|W_n-\mu\|).
\]
Multiply by \( \sqrt{n} \). The linear image of \( \mathcal{N}(0,\Sigma) \) under \( Dg(\mu) \) is \( \mathcal{N}\bigl(0,Dg(\mu)\,\Sigma\,Dg(\mu)^{T}\bigr) \) by \cref{def:gaussian-vector}, and the remainder again vanishes in probability. Equivalently, apply Cramér--Wold: for each \( t\in\mathbb{R}^{k} \), the scalar map \( u\mapsto t^{T}g(u) \) falls under (i) with derivative \( t^{T}Dg(\mu) \).
\end{proof}

On the present Gamma model the input \( W_n \) is a scalar function of \( \overline{X} \). Write \( \alpha=2025 \). Then \( \E[X_1]=\alpha\theta \) and \( \Var(X_1)=\alpha\theta^{2} \), so \cref{thm:clt} gives \( \sqrt{n}\,(\overline{X}-\alpha\theta)\xrightarrow{(d)}\mathcal{N}(0,\alpha\theta^{2}) \). The MLE of the scale is \( \hat{\theta}=\overline{X}/\alpha \), and the MLE of \( 1/\theta \) is \( g(\hat{\theta}) \) with \( g(u)=1/u \). The only derivative needed is \( g'(\theta)=-1/\theta^{2} \).

\subsubsection{Solution of Problem 8(a)}

\paragraph{Answer.}
No UMP test of size \( 0.05 \) exists for \( H_0:\theta=1 \) versus \( H_1:\theta\neq 1 \).

\begin{proof}[Correct proof]
Write \( \alpha=2025 \) for the known shape; the letter is not the size, which is \( 0.05 \). On \( \{x_i>0\} \) the joint density is
\[
L(\theta,\mathbf{x})
=
\bigl(\Gamma(\alpha)\bigr)^{-n}\,
\theta^{-n\alpha}
\Bigl(\prod_{i=1}^{n}x_i^{\alpha-1}\Bigr)
\exp\bigl(-S/\theta\bigr),
\qquad
S
\coloneqq
\sum_{i=1}^{n}X_i.
\]
The product of powers of the \( x_i \) does not depend on \( \theta \). For \( 0<\theta_1<\theta_2 \),
\[
\frac{L(\theta_1,\mathbf{x})}{L(\theta_2,\mathbf{x})}
=
\Bigl(\frac{\theta_2}{\theta_1}\Bigr)^{n\alpha}
\exp\Bigl(-S\bigl(\tfrac{1}{\theta_1}-\tfrac{1}{\theta_2}\bigr)\Bigr).
\]
The coefficient \( 1/\theta_1-1/\theta_2 \) is strictly positive, so the ratio is a strictly decreasing function of \( S \). The family therefore has monotone decreasing likelihood ratio in \( S \) (\cref{def:mlr}): large \( S \) favours large \( \theta \).

Each \( X_i \) is \( \mathrm{Gamma}(\alpha,\theta) \) in the scale parameterisation of the exam, so \( S\sim\mathrm{Gamma}(n\alpha,\theta) \). In particular \( S \) is absolutely continuous on \( (0,\infty) \) under every \( \theta>0 \). Neyman--Pearson tests of exact size \( 0.05 \) therefore exist without randomisation, and the \( 0.05 \)- and \( 0.95 \)-quantiles of \( S \) under \( \theta=1 \) are unique.

\emph{Most powerful test against a point \( \theta''>1 \).}
Fix \( \theta''>1 \). Neyman--Pearson \citep[Theorem~8.1.1]{hogg2019} rejects \( \theta=1 \) in favour of \( \theta=\theta'' \) on the set \( \{L(1)/L(\theta'')\le k\} \), with \( k \) chosen to meet the size. The ratio is
\[
\frac{L(1)}{L(\theta'')}
=
(\theta'')^{n\alpha}
\exp\bigl(-S(1-1/\theta'')\bigr).
\]
Here \( 1-1/\theta''>0 \), so \( L(1)/L(\theta'')\le k \) rewrites as \( S\ge c \) for a threshold \( c=c(k,\theta'') \). The size constraint
\[
\Pbb_{\theta=1}(S\ge c)
=
0.05
\]
forces \( c \) to be the \( 0.95 \)-quantile of \( \mathrm{Gamma}(n\alpha,1) \). That quantile does not depend on the particular \( \theta''>1 \). Hence every simple alternative on the right of \( 1 \) has the same most powerful size-\( 0.05 \) critical region
\[
C_+
=
\{S\ge c_+\},
\qquad
\Pbb_{\theta=1}(S\ge c_+)
=
0.05.
\]
Because \( S \) is continuous, \( \Pbb_{\theta=1}(S=c_+)=0 \), so the Neyman--Pearson boundary set is null: any size-\( 0.05 \) most powerful test against any \( \theta''>1 \) coincides with \( C_+ \) up to a \( \Pbb_{\theta=1} \)-null set. (The same region is the Karlin--Rubin test of \cref{thm:karlin-rubin}(i) for the one-sided pair \( H_0:\theta\le 1 \) versus \( H_1:\theta>1 \); only the simple-alternative half is used here.)

\emph{Most powerful test against a point \( \theta''<1 \).}
Now fix \( \theta''\in(0,1) \). The same likelihood ratio has \( 1-1/\theta''<0 \), so \( L(1)/L(\theta'')\le k \) reverses to \( S\le c \). Size \( 0.05 \) forces \( c \) to be the \( 0.05 \)-quantile of \( \mathrm{Gamma}(n\alpha,1) \). Every simple alternative on the left of \( 1 \) therefore has the same most powerful size-\( 0.05 \) region
\[
C_-
=
\{S\le c_-\},
\qquad
\Pbb_{\theta=1}(S\le c_-)
=
0.05,
\]
again unique up to a null set.

\emph{The two regions do not coincide.}
Under \( \theta=1 \), the density of \( S \) is strictly positive on \( (0,\infty) \). The two quantiles therefore satisfy \( 0<c_-<c_+<\infty \), the sets \( C_+ \) and \( C_- \) are disjoint, and
\[
\Pbb_{\theta=1}(C_+\triangle C_-)
=
\Pbb_{\theta=1}(C_+)
+
\Pbb_{\theta=1}(C_-)
=
0.10
>
0.
\]
In particular \( C_+ \) and \( C_- \) are not equal \( \Pbb_{\theta=1} \)-almost surely.

\emph{No UMP test against \( \theta\neq 1 \).}
Suppose some critical region \( C \) were UMP of size \( 0.05 \) for \( H_0:\theta=1 \) versus \( H_1:\theta\neq 1 \), in the sense of \cref{def:ump-size-alpha}. Then \( \Pbb_{\theta=1}(\mathbf{X}\in C)=0.05 \), and the same \( C \) would maximise power at every \( \theta''\neq 1 \). Taking any \( \theta''>1 \), uniqueness of the Neyman--Pearson region forces \( C=C_+ \) almost surely under \( \theta=1 \). Taking any \( \theta''<1 \) forces \( C=C_- \) almost surely. These two requirements cannot hold at once. Therefore no UMP size-\( 0.05 \) test exists.

The equal-tailed two-sided region \( \{S\le a\}\cup\{S\ge b\} \) with each tail of null probability \( 0.025 \) is a legitimate size-\( 0.05 \) test (\cref{def:test-size}). It is not most powerful against any simple alternative \( \theta''\neq 1 \), and it is not UMP. This is Hogg's two-sided normal-mean example \citep[Example~8.2.3]{hogg2019}, with \( S \) in place of \( \overline{X} \).
\end{proof}

\begin{remark}[One-sided contrast]
\label{rem:p8a-onesided}
The same calculation produces UMP tests for the two one-sided problems. Against \( H_1:\theta>1 \) (or \( H_0:\theta\le 1 \) versus \( H_1:\theta>1 \)), reject on \( C_+ \). Against \( H_1:\theta<1 \), reject on \( C_- \). Both statements are \cref{thm:karlin-rubin}; the power function of \( C_+ \) is increasing in \( \theta \), so the size over \( \{\theta\le 1\} \) equals the Type~I error at the boundary \( \theta=1 \).

At \( \theta=1 \) one may mark the cutoff via the chi-square representation \( 2S\sim\chi^{2}_{2n\alpha} \) (scale parameterisation: \( 2X_1/\theta\sim\chi^{2}_{2\alpha} \)). That is a computational convenience for the one-sided tests. It does not manufacture a two-sided UMP.
\end{remark}

\subsubsection{Manuscript proof idea for Problem 8(b)}

The calculation below follows the handwritten notes. \mserr{Red} marks a false claim or a missing justification. The skeleton is the right exam machine: maximise the Gamma likelihood in the scale, pass to \( 1/\hat{\theta} \) by invariance, feed \( \hat{\theta} \) to the CLT, then apply the one-dimensional delta method. The algebra on that skeleton is correct. Two justifications are missing, and one computation is not an input.

\textbf{The MLE of the scale.}
The notes write the joint density in exponential form
\[
L(\theta,\mathbf{x})
=
\exp\Bigl\{
-\frac{1}{\theta}\sum_{i=1}^{n}x_i
+
(\alpha-1)\sum_{i=1}^{n}\log x_i
-
n\log\Gamma(\alpha)
-
n\alpha\log\theta
\Bigr\}
\]
and differentiate the log-likelihood
\[
\frac{\partial}{\partial\theta}\ell(\theta)
=
\frac{1}{\theta^{2}}\sum_{i=1}^{n}x_i
-
\frac{n\alpha}{\theta}.
\]
The critical point is \( \hat{\theta}=\overline{X}/\alpha \). This is correct (and matches the first-order condition already used in part~(a)). A one-line second-derivative or boundary check that the critical point is a maximum is missing; it is cheap and should be written.

\textbf{Passing to \( 1/\theta \).}
The notes record
\[
\Bigl(\widehat{\tfrac{1}{\theta}}\Bigr)_{\mathrm{MLE}}
=
\frac{1}{\hat{\theta}_{\mathrm{MLE}}}
=
\frac{2025}{\overline{X}}.
\]
The identity is the right answer. \mserr{The step \( \widehat{g(\theta)}=g(\hat{\theta}) \) is invariance of the MLE under the bijection \( g(\theta)=1/\theta \), and needs one sentence (\cref{prop:mle-invariance}).} Writing the numerical shape \( 2025 \) in place of \( \alpha \) is fine; keeping the letter until the last display is cleaner.

\textbf{Moments and the CLT.}
The notes compute \( \E[\hat{\theta}]=\theta \) by integrating the Gamma density, then \( \sum X_i\sim\mathrm{Gamma}(n\alpha,\theta) \) and \( \Var(\hat{\theta})=\theta^{2}/(n\alpha) \), and conclude
\[
\sqrt{n}\,(\hat{\theta}-\theta)
\xrightarrow{(d)}
\mathcal{N}\bigl(0,\theta^{2}/\alpha\bigr).
\]
The variance and the Gaussian limit are correct. \mserr{Unbiasedness of \( \hat{\theta} \) is not an input to the CLT, and is not asked.} The CLT uses \( \E[X_1]=\alpha\theta \) and \( \Var(X_1)=\alpha\theta^{2} \); both follow from the Gamma integrals (or from \( S\sim\mathrm{Gamma}(n\alpha,\theta) \)) without a separate unbiasedness paragraph. The estimator of \( 1/\theta \) is in any case biased, by Jensen.

\textbf{Delta method.}
The notes set \( g(u)=1/u \), record \( g'(u)=-1/u^{2} \), and multiply:
\[
\bigl(g'(\theta)\bigr)^{2}\cdot\frac{\theta^{2}}{\alpha}
=
\frac{1}{\theta^{4}}\cdot\frac{\theta^{2}}{\alpha}
=
\frac{1}{\alpha\theta^{2}}.
\]
This is the right arithmetic for \cref{thm:delta-method}(i). \mserr{Name the theorem and the hypothesis: \( g \) is \( C^{1} \) on \( (0,\infty) \), and \( \hat{\theta}\xrightarrow{\Pbb}\theta \) is already implied by the CLT.} Do not derive the exact Inverse-Gamma law of \( 1/\overline{X} \).

\errpanel{%
\textbf{\mserr{The gap, in one box.}}

The notes have the correct MLE \( \alpha/\overline{X} \) and the correct limit \( \mathcal{N}(0,1/(\alpha\theta^{2})) \).

What is missing is the one-sentence invariance \( \widehat{1/\theta}=1/\hat{\theta} \), and the name of the delta method with \( g(u)=1/u \). Computing \( \E[\hat{\theta}]=\theta \) by hand is not needed and does not justify the CLT.
}

\begin{proposition}[Invariance of the MLE]
\label{prop:mle-invariance}
Let \( L(\theta,\mathbf{x}) \) be a likelihood on a parameter set \( \Theta \), and let \( \hat{\theta}(\mathbf{x}) \) be a maximiser. If \( g:\Theta\to\Lambda \) is a bijection, then \( g(\hat{\theta}(\mathbf{x})) \) maximises the reparameterised likelihood \( \lambda\mapsto L(g^{-1}(\lambda),\mathbf{x}) \), hence is an MLE of \( g(\theta) \). This is \citep[Theorem~6.1.2]{hogg2019}. The map need not be linear; here \( g(\theta)=1/\theta \) on \( (0,\infty) \).
\end{proposition}

\subsubsection{Solution of Problem 8(b)}

\paragraph{Answer.}
Write \( \alpha=2025 \). The MLE of the scale is \( \hat{\theta}=\overline{X}/\alpha \). The MLE of \( 1/\theta \) is
\[
\widehat{\tfrac{1}{\theta}}
=
\frac{\alpha}{\overline{X}}
=
\frac{2025}{\overline{X}}.
\]
Its asymptotic distribution is
\[
\sqrt{n}\Bigl(\frac{\alpha}{\overline{X}}-\frac{1}{\theta}\Bigr)
\xrightarrow{(d)}
\mathcal{N}\Bigl(0,\,\frac{1}{\alpha\theta^{2}}\Bigr)
=
\mathcal{N}\Bigl(0,\,\frac{1}{2025\,\theta^{2}}\Bigr).
\]

\begin{proof}[Correct proof]
Write \( \alpha=2025 \) for the known shape, and write \( S=\sum_{i=1}^{n}X_i \), \( \overline{X}=S/n \). On \( \{x_i>0\} \) the log-likelihood is
\[
\ell(\theta)
=
-\frac{S}{\theta}
-
n\alpha\log\theta
-
n\log\Gamma(\alpha)
+
(\alpha-1)\sum_{i=1}^{n}\log x_i.
\]
Differentiating in the scale,
\[
\ell'(\theta)
=
\frac{S}{\theta^{2}}
-
\frac{n\alpha}{\theta},
\qquad
\ell''(\theta)
=
-\frac{2S}{\theta^{3}}
+
\frac{n\alpha}{\theta^{2}}.
\]
The unique critical point in \( (0,\infty) \) is \( \hat{\theta}=S/(n\alpha)=\overline{X}/\alpha \). At that point \( S=n\alpha\hat{\theta} \), so \( \ell''(\hat{\theta})=-n\alpha/\hat{\theta}^{2}<0 \). Alternatively, \( \ell(\theta)\to-\infty \) as \( \theta\to 0^{+} \) and as \( \theta\to\infty \). Thus \( \hat{\theta} \) is the MLE of \( \theta \).

The map \( g(\theta)=1/\theta \) is a bijection of \( (0,\infty) \). \Cref{prop:mle-invariance} therefore yields that the MLE of \( 1/\theta \) is
\[
g(\hat{\theta})
=
\frac{1}{\hat{\theta}}
=
\frac{\alpha}{\overline{X}}.
\]

\emph{Moments.}
The exam density is the scale-Gamma law. One has
\[
\E[X_1]
=
\int_{0}^{\infty}x\cdot\frac{1}{\Gamma(\alpha)\theta^{\alpha}}x^{\alpha-1}e^{-x/\theta}\,\mathrm{d}x
=
\frac{\Gamma(\alpha+1)}{\Gamma(\alpha)}\,\theta
=
\alpha\theta,
\]
and likewise
\[
\E[X_1^{2}]
=
\frac{\Gamma(\alpha+2)}{\Gamma(\alpha)}\,\theta^{2}
=
\alpha(\alpha+1)\theta^{2},
\]
hence \( \Var(X_1)=\alpha\theta^{2} \). (Equivalently: \( S\sim\mathrm{Gamma}(n\alpha,\theta) \), so \( \E[S]=n\alpha\theta \) and \( \Var(S)=n\alpha\theta^{2} \).) In particular \( \E[\hat{\theta}]=\theta \) and \( \Var(\hat{\theta})=\theta^{2}/(n\alpha) \), but those two identities are not used below except as a check on the CLT variance.

\emph{CLT for the scale.}
\Cref{thm:clt} applied to the i.i.d.\ sample \( X_1,\dots,X_n \) gives
\[
\sqrt{n}\,(\overline{X}-\alpha\theta)
\xrightarrow{(d)}
\mathcal{N}(0,\alpha\theta^{2}).
\]
Dividing by the constant \( \alpha \) yields
\[
\sqrt{n}\,(\hat{\theta}-\theta)
=
\frac{1}{\alpha}\,\sqrt{n}\,(\overline{X}-\alpha\theta)
\xrightarrow{(d)}
\mathcal{N}\bigl(0,\,\theta^{2}/\alpha\bigr).
\]

\emph{Delta method.}
The map \( g(u)=1/u \) is \( C^{1} \) on \( (0,\infty) \), with \( g'(u)=-1/u^{2} \) and in particular \( g'(\theta)=-1/\theta^{2}\neq 0 \). Convergence in distribution to a tight limit forces \( \hat{\theta}\xrightarrow{\Pbb}\theta \), so one is eventually inside the domain of \( g \). \Cref{thm:delta-method}(i) therefore yields
\[
\sqrt{n}\bigl(g(\hat{\theta})-g(\theta)\bigr)
\xrightarrow{(d)}
\mathcal{N}\Bigl(0,\,\frac{\theta^{2}}{\alpha}\bigl(g'(\theta)\bigr)^{2}\Bigr)
=
\mathcal{N}\Bigl(0,\,\frac{\theta^{2}}{\alpha}\cdot\frac{1}{\theta^{4}}\Bigr)
=
\mathcal{N}\Bigl(0,\,\frac{1}{\alpha\theta^{2}}\Bigr).
\]
Substituting \( g(\hat{\theta})=\alpha/\overline{X} \) and \( \alpha=2025 \) is the displayed answer.
\end{proof}

\begin{remark}[What (b) is not]
\label{rem:p8b-exact}
The exact law is available and is not asked. Since \( S\sim\mathrm{Gamma}(n\alpha,\theta) \), one has \( 1/\hat{\theta}=n\alpha/S \), an Inverse-Gamma. The exam asks only for the \( \sqrt{n} \)-limit.

The same Gaussian limit for \( \hat{\theta} \) is the regular MLE expansion: the per-observation Fisher information is \( I(\theta)=\alpha/\theta^{2} \), so \( \sqrt{n}(\hat{\theta}-\theta)\xrightarrow{(d)}\mathcal{N}(0,1/I(\theta)) \). That is a check, not a replacement for the CLT-plus-delta write-up.

Jensen's inequality for the convex map \( u\mapsto 1/u \) gives \( \E[1/\hat{\theta}]>1/\theta \) for every finite \( n \). Bias of \( 1/\hat{\theta} \) is compatible with the \( \sqrt{n} \)-central limit theorem.
\end{remark}

\subsection{Problem 9. [15 pts.]}
Let \( (X_1, Y_1), \dots, (X_n, Y_n) \) be iid random bivariate vectors. Let \( \rho \) be the correlation coefficient between \( X_i \) and \( Y_i \). Let

\[
\hat{\rho} = \frac{1}{(n-1)\sqrt{S_X^2 S_Y^2}} \sum_{i=1}^n (X_i - \overline{X})(Y_i - \overline{Y}),
\]

where \( \overline{X} = \frac{\sum_{i=1}^n X_i}{n} \), \( \overline{Y} = \frac{\sum_{i=1}^n Y_i}{n} \), \( S_X^2 = \frac{\sum_{i=1}^n (X_i - \overline{X})^2}{n-1} \), \( S_Y^2 = \frac{\sum_{i=1}^n (Y_i - \overline{Y})^2}{n-1} \).
\begin{enumerate}[label=(\alph*)]
    \item Assume that \( \mathbb{E}|X_i|^4 < \infty \) and \( \mathbb{E}|Y_i|^4 < \infty \), prove that \( \sqrt{n}(\hat{\rho} - \rho) \longrightarrow N(0, c) \) for some value \( c \). Please identify \( c \) in terms of moments of \( (X_i, Y_i) \), however, there is no need to simplify your calculation on \( c \).
    \item Further, assume that \( (X_i, Y_i) \) follows a bivariate normal distribution that

    \[
    \begin{pmatrix} X_i \\ Y_i \end{pmatrix}
    \sim
    N\left(
    \begin{pmatrix} 0 \\ 0 \end{pmatrix},
    \begin{pmatrix} 1 & \rho \\ \rho & 1 \end{pmatrix}
    \right).
    \]

    The probability density function of the above bivariate normal distribution is given by

    \[
    f(x, y \mid \rho) = \frac{1}{2\pi \sqrt{1 - \rho^2}} \exp\left\{ -\frac{1}{2(1 - \rho^2)} [x^2 - 2\rho xy + y^2] \right\}.
    \]

    Find a minimal sufficient statistic for \( \rho \). Is this minimal sufficient statistic complete or not? Please explain.
\end{enumerate}

\setcounter{section}{9}

\subsubsection{Definitions used in Problem 9}

Part~(a) is not a one-dimensional average. The sample correlation is a smooth function of five empirical moments, so the machine is the multivariate CLT on that moment vector, followed by the delta method. The one-dimensional CLT is the input to Cramér--Wold; it is not by itself the argument for \( \sqrt{n}(\hat{\rho}-\rho) \). The fourth-moment hypotheses \( \E|X_i|^{4}<\infty \) and \( \E|Y_i|^{4}<\infty \) are exactly the second-moment hypotheses for the vector \( (X,Y,X^{2},Y^{2},XY) \).

\begin{definition}[Convergence in distribution]
\label{def:conv-d}
A sequence of \(\mathbb{R}^{d}\)-valued random vectors \( Z_n \) \emph{converges in distribution} to \( Z \), written \( Z_n\xrightarrow{(d)}Z \) or \( Z_n\Rightarrow Z \), if \( \E[g(Z_n)]\to\E[g(Z)] \) for every bounded continuous \( g:\mathbb{R}^{d}\to\mathbb{R} \). Equivalently, \( \Pbb(Z_n\in A)\to\Pbb(Z\in A) \) for every Borel set \( A \) with \( \Pbb(Z\in\partial A)=0 \). In dimension one this is \( F_{Z_n}(x)\to F_Z(x) \) at every continuity point of \( F_Z \). See \citep[Theorems~3.2.3 and~3.2.5]{durrett2010}. Convergence in probability (\cref{def:conv-p}) implies convergence in distribution; the converse holds when the limit is a constant.
\end{definition}

\begin{theorem}[Central Limit Theorem]
\label{thm:clt}
Let \( \xi_1,\xi_2,\dots \) be i.i.d.\ real random variables with \( \E[\xi_1]=\mu\in\mathbb{R} \) and \( 0<\Var(\xi_1)=\sigma^{2}<\infty \). Write \( S_n=\sum_{i=1}^{n}\xi_i \) and \( \overline{\xi}_n=S_n/n \). Then
\[
\frac{S_n-n\mu}{\sigma\sqrt{n}}
=
\frac{\sqrt{n}\,(\overline{\xi}_n-\mu)}{\sigma}
\xrightarrow{(d)}
\mathcal{N}(0,1)
\]
in the sense of \cref{def:conv-d}. Equivalently, \( \sqrt{n}\,(\overline{\xi}_n-\mu)\xrightarrow{(d)}\mathcal{N}(0,\sigma^{2}) \). This is \citep[Theorem~5.3.1]{hogg2019} and \citep[Theorem~3.4.1]{durrett2010}.
\end{theorem}

\begin{proof}
Centre \( \eta_i=\xi_i-\mu \), so \( \E[\eta_1]=0 \) and \( \E[\eta_1^{2}]=\sigma^{2} \). Let \( \varphi(u)=\E[e^{iu\eta_1}] \). Finite second moment gives \( \varphi(u)=1-\sigma^{2}u^{2}/2+o(u^{2}) \) as \( u\to 0 \). The characteristic function of \( S_n^{\eta}/(\sigma\sqrt{n}) \) is
\[
\Bigl(\varphi\Bigl(\frac{t}{\sigma\sqrt{n}}\Bigr)\Bigr)^{n}
=
\Bigl(1-\frac{t^{2}}{2n}+o(n^{-1})\Bigr)^{n}
\to
e^{-t^{2}/2}.
\]
\Cref{prop:levy-gaussian} yields the \( \mathcal{N}(0,1) \) limit. (Hogg's write-up in \S5.3 uses the moment generating function under an extra existence assumption; replacing the mgf by \( \varphi \) removes that assumption and is the Durrett argument.)
\end{proof}

\begin{theorem}[Cramér--Wold]
\label{thm:cramer-wold}
Let \( Z_n,Z \) be \(\mathbb{R}^{d}\)-valued. Then \( Z_n\xrightarrow{(d)}Z \) if and only if \( t^{T}Z_n\xrightarrow{(d)}t^{T}Z \) for every \( t\in\mathbb{R}^{d} \). See \citep[\S3.9]{durrett2010}. In particular, \( Z_n\xrightarrow{(d)}\mathcal{N}(\mu,\Sigma) \) if and only if, for every \( t \),
\[
t^{T}Z_n
\xrightarrow{(d)}
\mathcal{N}(t^{T}\mu,\,t^{T}\Sigma t),
\]
which is the one-dimensional statement that every linear image is Gaussian with the correct mean and variance (\cref{def:gaussian-vector}).
\end{theorem}

\begin{theorem}[Multivariate Central Limit Theorem]
\label{thm:clt-mv}
Let \( V_1,V_2,\dots \) be i.i.d.\ random vectors in \( \mathbb{R}^{d} \) with \( \E[V_1]=\mu \) and \( \operatorname{Cov}(V_1)=\Sigma \) finite (positive semidefinite, not necessarily invertible). Write \( \overline{V}_n=n^{-1}\sum_{i=1}^{n}V_i \). Then
\[
\sqrt{n}\,(\overline{V}_n-\mu)
\xrightarrow{(d)}
\mathcal{N}(0,\Sigma).
\]
This is \citep[Theorem~5.4.4]{hogg2019}, stated without an mgf hypothesis and without requiring \( \Sigma \) to be definite. The exam's \( \longrightarrow N(0,c) \) is the \( d=1 \) special case of this limit after a scalar map.
\end{theorem}

\begin{proof}
Fix \( t\in\mathbb{R}^{d} \). The scalars \( t^{T}V_i \) are i.i.d.\ with mean \( t^{T}\mu \) and variance \( t^{T}\Sigma t \). If \( t^{T}\Sigma t=0 \) then \( t^{T}(\overline{V}_n-\mu)=0 \) almost surely, which is \( \mathcal{N}(0,0) \). If \( t^{T}\Sigma t>0 \), \cref{thm:clt} gives
\[
\sqrt{n}\,t^{T}(\overline{V}_n-\mu)
\xrightarrow{(d)}
\mathcal{N}(0,t^{T}\Sigma t).
\]
\Cref{thm:cramer-wold} upgrades every one-dimensional projection to the Gaussian vector \( \mathcal{N}(0,\Sigma) \).
\end{proof}

The several-dimensional delta method used below is \cref{thm:delta-method}(ii).

\begin{theorem}[Slutsky]
\label{thm:slutsky}
If \( A_n\xrightarrow{(d)}A \) in \( \mathbb{R}^{k} \) and \( B_n\xrightarrow{\Pbb}b \) in \( \mathbb{R}^{m} \), then \( (A_n,B_n)\xrightarrow{(d)}(A,b) \). Consequently, if \( g \) is continuous at the support of \( (A,b) \), one has \( g(A_n,B_n)\xrightarrow{(d)}g(A,b) \). The usual calculus forms are: \( A_n+B_n\xrightarrow{(d)}A+b \), \( B_n A_n\xrightarrow{(d)}bA \), and, when \( b\neq 0 \) is a scalar, \( A_n/B_n\xrightarrow{(d)}A/b \). See \citep[Theorem~5.2.5]{hogg2019}.
\end{theorem}

The exam map for part~(a) is then: put
\[
V_i
=
\bigl(X_i,\,Y_i,\,X_i^{2},\,Y_i^{2},\,X_i Y_i\bigr)^{T},
\]
apply \cref{thm:clt-mv} (the fourth-moment assumptions make \( \E\|V_1\|^{2}<\infty \)), rewrite \( \hat{\rho} \) as \( g(\overline{V}_n) \), and \emph{substitute} \( \nabla g(\mu) \) into \( \Sigma \): the exam wants the expanded quadratic form \( c=\sum_{j,k}\gamma_j\gamma_k\Sigma_{jk} \) with every \( \gamma_j \) and every \( \Sigma_{jk} \) written as a moment of \( (X,Y) \). Do not stop at the matrix slogan \( \nabla g^{T}\Sigma\nabla g \). Do not cancel the expression down to \( (1-\rho^{2})^{2} \): that identity uses Gaussian moment relations and is not available in (a). The factor \( n-1 \) in the exam formula cancels identically against the two sample variances; one does not need \cref{thm:slutsky} for that cancellation.

\subsubsection{Hint for Problem 9(b)}

The model in (b) is not the five-parameter bivariate normal. Means are known to be \( 0 \) and variances known to be \( 1 \); the only unknown is \( \rho\in(-1,1) \). The joint density depends on the sample only through
\[
T
=
\Bigl(\sum_{i=1}^{n}(X_i^{2}+Y_i^{2}),\,\sum_{i=1}^{n}X_i Y_i\Bigr).
\]
That pair is sufficient by factorization and minimal by the likelihood-ratio criterion. It is \emph{not} complete: the family of laws of \( T \) is a curved exponential family (two-dimensional sufficient statistic, one-dimensional parameter), and an explicit nonzero unbiased estimator of zero is \( T_1-2n \).

Three things not to do: (i) list \( \bigl(\sum X_i,\sum Y_i,\sum X_i^{2},\sum Y_i^{2},\sum X_i Y_i\bigr) \) and call it minimal sufficient --- those five statistics are jointly complete sufficient for the \emph{five}-parameter bivariate normal \citep[Exercise~7.5.8]{hogg2019}, which is a different model; (ii) quote that exercise and conclude that \( T \) is complete on the present submodel; (iii) claim that \( \hat{\rho} \) itself is minimal sufficient. Completeness is \cref{def:complete}; minimal sufficiency is \cref{def:minimal-sufficient,thm:lehmann-scheffe-ratio}.

\begin{definition}[Sufficient statistic]
\label{def:sufficient}
A statistic \( T=T(X) \) is \emph{sufficient} for \( \theta \) if the conditional law of the sample \( X \) given \( T \) does not depend on \( \theta \). See \citep[\S7.2]{hogg2019}.
\end{definition}

\begin{theorem}[Neyman factorization]
\label{thm:factorization}
\( T \) is sufficient for \( \theta \) if and only if the joint density (or mass function) factors as
\[
f(x\mid\theta)
=
k_1\bigl(T(x),\theta\bigr)\,k_2(x),
\]
with \( k_2 \) free of \( \theta \). This is \citep[Theorem~7.2.1]{hogg2019}.
\end{theorem}

\begin{definition}[Minimal sufficient statistic]
\label{def:minimal-sufficient}
A sufficient statistic \( T \) is \emph{minimal sufficient} if, for every other sufficient statistic \( T' \), there is a function \( h \) with \( T=h(T') \) almost surely. Equivalently: \( T \) is a coarsest sufficient reduction of the data. See \citep[\S7.3]{hogg2019}.
\end{definition}

\begin{theorem}[Lehmann--Scheff\'e ratio criterion]
\label{thm:lehmann-scheffe-ratio}
Assume the model has a joint density \( f(x\mid\theta) \) with respect to a \( \theta \)-free dominating measure. A statistic \( T \) is minimal sufficient if and only if, for almost every pair of sample points \( x,y \), the ratio \( f(x\mid\theta)/f(y\mid\theta) \) is independent of \( \theta \) precisely when \( T(x)=T(y) \). This is \citep[Theorem~7.3.1]{hogg2019}.
\end{theorem}

\begin{definition}[Completeness]
\label{def:complete}
A family \( \{f_T(t;\theta):\theta\in\Omega\} \) is \emph{complete} if \( \E_\theta[u(T)]=0 \) for every \( \theta\in\Omega \) forces \( u(T)=0 \) almost surely under every member of the family. See \citep[\S7.4]{hogg2019}. Sufficiency compresses the data; completeness forbids extra unbiased-zero functions of that compression.
\end{definition}

\begin{theorem}[Regular exponential family, completeness]
\label{thm:expfam-complete}
If the observations are i.i.d.\ from a regular \( k \)-parameter exponential family
\[
f(x;\theta)
=
\exp\bigl(p(\theta)^{T}K(x)+q(\theta)+S(x)\bigr)
\]
whose natural parameter space contains a nonempty open set in \( \mathbb{R}^{k} \), and whose support does not depend on \( \theta \), then \( T=\sum_{i=1}^{n}K(X_i) \) is a complete sufficient statistic. This is \citep[Theorems~7.5.1--7.5.2]{hogg2019} in the one-parameter case, and the same argument in \( k \) parameters. If \( p(\theta) \) traces only a lower-dimensional curve in \( \mathbb{R}^{k} \), the hypothesis fails and completeness may fail with it.
\end{theorem}

\subsubsection{Solution of Problem 9}

\paragraph{Answer.}
Write \( V=(X,Y,X^{2},Y^{2},XY)^{T} \) for a generic pair, \( \mu=\E[V] \), and \( \Sigma=\operatorname{Cov}(V) \). Let
\[
g(u,v,a,b,c)
=
\frac{c-uv}{\sqrt{(a-u^{2})(b-v^{2})}}
\]
on \( \{a>u^{2},\,b>v^{2}\} \). Then \( \hat{\rho}=g(\overline{V}_n) \) and \( \rho=g(\mu) \), and
\[
\sqrt{n}\,(\hat{\rho}-\rho)
\xrightarrow{(d)}
\mathcal{N}(0,c)
\]
with the expanded moment formula for \( c \) recorded in the proof: \( c=\sum_{j,k=1}^{5}\gamma_j\gamma_k\Sigma_{jk} \), each \( \gamma_j=\partial_j g(\mu) \) and each \( \Sigma_{jk}=\operatorname{Cov}(V^{(j)},V^{(k)}) \) written out. For the centred unit-variance bivariate normal, the pair
\[
T
=
\Bigl(\sum_{i=1}^{n}(X_i^{2}+Y_i^{2}),\,\sum_{i=1}^{n}X_i Y_i\Bigr)
\]
is minimal sufficient for \( \rho \) and is not complete.

\begin{proof}[Correct proof]
\textbf{(a), the moment vector.}
Write a generic pair as \( (X,Y) \) and set
\[
V
\coloneqq
\bigl(X,\,Y,\,X^{2},\,Y^{2},\,XY\bigr)^{T}
\in
\mathbb{R}^{5}.
\]
The hypothesis \( \E[X^{4}]<\infty \) and \( \E[Y^{4}]<\infty \) must be upgraded to \( \E\|V\|^{2}<\infty \), which is the second-moment hypothesis of \cref{thm:clt-mv}. Lyapunov gives \( \E[X^{2}]\le(\E[X^{4}])^{1/2}<\infty \) and \( \E[Y^{2}]<\infty \). Cauchy--Schwarz on \( X^{2} \) and \( Y^{2} \) gives
\[
\E[X^{2}Y^{2}]
\le
\bigl(\E[X^{4}]\bigr)^{1/2}\bigl(\E[Y^{4}]\bigr)^{1/2}
<
\infty.
\]
The mixed moments of order four that appear in \( \E[VV^{T}] \) are then finite as well: \( \E[|X^{3}Y|]=\E[|X|^{2}\cdot|XY|]\le(\E[X^{4}])^{1/2}(\E[X^{2}Y^{2}])^{1/2}<\infty \), and likewise \( \E[|XY^{3}|]<\infty \). Thus \( \E\|V\|^{2}<\infty \). Write
\[
\mu
\coloneqq
\E[V]
=
\bigl(\mu_X,\,\mu_Y,\,\E[X^{2}],\,\E[Y^{2}],\,\E[XY]\bigr)^{T},
\qquad
\Sigma
\coloneqq
\operatorname{Cov}(V)
=
\E[VV^{T}]-\mu\mu^{T}.
\]
The entries of \( \Sigma \) are moments of \( (X,Y) \) of order at most four; the exam does not ask that they be expanded. Assume throughout that \( \sigma_X^{2}=\Var(X)>0 \) and \( \sigma_Y^{2}=\Var(Y)>0 \), so that \( \rho \) is defined and the map below is differentiable at \( \mu \). (If either variance vanishes the correlation coefficient is not defined.)

\textbf{(a), algebraic identity.}
The exam formula uses the unbiased sample variances, but the factor \( n-1 \) cancels. Write \( A_n=\sum_{i=1}^{n}(X_i-\overline{X})^{2} \) and \( B_n=\sum_{i=1}^{n}(Y_i-\overline{Y})^{2} \) and \( C_n=\sum_{i=1}^{n}(X_i-\overline{X})(Y_i-\overline{Y}) \). Then \( S_X^{2}=A_n/(n-1) \), \( S_Y^{2}=B_n/(n-1) \), and
\[
\hat{\rho}
=
\frac{C_n}{(n-1)\sqrt{S_X^{2}S_Y^{2}}}
=
\frac{C_n}{\sqrt{A_n B_n}}.
\]
The usual expansion \( C_n=n\bigl(n^{-1}\sum X_i Y_i-\overline{X}\,\overline{Y}\bigr) \), and likewise for \( A_n \) and \( B_n \), therefore yields
\[
\hat{\rho}
=
g(\overline{V}_n),
\qquad
g(u,v,a,b,c)
\coloneqq
\frac{c-uv}{\sqrt{(a-u^{2})(b-v^{2})}},
\]
on the set where \( a>u^{2} \) and \( b>v^{2} \). The same formula at the population point is the definition of \( \rho \):
\[
g(\mu)
=
\frac{\E[XY]-\mu_X\mu_Y}{\sqrt{(\E[X^{2}]-\mu_X^{2})(\E[Y^{2}]-\mu_Y^{2})}}
=
\rho.
\]
The event \( \{A_n=0\}\cup\{B_n=0\} \) has probability tending to zero (by the WLLN for the second-moment coordinates, or because \( \overline{V}_n\xrightarrow{\Pbb}\mu \) from the CLT below), so \( g(\overline{V}_n) \) is eventually well-defined and one may redefine \( \hat{\rho} \) arbitrarily on that null-probability set.

\textbf{(a), multivariate CLT.}
The vectors \( V_1,\dots,V_n \) are i.i.d.\ in \( \mathbb{R}^{5} \) with mean \( \mu \) and finite covariance \( \Sigma \). \Cref{thm:clt-mv} gives
\[
\sqrt{n}\,(\overline{V}_n-\mu)
\xrightarrow{(d)}
\mathcal{N}(0,\Sigma).
\]
(The one-dimensional CLT of \cref{thm:clt} is the input to this statement, via Cram\'er--Wold: every linear image \( t^{T}V_i \) is a real i.i.d.\ sequence with finite variance \( t^{T}\Sigma t \). It is not by itself an argument for \( \sqrt{n}(\hat{\rho}-\rho) \).)

\textbf{(a), delta method.}
The map \( g \) is \( C^{1} \) on the open set \( \{a>u^{2},\,b>v^{2}\} \), which contains \( \mu \). Write \( N=c-uv \), \( A=a-u^{2} \), \( B=b-v^{2} \), and \( D=\sqrt{AB} \), so \( g=N/D \). Differentiating,
\begin{align*}
\partial_u g
&=
\frac{B(uN-vA)}{D^{3}},
&
\partial_v g
&=
\frac{A(vN-uB)}{D^{3}},
\\
\partial_a g
&=
-\frac{NB}{2D^{3}},
&
\partial_b g
&=
-\frac{NA}{2D^{3}},
&
\partial_c g
&=
\frac{1}{D}.
\end{align*}
Evaluating at \( \mu \) one has \( N=\operatorname{Cov}(X,Y) \), \( A=\sigma_X^{2} \), \( B=\sigma_Y^{2} \), \( D=\sigma_X\sigma_Y \). Write \( \gamma=\nabla g(\mu) \) for the column of these five numbers:
\begin{align*}
\gamma_1
&=
\frac{\sigma_Y^{2}\bigl(\mu_X\operatorname{Cov}(X,Y)-\mu_Y\sigma_X^{2}\bigr)}{(\sigma_X\sigma_Y)^{3}},
&
\gamma_2
&=
\frac{\sigma_X^{2}\bigl(\mu_Y\operatorname{Cov}(X,Y)-\mu_X\sigma_Y^{2}\bigr)}{(\sigma_X\sigma_Y)^{3}},
\\
\gamma_3
&=
-\frac{\operatorname{Cov}(X,Y)\,\sigma_Y^{2}}{2(\sigma_X\sigma_Y)^{3}},
&
\gamma_4
&=
-\frac{\operatorname{Cov}(X,Y)\,\sigma_X^{2}}{2(\sigma_X\sigma_Y)^{3}},
&
\gamma_5
&=
\frac{1}{\sigma_X\sigma_Y}.
\end{align*}
Each of \( \sigma_X^{2}=\E[X^{2}]-\mu_X^{2} \), \( \sigma_Y^{2}=\E[Y^{2}]-\mu_Y^{2} \), and \( \operatorname{Cov}(X,Y)=\E[XY]-\mu_X\mu_Y \) is already a moment. The covariance matrix of \( V \) is the \( 5\times 5 \) array of moments of order at most four. Label the coordinates of \( V \) as \( (V^{(1)},\dots,V^{(5)})=(X,Y,X^{2},Y^{2},XY) \). Then \( \Sigma_{jk}=\operatorname{Cov}(V^{(j)},V^{(k)})=\E[V^{(j)}V^{(k)}]-\mu_j\mu_k \), namely
\begin{align*}
\Sigma_{11}
&=
\E[X^{2}]-\mu_X^{2},
&
\Sigma_{22}
&=
\E[Y^{2}]-\mu_Y^{2},
&
\Sigma_{12}
&=
\E[XY]-\mu_X\mu_Y,
\\
\Sigma_{13}
&=
\E[X^{3}]-\mu_X\E[X^{2}],
&
\Sigma_{14}
&=
\E[XY^{2}]-\mu_X\E[Y^{2}],
&
\Sigma_{15}
&=
\E[X^{2}Y]-\mu_X\E[XY],
\\
\Sigma_{23}
&=
\E[X^{2}Y]-\mu_Y\E[X^{2}],
&
\Sigma_{24}
&=
\E[Y^{3}]-\mu_Y\E[Y^{2}],
&
\Sigma_{25}
&=
\E[XY^{2}]-\mu_Y\E[XY],
\\
\Sigma_{33}
&=
\E[X^{4}]-(\E[X^{2}])^{2},
&
\Sigma_{44}
&=
\E[Y^{4}]-(\E[Y^{2}])^{2},
&
\Sigma_{55}
&=
\E[X^{2}Y^{2}]-(\E[XY])^{2},
\\
\Sigma_{34}
&=
\E[X^{2}Y^{2}]-\E[X^{2}]\E[Y^{2}],
&
\Sigma_{35}
&=
\E[X^{3}Y]-\E[X^{2}]\E[XY],
&
\Sigma_{45}
&=
\E[XY^{3}]-\E[Y^{2}]\E[XY].
\end{align*}
\Cref{thm:delta-method}(ii) --- or the first-order expansion \( g(\overline{V}_n)-g(\mu)=\gamma^{T}(\overline{V}_n-\mu)+o_{\Pbb}(\|\overline{V}_n-\mu\|) \), multiplied by \( \sqrt{n} \) --- yields \( \sqrt{n}(\hat{\rho}-\rho)\xrightarrow{(d)}\mathcal{N}(0,c) \) with the quadratic form substituted, not left as a matrix slogan:
\[
c
=
\sum_{j=1}^{5}\sum_{k=1}^{5}
\gamma_j\,\Sigma_{jk}\,\gamma_k
=
\sum_{j=1}^{5}\gamma_j^{2}\,\Sigma_{jj}
+
2\sum_{1\le j<k\le 5}
\gamma_j\gamma_k\,\Sigma_{jk}.
\]
Writing the fifteen terms out,
\begin{align*}
c
&=
\gamma_1^{2}\Sigma_{11}
+
\gamma_2^{2}\Sigma_{22}
+
\gamma_3^{2}\Sigma_{33}
+
\gamma_4^{2}\Sigma_{44}
+
\gamma_5^{2}\Sigma_{55}
\\
&\qquad
+
2\gamma_1\gamma_2\Sigma_{12}
+
2\gamma_1\gamma_3\Sigma_{13}
+
2\gamma_1\gamma_4\Sigma_{14}
+
2\gamma_1\gamma_5\Sigma_{15}
\\
&\qquad
+
2\gamma_2\gamma_3\Sigma_{23}
+
2\gamma_2\gamma_4\Sigma_{24}
+
2\gamma_2\gamma_5\Sigma_{25}
\\
&\qquad
+
2\gamma_3\gamma_4\Sigma_{34}
+
2\gamma_3\gamma_5\Sigma_{35}
+
2\gamma_4\gamma_5\Sigma_{45}.
\end{align*}
Every factor on the right is a moment of \( (X,Y) \) of order at most four, via the displays for \( \gamma \) and \( \Sigma \) above. That is the identification the exam asks for. One is not asked to cancel this expression; in particular one must not replace it by the Gaussian special value \( (1-\rho^{2})^{2} \).

\textbf{(b), sufficiency.}
The joint density of the sample is
\[
L(\rho)
=
(2\pi)^{-n}(1-\rho^{2})^{-n/2}
\exp\biggl(
-\frac{1}{2(1-\rho^{2})}
\sum_{i=1}^{n}\bigl(X_i^{2}-2\rho X_i Y_i+Y_i^{2}\bigr)
\biggr)
\]
\[
=
(2\pi)^{-n}(1-\rho^{2})^{-n/2}
\exp\biggl(
-\frac{T_1}{2(1-\rho^{2})}
+
\frac{\rho\,T_2}{1-\rho^{2}}
\biggr),
\]
where
\[
T_1
\coloneqq
\sum_{i=1}^{n}(X_i^{2}+Y_i^{2}),
\qquad
T_2
\coloneqq
\sum_{i=1}^{n}X_i Y_i,
\qquad
T
\coloneqq
(T_1,T_2).
\]
The factor \( (2\pi)^{-n} \) is free of \( \rho \), and the remaining factor depends on the data only through \( T \). \Cref{thm:factorization} yields that \( T \) is sufficient for \( \rho \).

\textbf{(b), minimal sufficiency.}
Let \( (x,y) \) and \( (x',y') \) be two sample points, with associated \( (T_1,T_2) \) and \( (T_1',T_2') \). The likelihood ratio is
\[
\frac{L(\rho;x,y)}{L(\rho;x',y')}
=
\exp\Bigl\{
-\frac{T_1-T_1'}{2(1-\rho^{2})}
+
\frac{\rho(T_2-T_2')}{1-\rho^{2}}
\Bigr\}.
\]
Write \( a=T_1-T_1' \) and \( b=T_2-T_2' \). The exponent is \( \bigl(-a/2+\rho b\bigr)/(1-\rho^{2}) \). This function of \( \rho\in(-1,1) \) is constant if and only if \( a=b=0 \). Indeed, if it equals a constant \( k \), then \( -a/2+\rho b=k(1-\rho^{2}) \) as polynomials in \( \rho \), so comparing coefficients of \( \{1,\rho,\rho^{2}\} \) forces \( k=0 \), \( b=0 \), and \( a=0 \). Thus the ratio is free of \( \rho \) if and only if \( T(x,y)=T(x',y') \). \Cref{thm:lehmann-scheffe-ratio} yields that \( T \) is minimal sufficient.

\textbf{(b), not complete.}
Each marginal is \( X_i\sim\mathcal{N}(0,1) \) and \( Y_i\sim\mathcal{N}(0,1) \), for every \( \rho\in(-1,1) \). Therefore
\[
\E_\rho[T_1]
=
\E_\rho\Bigl[\sum_{i=1}^{n}X_i^{2}\Bigr]
+
\E_\rho\Bigl[\sum_{i=1}^{n}Y_i^{2}\Bigr]
=
n+n
=
2n
\]
for every \( \rho \). The function \( u(t_1,t_2)=t_1-2n \) then satisfies \( \E_\rho[u(T)]=0 \) for every \( \rho\in(-1,1) \). It is not almost surely zero: \( T_1 \) is absolutely continuous (e.g.\ at \( \rho=0 \) one has \( T_1\sim\chi^{2}_{2n} \), and for \( \rho\neq 0 \) the law of \( T_1 \) is still nondegenerate, as \( \Var_\rho(T_1)>0 \)). \Cref{def:complete} therefore fails: the family of laws of \( T \) is not complete.

(The same conclusion is the curved-exponential-family obstruction of \cref{thm:expfam-complete}. The density is of the form \( \exp\bigl(\eta_1(\rho)T_1+\eta_2(\rho)T_2-A(\rho)\bigr) \) with
\[
\eta_1(\rho)
=
-\frac{1}{2(1-\rho^{2})},
\qquad
\eta_2(\rho)
=
\frac{\rho}{1-\rho^{2}}.
\]
The map \( \rho\mapsto(\eta_1(\rho),\eta_2(\rho)) \) traces a curve in \( \mathbb{R}^{2} \), which contains no nonempty open set. The regular two-parameter completeness theorem does not apply, and the explicit annihilator \( T_1-2n \) shows that completeness indeed fails.)
\end{proof}

\begin{remark}[What (a) is not]
\label{rem:p9-normal-var}
Under the additional bivariate-normal hypothesis of part~(b), the limit variance collapses to the classical value \( c=(1-\rho^{2})^{2} \). That identity uses Gaussian moment relations (Isserlis / Wick) and is not available in (a). Writing \( (1-\rho^{2})^{2} \) as the answer to (a) is the wrong model.

The vector \( V \) cannot be shortened to \( (XY) \) or to \( (\hat{\rho}) \). The sample correlation recentres and rescales, so the CLT must be run on the five raw moments that enter \( g \).
\end{remark}

\begin{remark}[What (b) is not]
\label{rem:p9-five-param}
If all five parameters \( (\mu_X,\mu_Y,\sigma_X^{2},\sigma_Y^{2},\rho) \) were unknown, the five sums
\[
\Bigl(\sum_{i}X_i,\,\sum_{i}Y_i,\,\sum_{i}X_i^{2},\,\sum_{i}Y_i^{2},\,\sum_{i}X_i Y_i\Bigr)
\]
would be jointly complete sufficient \citep[Exercise~7.5.8]{hogg2019}. Restricting to \( \mu_X=\mu_Y=0 \) and \( \sigma_X^{2}=\sigma_Y^{2}=1 \) drops three coordinates of the natural parameter and produces a curved subfamily. Completeness is lost; minimal sufficiency of the remaining pair \( (T_1,T_2) \) is not.

The statistic \( \bigl(\sum X_i^{2},\sum Y_i^{2},\sum X_i Y_i\bigr) \) is sufficient for \( \rho \) on the present model (the likelihood does not use a finer split of \( T_1 \)), but it is not minimal: two samples with the same \( (T_1,T_2) \) and a different split of \( \sum X_i^{2} \) versus \( \sum Y_i^{2} \) have proportional likelihoods.
\end{remark}

\subsection{Problem 10. [10 pts.]}
Let \( X_1, \dots, X_n \) be iid from \( N(0, \sigma^2) \) with \( \sigma > 0 \) unknown. Let \( W = \frac{1}{n} \sum_{i=1}^n |X_i| \). Recall \( \chi^2_k \) has probability density function

\[
\frac{1}{2^{k/2} \Gamma(k/2)} x^{k/2 - 1} e^{-x/2}, \qquad x \ge 0.
\]

\begin{enumerate}[label=(\alph*)]
    \item Find an unbiased estimator of \( \sigma \) based on \( W \) and identify its limiting distribution.
    \item Is it admissible under the square error loss? Justify your answer.
\end{enumerate}

\setcounter{section}{10}

\subsubsection{Definitions used in Problem 10(b)}

Part~(b) is a comparison of risk \emph{functions}, not a statement about unbiasedness or about a single value of \( \sigma \). The loss is the same squared error as in Problem~7. An estimator is admissible when nothing dominates it uniformly on \( \sigma>0 \). The definition does not produce a test; to decide the present \( \delta \) one still needs a competitor (typically Rao--Blackwell of \( \delta \) onto \( \sum X_i^{2} \)).

\begin{definition}[Squared error loss]
\label{def:se-loss}
The \emph{squared error loss} for estimating a real parameter \( g(\theta) \) by an action \( a\in\mathbb{R} \) is
\[
L\bigl(\theta,a\bigr)
=
\bigl(a-g(\theta)\bigr)^{2}.
\]
On this problem \( \theta=\sigma>0 \) and \( g(\sigma)=\sigma \), so \( L(\sigma,a)=(a-\sigma)^{2} \). This is the loss already used in Problem~7 and in \cref{prop:l2-posterior-mean}.
\end{definition}

\begin{definition}[Risk]
\label{def:risk}
The \emph{risk} of a decision rule \( \delta \) at \( \theta \) is the expected loss under the sampling experiment:
\[
R(\theta,\delta)
=
\E_\theta\bigl[L\bigl(\theta,\delta(X)\bigr)\bigr].
\]
Under \cref{def:se-loss} this is mean squared error,
\[
R(\sigma,\delta)
=
\E_\sigma\bigl[(\delta-\sigma)^{2}\bigr]
=
\Var_\sigma(\delta)
+
\bigl(\E_\sigma[\delta]-\sigma\bigr)^{2}.
\]
Risk is a function of \( \sigma \), not a number. See \citep[\S7.1]{hogg2019}.
\end{definition}

\begin{definition}[Domination]
\label{def:dominates}
An estimator \( \delta' \) \emph{dominates} \( \delta \) (under the given loss) if
\[
R(\theta,\delta')
\le
R(\theta,\delta)
\qquad\text{for every }\theta\in\Theta,
\]
and the inequality is strict for at least one \( \theta \). Pointwise improvement at a single parameter value is not domination.
\end{definition}

\begin{definition}[Admissible estimator]
\label{def:admissible}
An estimator \( \delta \) of \( g(\theta) \) is \emph{admissible} under the loss \( L \) if no estimator dominates it in the sense of \cref{def:dominates}. Equivalently: if \( R(\theta,\delta')\le R(\theta,\delta) \) for every \( \theta \), then \( R(\,\cdot\,,\delta')=R(\,\cdot\,,\delta) \). The estimator is \emph{inadmissible} if some \( \delta' \) dominates it.

Admissibility is a statement about the whole risk function on the parameter space of the problem (here \( \sigma>0 \)). It is not implied by unbiasedness, and it is not implied by having a \( \sqrt{n} \)-Gaussian limit. A unique Bayes rule for a prior with full support is admissible; that sufficient criterion is not needed to read the definition, and is not the expected write-up of part~(b).
\end{definition}

\begin{theorem}[Rao--Blackwell]
\label{thm:rao-blackwell}
Let \( T \) be sufficient for \( \theta \), and let \( \delta \) be unbiased for \( g(\theta) \) with \( \E_\theta[\delta^{2}]<\infty \). Set \( \delta^{*}= \E[\delta\mid T] \). Then \( \delta^{*} \) is unbiased for \( g(\theta) \), is a function of \( T \), and
\[
\Var_\theta(\delta^{*})
\le
\Var_\theta(\delta)
\]
for every \( \theta \), with equality if and only if \( \delta \) is a function of \( T \) almost surely. Under squared error this is \( R(\theta,\delta^{*})\le R(\theta,\delta) \), strict unless \( \delta \) is already a function of \( T \). This is \citep[Theorem~7.3.1]{hogg2019}.
\end{theorem}

\begin{theorem}[Lehmann--Scheff\'e]
\label{thm:ls}
If \( T \) is complete sufficient and \( \varphi(T) \) is unbiased for \( g(\theta) \), then \( \varphi(T) \) is the unique MVUE of \( g(\theta) \). See \citep[Theorem~7.4.1]{hogg2019}. Combined with \cref{thm:expfam-complete}: in a regular one-parameter exponential family, the unique unbiased function of \( T=\sum K(X_i) \) is the unique MVUE.
\end{theorem}

\subsubsection{Hint for Problem 10(a)}

The target is \( \sigma \), not \( \sigma^{2} \). The usable machine is a scale multiple of \( W \) itself. Each \( |X_i| \) is half-normal: \( \E[|X_1|]=\sigma\sqrt{2/\pi} \), so \( \delta=\sqrt{\pi/2}\,W \) is unbiased. The CLT for the i.i.d.\ sequence \( |X_i| \), followed by the constant factor \( \sqrt{\pi/2} \), is the limiting distribution. Do not expand \( \E[W^{2}] \), and do not construct an unbiased estimator of \( \sigma^{2} \). The displayed \( \chi^{2} \) density is for part~(b): \( \sum X_i^{2}/\sigma^{2}\sim\chi^{2}_{n} \).

\subsubsection{Manuscript proof idea}

The calculation below follows the handwritten notes. \mserr{Red} marks a false claim or a missing justification. The notes treat a different estimand and drop the scale in the half-normal mean.

\textbf{The estimand.}
The exam asks for an unbiased estimator of \( \sigma \) based on \( W \), and for the limiting distribution of that estimator. The notes compute \( \E[W^{2}] \) and conclude
\[
\mserr{nW^{2}-(n-1)\tfrac{2}{\pi}\text{ is unbiased for }\sigma^{2}}.
\]
That is a (broken) construction for \( \sigma^{2} \). Part~(a) is linear in \( W \), not quadratic.

\textbf{The half-normal mean.}
The integral setup is right:
\[
\E[|X_1|]
=
2\int_{0}^{\infty}x\cdot\frac{1}{\sqrt{2\pi\sigma^{2}}}\,e^{-x^{2}/(2\sigma^{2})}\,\mathrm{d}x.
\]
The notes record
\[
\mserr{\E[|X_1|]=\sqrt{2/\pi}}.
\]
The factor \( \sigma \) is missing. The substitution \( u=x^{2}/(2\sigma^{2}) \) produces \( \E[|X_1|]=\sigma\sqrt{2/\pi} \). Consequently the independence computation should be \( \E[|X_i||X_j|]=\sigma^{2}\cdot(2/\pi) \) for \( i\neq j \), not \( 2/\pi \). The notes keep \( \E[X_i^{2}]=\sigma^{2} \) while dropping \( \sigma \) from the first-moment identities, so the later expansion of \( \E[W^{2}] \) is internally inconsistent.

\textbf{What is not asked.}
Even after repairing the scale, \( \E[W^{2}]=\sigma^{2}\bigl(n^{-1}+2(n-1)/(n\pi)\bigr) \) only yields an unbiased estimator of \( \sigma^{2} \). The limiting law of \( W \) (or of \( \sqrt{\pi/2}\,W \)) is not computed. The \( \chi^{2} \) density is copied and unused; it is not needed for (a).

\errpanel{%
\textbf{\mserr{The error, in one box.}}

Part~(a) estimates \( \sigma \), not \( \sigma^{2} \). Use \( \E[|X_1|]=\sigma\sqrt{2/\pi} \), hence \( \delta=\sqrt{\pi/2}\,W \).

Do not drop \( \sigma \) from the half-normal mean, and do not expand \( \E[W^{2}] \).

The limit is the CLT for \( |X_i| \), scaled by \( \sqrt{\pi/2} \).
}

\subsubsection{Solution of Problem 10}

\paragraph{Answer.}
The estimator \( \delta=\sqrt{\pi/2}\,W \) is unbiased for \( \sigma \), and
\[
\sqrt{n}\,(\delta-\sigma)
\xrightarrow{(d)}
\mathcal{N}\bigl(0,\,\sigma^{2}(\pi/2-1)\bigr).
\]
For \( n\ge 2 \) it is inadmissible under squared error loss.

\begin{proof}[Correct proof]
\textbf{(a), unbiasedness.}
Write \( Z\sim\mathcal{N}(0,1) \), so \( X_1=\sigma Z \). Then \( |X_1|=\sigma|Z| \) and it is enough to compute \( \E[|Z|] \). The even integrand gives
\[
\E[|X_1|]
=
2\int_{0}^{\infty}x\cdot\frac{1}{\sqrt{2\pi\sigma^{2}}}\,e^{-x^{2}/(2\sigma^{2})}\,\mathrm{d}x.
\]
The substitution \( u=x^{2}/(2\sigma^{2}) \) yields \( \mathrm{d}u=x\,\mathrm{d}x/\sigma^{2} \), hence
\[
\int_{0}^{\infty}x\,e^{-x^{2}/(2\sigma^{2})}\,\mathrm{d}x
=
\sigma^{2},
\]
and
\[
\E[|X_1|]
=
\frac{2\sigma^{2}}{\sqrt{2\pi\sigma^{2}}}
=
\sigma\sqrt{\frac{2}{\pi}}.
\]
Linearity and identical distribution produce \( \E[W]=\sigma\sqrt{2/\pi} \). Therefore
\[
\delta
\coloneqq
\sqrt{\frac{\pi}{2}}\,W
=
\frac{1}{n}\sqrt{\frac{\pi}{2}}\sum_{i=1}^{n}|X_i|
\]
satisfies \( \E_\sigma[\delta]=\sigma \) for every \( \sigma>0 \). This is the unbiased estimator of \( \sigma \) based on \( W \).

\textbf{(a), limiting distribution.}
The variables \( |X_1|,\dots,|X_n| \) are i.i.d.\ with mean \( \sigma\sqrt{2/\pi} \) and
\[
\Var(|X_1|)
=
\E[X_1^{2}]
-
\bigl(\E[|X_1|]\bigr)^{2}
=
\sigma^{2}
-
\sigma^{2}\cdot\frac{2}{\pi}
=
\sigma^{2}\Bigl(1-\frac{2}{\pi}\Bigr)
\in
(0,\infty).
\]
\Cref{thm:clt} yields
\[
\sqrt{n}\Bigl(W-\sigma\sqrt{\frac{2}{\pi}}\Bigr)
\xrightarrow{(d)}
\mathcal{N}\Bigl(0,\,\sigma^{2}\bigl(1-2/\pi\bigr)\Bigr).
\]
Multiplying by the constant \( \sqrt{\pi/2} \) gives
\[
\sqrt{n}\,(\delta-\sigma)
=
\sqrt{\frac{\pi}{2}}\cdot\sqrt{n}\Bigl(W-\sigma\sqrt{\frac{2}{\pi}}\Bigr)
\xrightarrow{(d)}
\mathcal{N}\Bigl(0,\,\sigma^{2}\cdot\frac{\pi}{2}\Bigl(1-\frac{2}{\pi}\Bigr)\Bigr)
=
\mathcal{N}\bigl(0,\,\sigma^{2}(\pi/2-1)\bigr).
\]

\textbf{(b).}
The model is a regular one-parameter exponential family in the scale. The density is
\[
f(x\mid\sigma)
=
(2\pi\sigma^{2})^{-1/2}
\exp\bigl(-x^{2}/(2\sigma^{2})\bigr)
=
\exp\bigl(-x^{2}/(2\sigma^{2})-\tfrac12\log(2\pi\sigma^{2})\bigr),
\]
so \( K(x)=x^{2} \). \Cref{thm:expfam-complete} yields that
\[
T
\coloneqq
\sum_{i=1}^{n}X_i^{2}
\]
is complete sufficient for \( \sigma^{2} \), hence also for \( \sigma \) (a one-to-one reparametrisation). Equivalently, \( T/\sigma^{2}\sim\chi^{2}_{n} \) by the exam density, which is a regular scale family in \( \sigma^{2} \).

The statistic \( W=n^{-1}\sum|X_i| \) is not a function of \( T \) when \( n\ge 2 \). On the sphere \( \{x: \sum x_i^{2}=t\} \) the value of \( \sum|x_i| \) is not constant: it ranges between \( \sqrt{t} \) (mass on a coordinate axis) and \( \sqrt{n t} \) (equal absolute coordinates). Thus \( \delta=\sqrt{\pi/2}\,W \) is unbiased with finite variance and is not a function of \( T \). \Cref{thm:rao-blackwell} produces the competitor \( \delta^{*}=\E[\delta\mid T] \), which is unbiased and satisfies
\[
R(\sigma,\delta^{*})
=
\Var_\sigma(\delta^{*})
<
\Var_\sigma(\delta)
=
R(\sigma,\delta)
\qquad\text{for every }\sigma>0.
\]
Hence \( \delta^{*} \) dominates \( \delta \) in the sense of \cref{def:dominates}, and \( \delta \) is inadmissible under \cref{def:admissible,def:se-loss}.

(The same conclusion follows from \cref{thm:ls} without computing \( \delta^{*} \): the unique MVUE of \( \sigma \) is the unbiased multiple of \( \sqrt{T} \),
\[
\varphi(T)
=
\frac{\Gamma(n/2)}{\sqrt{2}\,\Gamma((n+1)/2)}\,\sqrt{T},
\]
because \( \E[\sqrt{T}]=\sigma\E[\sqrt{\chi^{2}_{n}}]=\sigma\sqrt{2}\,\Gamma((n+1)/2)/\Gamma(n/2) \). For \( n\ge 2 \) one has \( \varphi(T)\neq\delta \) on a set of positive probability, so \( \delta \) is not the MVUE and cannot be admissible among unbiased estimators, a fortiori not admissible in the larger class.)
\end{proof}

\begin{remark}[What (a) is not]
\label{rem:p10-sigma2}
An unbiased estimator of \( \sigma^{2} \) based on \( W^{2} \) is a different problem. After restoring the scale one has \( \E[W^{2}]=\sigma^{2}\bigl(n^{-1}+2(n-1)/(n\pi)\bigr) \), so \( W^{2} \) over that constant is unbiased for \( \sigma^{2} \). That estimator is still a function of \( W \), still not a function of \( T \) for \( n\ge 2 \), and is likewise inadmissible for \( \sigma^{2} \). It does not answer (a).

For \( n=1 \), \( W=|X_1|=\sqrt{T} \) and \( \delta \) coincides with \( \varphi(T) \). The domination argument needs \( n\ge 2 \), which is the intended range once an average of several \( |X_i| \) is written.
\end{remark}

\subsection{Problem 11. [15 pts.]}
Let \( X_1, \dots, X_n \) be iid from exponential distribution with density \( f(x; \lambda, \alpha) = \lambda e^{-\lambda(x-\alpha)} I_{\{x \ge \alpha\}} \), where \( \lambda > 0 \) and \( -\infty < \alpha < \infty \) are unknown parameters.
\begin{enumerate}[label=(\alph*)]
    \item Find the MLE of \( (\lambda, \alpha) \), say \( (\hat{\lambda}, \hat{\alpha}) \).
    \item What distribution does \( \frac{\sqrt{n}(\hat{\lambda} - \lambda)}{n(\hat{\alpha} - \alpha)} \) converge to? You will receive partial credits if you can identify the distribution(s) of the numerator and/or denominator.
\end{enumerate}

\setcounter{section}{11}

\subsubsection{Hint for Problem 11}

Assume \( n\ge 2 \), otherwise \( \overline{X}=\min_i X_i \) and \( \hat{\lambda} \) is undefined. The support depends on \( \alpha \), so \( \alpha \) is not a regular interior parameter: do not differentiate through the indicator and set \( \partial\ell/\partial\alpha=0 \). On the half-space \( \{\alpha\le X_{(1)}\} \) the log-likelihood is strictly increasing in \( \alpha \), hence the maximum is on the boundary \( \hat{\alpha}=X_{(1)} \). Then maximise in \( \lambda \).

Part~(b) is not \( \E[\overline{X}-X_{(1)}] \). The denominator is already at the extreme-value scale \( n(\hat{\alpha}-\alpha) \), which is exactly Exponential and does not need a limit. The numerator is a regular \( \sqrt{n} \)-CLT for the rate, after the \( n-1 \) spacings above the minimum. The two pieces are independent, so the ratio converges to a normal over an independent exponential. Partial credit is exactly that identification.

\begin{theorem}[Exponential spacings]
\label{thm:exp-spacings}
Let \( Z_1,\dots,Z_n \) be i.i.d.\ with density \( \lambda e^{-\lambda z} \) on \( (0,\infty) \). Write \( Z_{(1)}<\cdots<Z_{(n)} \) for the order statistics, \( Z_{(0)}\coloneqq 0 \), and
\[
\Delta_1
\coloneqq
n Z_{(1)},
\qquad
\Delta_k
\coloneqq
(n-k+1)\bigl(Z_{(k)}-Z_{(k-1)}\bigr),
\qquad
k=2,\dots,n.
\]
Then \( \Delta_1,\dots,\Delta_n \) are i.i.d.\ with the same Exponential\( (\lambda) \) law (rate \( \lambda \), mean \( 1/\lambda \)). In particular
\[
n Z_{(1)}
=
\Delta_1
\sim
\mathrm{Exp}(\lambda),
\qquad
\sum_{i=1}^{n}\bigl(Z_i-Z_{(1)}\bigr)
=
\sum_{k=2}^{n}\Delta_k
\sim
\mathrm{Gamma}(n-1,\lambda)
\]
are independent. This is the standard transformation of the order-statistic density in \citep[\S4.4]{hogg2019}.
\end{theorem}

\begin{proof}
The joint density of \( (Z_{(1)},\dots,Z_{(n)}) \) is
\[
n!\,\lambda^{n}\exp\Bigl(-\lambda\sum_{i=1}^{n}z_{(i)}\Bigr)
\]
on \( 0<z_{(1)}<\cdots<z_{(n)} \). The map to \( (\Delta_1,\dots,\Delta_n) \) is linear and triangular, with Jacobian \( 1/n! \). The identity
\[
\sum_{i=1}^{n}z_{(i)}
=
\sum_{k=1}^{n}\Delta_k
\]
converts the joint density into \( \lambda^{n}\exp(-\lambda\sum_{k}\Delta_k) \) on \( (0,\infty)^{n} \), which is the product of \( n \) Exponential\( (\lambda) \) densities. The two displayed sums are \( \Delta_1 \) and \( \sum_{k\ge 2}\Delta_k \), hence independent, and the second is \( \mathrm{Gamma}(n-1,\lambda) \) (rate \( \lambda \)).
\end{proof}

\subsubsection{Manuscript proof idea}

The calculation below follows the handwritten notes. \mserr{Red} marks a false claim or a missing justification. The skeleton of (a) is the right exam machine: the log-likelihood increases in \( \alpha \) up to the constraint, then the score in \( \lambda \) is solved at that boundary. Part~(b) computes the right means and stops one step short of the question.

\textbf{(a).}
The notes write
\[
\ell
=
-\lambda\sum_{i}x_i
+
n\lambda\alpha
+
n\log\lambda
\]
and \( \partial\ell/\partial\lambda=-\sum x_i+n\alpha+n/\lambda \), hence \( \hat{\lambda}=1/(\overline{x}-\alpha) \) for fixed \( \alpha \), then \( \partial\ell/\partial\alpha=n\lambda>0 \) so \( \hat{\alpha}=\min_i X_i \) and \( \hat{\lambda}=1/(\overline{X}-\min_i X_i) \). The two displayed estimators are the correct MLE. \mserr{The formula for \( \ell \) is valid only on \( \{\alpha\le\min_i x_i\} \); the indicators must be written.} Once that restriction is in place, \( \partial\ell/\partial\alpha=n\lambda \) is an interior derivative on \( \alpha<\min_i x_i \), not a critical-point equation. The maximum is on the boundary because \( \ell \) is increasing in \( \alpha \) there. Do not set \( \partial\ell/\partial\alpha=0 \).

\textbf{(b).}
The survival function \( \Pbb(\min_i X_i\ge y)=e^{-n\lambda(y-\alpha)} \) for \( y>\alpha \) is correct, and so are
\[
\E[\min_i X_i]
=
\alpha+\frac{1}{n\lambda},
\qquad
\E[X_1]
=
\alpha+\frac{1}{\lambda},
\qquad
\E[\overline{X}-\min_i X_i]
=
\frac{1}{\lambda}\Bigl(1-\frac{1}{n}\Bigr).
\]
Those identities identify the means of the two pieces of \( (\hat{\lambda},\hat{\alpha}) \). \mserr{Part~(b) asks for the limiting law of the ratio \( \sqrt{n}(\hat{\lambda}-\lambda)/\bigl(n(\hat{\alpha}-\alpha)\bigr) \), not for \( \E[\overline{X}-\min_i X_i] \).} The denominator is already tight at the \( n \)-scale and does not concentrate; an expectation of order \( 1/\lambda \) does not name its limit.

\errpanel{%
\textbf{\mserr{The gap, in one box.}}

The MLE is \( \hat{\alpha}=X_{(1)} \), \( \hat{\lambda}=1/(\overline{X}-X_{(1)}) \), after maximising on \( \{\alpha\le X_{(1)}\} \).

Part~(b) is the joint limit of \( \bigl(\sqrt{n}(\hat{\lambda}-\lambda),\,n(\hat{\alpha}-\alpha)\bigr) \). The second coordinate is exactly \( \mathrm{Exp}(\lambda) \); the first is asymptotically \( \mathcal{N}(0,\lambda^{2}) \); they are independent. The ratio converges to that normal over that exponential. Do not stop at \( \E[\overline{X}-X_{(1)}] \).
}

\subsubsection{Solution of Problem 11}

\paragraph{Answer.}
Write \( X_{(1)}=\min_i X_i \). The MLE is
\[
\hat{\alpha}
=
X_{(1)},
\qquad
\hat{\lambda}
=
\frac{1}{\overline{X}-X_{(1)}}
=
\frac{n}{\sum_{i=1}^{n}(X_i-X_{(1)})},
\]
for \( n\ge 2 \). The ratio converges in distribution to \( Z/E \), where \( Z\sim\mathcal{N}(0,\lambda^{2}) \) is independent of \( E\sim\mathrm{Exp}(\lambda) \) (rate \( \lambda \)). Equivalently, the limit is \( \lambda^{2}\,G/E_1 \) with \( G\sim\mathcal{N}(0,1) \) independent of \( E_1\sim\mathrm{Exp}(1) \).

\begin{proof}[Correct proof]
\textbf{(a).}
The joint density on \( \mathbb{R}^{n} \) is
\[
L(\lambda,\alpha)
=
\lambda^{n}
\exp\Bigl(-\lambda\sum_{i=1}^{n}(x_i-\alpha)\Bigr)
\ind_{\{\alpha\le x_{(1)}\}},
\]
and \( L=0 \) if \( \alpha>x_{(1)} \). On the half-space \( \alpha\le x_{(1)} \) the log-likelihood is
\[
\ell(\lambda,\alpha)
=
n\log\lambda
-
\lambda\sum_{i=1}^{n}x_i
+
n\lambda\alpha.
\]
For fixed \( \lambda>0 \) one has \( \partial\ell/\partial\alpha=n\lambda>0 \), so \( \ell \) is strictly increasing in \( \alpha \) up to the boundary \( \alpha=x_{(1)} \). The maximiser in \( \alpha \) is therefore \( \hat{\alpha}=x_{(1)} \), independently of \( \lambda \). Substituting \( \alpha=x_{(1)} \) leaves
\[
\ell(\lambda,x_{(1)})
=
n\log\lambda
-
\lambda\sum_{i=1}^{n}(x_i-x_{(1)}).
\]
The score in \( \lambda \) is \( n/\lambda-\sum(x_i-x_{(1)}) \), with unique zero
\[
\hat{\lambda}
=
\frac{n}{\sum_{i=1}^{n}(x_i-x_{(1)})}
=
\frac{1}{\overline{x}-x_{(1)}}
\]
provided \( n\ge 2 \) (for \( n=1 \) the denominator vanishes). The second derivative \( -n/\lambda^{2}<0 \) confirms a maximum. Thus \( (\hat{\lambda},\hat{\alpha}) \) is the displayed MLE.

\textbf{(b), exact laws.}
Write \( Z_i=X_i-\alpha \), so the \( Z_i \) are i.i.d.\ Exponential\( (\lambda) \) on \( (0,\infty) \), \( \hat{\alpha}-\alpha=Z_{(1)} \), and \( \sum(X_i-\hat{\alpha})=\sum(Z_i-Z_{(1)}) \). \Cref{thm:exp-spacings} gives the exact pair
\[
n(\hat{\alpha}-\alpha)
=
n Z_{(1)}
\sim
\mathrm{Exp}(\lambda),
\qquad
S
\coloneqq
\sum_{i=1}^{n}(X_i-\hat{\alpha})
\sim
\mathrm{Gamma}(n-1,\lambda)
\quad\text{(rate \(\lambda\))},
\]
and \( n(\hat{\alpha}-\alpha) \) is independent of \( S \). In particular \( \hat{\lambda}=n/S \) is independent of the denominator \( n(\hat{\alpha}-\alpha) \), for every finite \( n\ge 2 \). (The notes' identities \( \E[X_{(1)}]=\alpha+1/(n\lambda) \) and \( \E[\overline{X}-X_{(1)}]=(n-1)/(n\lambda) \) are the means of these two laws.)

\textbf{(b), numerator.}
The sum \( S \) is the sum of \( n-1 \) i.i.d.\ Exponential\( (\lambda) \) variables, so \( \E[S]=(n-1)/\lambda \) and \( \Var(S)=(n-1)/\lambda^{2} \). Hence
\[
\sqrt{n}\Bigl(\frac{S}{n}-\frac{1}{\lambda}\Bigr)
=
\frac{S-\E[S]}{\sqrt{n}}
-
\frac{1}{\lambda\sqrt{n}}
=
\sqrt{\frac{n-1}{n}}\cdot\frac{S-\E[S]}{\sqrt{n-1}}
-
\frac{1}{\lambda\sqrt{n}}.
\]
\Cref{thm:clt} sends the standardised Gamma average to \( \mathcal{N}(0,1/\lambda^{2}) \), and the remaining two factors tend to \( 1 \) and \( 0 \). Therefore
\[
\sqrt{n}\Bigl(\frac{S}{n}-\frac{1}{\lambda}\Bigr)
\xrightarrow{(d)}
\mathcal{N}\bigl(0,1/\lambda^{2}\bigr).
\]
The map \( g(u)=1/u \) is \( C^{1} \) at \( u=1/\lambda \), with \( g'(1/\lambda)=-\lambda^{2} \). \Cref{thm:delta-method}(i) yields
\[
\sqrt{n}\,(\hat{\lambda}-\lambda)
=
\sqrt{n}\Bigl(g(S/n)-g(1/\lambda)\Bigr)
\xrightarrow{(d)}
\mathcal{N}\bigl(0,\,\lambda^{2}\bigr).
\]

\textbf{(b), the ratio.}
Write \( U_n=\sqrt{n}(\hat{\lambda}-\lambda) \) and \( V_n=n(\hat{\alpha}-\alpha) \). Then \( U_n \) is a (measurable) function of \( S \), \( V_n \) is a function of \( Z_{(1)} \), and \( S\perp Z_{(1)} \), so \( U_n\perp V_n \). Moreover \( V_n\sim\mathrm{Exp}(\lambda) \) exactly, hence \( V_n\Rightarrow E \) with \( E\sim\mathrm{Exp}(\lambda) \), while \( U_n\Rightarrow Z\sim\mathcal{N}(0,\lambda^{2}) \). Jointly,
\[
(U_n,V_n)
\xrightarrow{(d)}
(Z,E)
\]
with \( Z\perp E \). The map \( (u,v)\mapsto u/v \) is continuous on \( \mathbb{R}\times(0,\infty) \), and \( \Pbb(E>0)=1 \), so
\[
\frac{U_n}{V_n}
=
\frac{\sqrt{n}(\hat{\lambda}-\lambda)}{n(\hat{\alpha}-\alpha)}
\xrightarrow{(d)}
\frac{Z}{E}.
\]
Rescaling \( Z=\lambda G \) and \( E=E_1/\lambda \) with \( G\sim\mathcal{N}(0,1) \) independent of \( E_1\sim\mathrm{Exp}(1) \) rewrites the same limit as \( \lambda^{2}\,G/E_1 \).
\end{proof}

\begin{remark}[Why the two rates]
\label{rem:p11-rates}
The parameter \( \alpha \) sits in the support. The information for \( \alpha \) is infinite in the usual interior sense, and \( \hat{\alpha}-\alpha=O_{\Pbb}(n^{-1}) \), not \( O_{\Pbb}(n^{-1/2}) \). That is why the exam already multiplies the denominator by \( n \). The rate \( \lambda \) is regular after the minimum is removed, so the numerator is a genuine \( \sqrt{n} \)-CLT. Quoting a two-dimensional Fisher information matrix for \( (\lambda,\alpha) \) is the wrong machine.

The notes' computation of \( \E[\overline{X}-X_{(1)}] \) is the mean of \( S/n \). It checks the Gamma law; it does not name \( Z/E \).
\end{remark}

\begin{thebibliography}{99}

\bibitem[Durrett(2010)]{durrett2010}
Rick Durrett.
\newblock \emph{Probability: Theory and Examples}.
\newblock 4th~ed., Cambridge University Press, 2010.
\newblock Conditional expectation: \S4.1. Characteristic functions and uniqueness: Theorems~3.3.1, 3.3.2. L\'evy continuity: Theorem~3.3.6. Weak convergence: Theorems~3.2.3, 3.2.5. Classical CLT: Theorem~3.4.1. Gaussian vectors and Cram\'er--Wold: \S3.3, \S3.9. Stopping times: \S5.2. Optional stopping: Theorems~5.7.3--5.7.6. Simple random walk, gambler's ruin, and exponential martingales: \S4.1, \S5.7, \S5.10. Brownian motion, hitting times, and \( B_t^{2}-t \): \S8.1--\S8.5.

\bibitem[Hogg, McKean and Craig(2019)]{hogg2019}
Robert~V. Hogg, Joseph~W. McKean, and Allen~T. Craig.
\newblock \emph{Introduction to Mathematical Statistics}.
\newblock 8th~ed., Pearson, 2019.
\newblock Order statistics: \S4.4. Joint density of \( (Y_{1},Y_{n}) \): \( n(n-1)(F(y)-F(x))^{n-2}f(x)f(y) \). Exponential spacings: \S4.4. Convergence in probability: Definition~5.1.1. Consistency of an estimator: Definition~5.1.2. Consistency of \( Y_n=\max_i X_i \) for Uniform\( (0,\theta) \): Example~5.1.2 and (5.1.1). Bias versus consistency: \S5.1.1. Slutsky: Theorem~5.2.5. Delta method: Theorems~5.2.9, 5.4.6. Classical CLT: Theorem~5.3.1. Multivariate CLT: Theorem~5.4.4. Invariance of the MLE: Theorem~6.1.2. Sufficiency and factorization: \S7.2, Theorem~7.2.1. Minimal sufficiency and the likelihood-ratio criterion: \S7.3, Theorem~7.3.1. Completeness: \S7.4. Regular exponential-family completeness: Theorems~7.5.1--7.5.2. Five-parameter bivariate normal, joint complete sufficient statistics: Exercise~7.5.8. Size, UMP, monotone likelihood ratio, Karlin--Rubin: \S8.1--\S8.2, (8.1.3), Definitions~8.2.1--8.2.2, (8.2.3), (8.2.6). Two-sided point null has no UMP: Example~8.2.3. Risk function: \S7.1. Rao--Blackwell: Theorem~7.3.1. Lehmann--Scheff\'e: Theorem~7.4.1. Bayes estimator under squared error (posterior mean): \S11.1.

\bibitem[Ren and Liu(2024)]{ren2024}
任佳刚, 刘继成.
\newblock \emph{概率论教程}.
\newblock Lecture notes, compiled 10~September 2024, 300~pp.
\newblock Uncorrelated versus independent; bivariate normal: Theorem~6.2.1 and Example~3 in \S6.2. Characteristic functions and uniqueness: Theorem~8.4.4.

\end{thebibliography}

\end{document}
